\documentclass[letterpaper, 9
 pt, conference]{ieeeconf}

% Wider columns + taller text block to reduce page count
\geometry{left=0.55in, right=0.55in, top=0.75in, bottom=0.75in}
\setlength{\columnsep}{0.25in}

\usepackage{amsmath,amssymb,amsfonts,amsthm,bm}
\usepackage{mathtools}
\usepackage{booktabs}
\usepackage{enumitem}
\usepackage[numbers,sort&compress]{natbib}
\usepackage{hyperref}
\usepackage{xcolor}
\usepackage{multirow}
\usepackage{graphicx}
\graphicspath{{figs/}}
\usepackage{caption}
\setlength{\textfloatsep}{8pt plus 1pt minus 2pt}
\setlength{\floatsep}{6pt plus 1pt minus 2pt}
\setlength{\intextsep}{6pt plus 1pt minus 2pt}
\usepackage{stfloats}
\usepackage[section]{placeins}
\usepackage{algorithm}
\usepackage{algpseudocode}

\hypersetup{
    colorlinks=true,
    linkcolor=blue,
    citecolor=blue,
    urlcolor=blue
}


\title{Towards Certifiably Informative Measurement Selection for Robotic Calibration}
\author{Parth Mahajan\\Northeastern University\\mahajan.parth@northeastern.edu}

% Append the teaser to \@maketitle so it lands inside the one-column header
% that ieeeconf creates via \twocolumn[\@maketitle].  This puts Figure 1
% below the title/author block and above the two-column abstract.
\makeatletter
\let\orig@maketitle\@maketitle
\def\@maketitle{%
  \orig@maketitle
  \par\vskip 8pt
  {\centering
   \includegraphics[width=\linewidth]{Teaser}\par}%
  \vskip 4pt
  \captionof{figure}{%
    \textbf{Certifiably informative calibration data design in practice.}
    Given a finite pool of candidate calibration views
    (left: real-world ChArUco captures), we select a fixed-budget subset
    in a single-shot manner by optimizing information-theoretic criteria
    on the marginal Fisher information matrix of the calibration parameters.
    The selected views serve as acquisition references for a human operator,
    indicating which poses or viewpoints are most informative. As the budget
    $K$ increases, the minimum eigenvalue of the marginal information matrix
    improves, increasing excitation of the weakest calibration direction and
    reducing local parameter uncertainty. The selected subset also maintains
    broad image-plane coverage, as shown by the corner-density heatmap
    (center-right). Evaluation on held-out images yields subpixel reprojection
    RMSE, demonstrating that the information-selected subset remains predictive
    beyond the views used for calibration.}%
  \label{fig:teaser}%
  \par\vskip 8pt
}
\makeatother

\begin{document}

\maketitle

% \begin{abstract}



% Calibration is a fundamental prerequisite for robotic systems. Its quality depends not only on the backend estimator, but also on whether the collected data is sufficiently informative enough to excite the calibration parameters. In practice, calibration datasets are often gathered using operator intuition, motion coverage, or geometric heuristics. Existing calibration pipelines such as Kalibr
% primarily address the estimation problem after data collection, while active calibration and
% next best view methods guide acquisition sequentially using local information or coverage criteria. As a result, they provide limited assurance that the final dataset is globally informative
% for accurate parameter recovery. This motivates a central question: how can one certify that a
% chosen set of calibration views, trajectories, or motion segments is sufficiently informative for
% calibration?

% Inspired by OASIS, we formulate calibration data collection as a certifiable experimental design problem. We cast calibration measurement selection as a fixed budget subset selection problem under information theoretic objectives. Using the marginal Fisher information matrix obtained after eliminating nuisance variables, we study E , A , and D optimal subset selection together with Frank Wolfe relaxations. These relaxations provide post-hoc certificates through duality gaps and valid upper bounds on the discrete optimum for all three design criteria. On tractable instances, FW-E recovers solutions within 2.5\% of the true discrete E-optimum, FW-D is exact in all tested cases, and FW-A remains within 10.6\%.

% We evaluate the proposed framework by estimating camera intrinsics with two calibration targets, AprilGrid and ChArUco, in photorealistic Isaac Sim benchmarks spanning multiple corner-measurement noise levels 
% ($\sigma \in \{0.01, 0.5, 1.0\}$ px) and distortion regimes. We also validate the approach on two different real world camera types and compare its behavior under both OpenCV and Kalibr backends. Across these settings, information driven selection consistently outperforms random sampling and geometric heuristics, while training reprojection error is shown to be a poor proxy for true calibration quality. The results further reveal that criterion choice is regime dependent and interestingly depends on the backend solver, views budget K and measurement noise regime: greedy E optimal Or Wolf Frank E optimal are effective in param recovery under fixed pose OpenCV calibration backend, whereas greedy A optimality is usally a more safer bet and is robust under high corner noise, small view budgets, and joint bundle adjustment backends such as Kalibr. Overall, these findings support treating calibration as a certifiable data design problem in which the optimal selection criterion depends on view budget, noise level, and estimator structure.

% \end{abstract}

% \begin{abstract}
\begin{abstract}
Calibration accuracy in robotic systems depends not only on the backend estimator, but also on whether the collected data sufficiently excites the calibration parameters. In practice, calibration datasets are often acquired using operator intuition, motion coverage, or geometric heuristics, while pipelines such as Kalibr optimize parameters only after data collection. Active and next best view calibration methods guide acquisition sequentially, but provide limited post hoc assurance that the final dataset is globally informative for parameter recovery. This raises an important question: can we certify the information theoretic richness of a calibration dataset?

We formulate calibration data collection as a certifiable information theoretic experimental design problem. Given a finite pool of candidate views, poses, or motion segments, we select a fixed budget subset using the marginal Fisher information matrix obtained after eliminating nuisance variables. We study E, A, and D optimal design criteria and solve the resulting subset selection problem using greedy methods and Frank Wolfe relaxations, which provide upper bounds and post hoc duality gap certificates for the design surrogate.

We evaluate the framework on camera intrinsic calibration using AprilGrid and ChArUco targets in photorealistic Isaac Sim benchmarks, two real world capture sessions, and both OpenCV and Kalibr backends. Information driven selection consistently outperforms random sampling and geometric heuristics, while training reprojection error is shown to be a poor proxy for calibration quality. Our results show that criterion choice is regime dependent: E optimality performs well for the OpenCV intrinsics benchmark, while A optimality is more robust under high noise, small budgets, and joint bundle adjustment backends such as Kalibr.
\end{abstract}
% \begin{abstract}
% Calibration is a fundamental pre requisite for robotics systems. Its quality depends not only on the backend estimator, but also on
% whether the collected data is sufficiently informative enough to excite the calibration parameters. In practice, however, calibration datasets are still often gathered by operators intuition or coverage heuristics or motion coverage, while current classic pipelines like kalibr optimize the parameters only after the data have already been collected or unlike active calibration next best view methods like aprical or calib wizard that guide acquisition sequentially using local
% information gain criterias but provide no assurances on the informativeness of the collected data for good calibration. This leaves open an important question: how can one certify that a chosen
% set of calibration views, trajectories, or motion segments is sufficiently
% informative for accurate parameter recovery?
% taking inspiration from OASIS, we formulate calibration data collection as a certifiable
% experimental design problem.
% We cast the selection of calibration measurements as a fixed budget subset selection problem under an information theoretic objective that strengthens the weakest
% observable directions of the calibration parameter space. Although the resulting optimization problem is NP hard in general, we show
% how greedy selection together with convex relaxation based post hoc verification can enable the efficient recovery of certifiably informative calibration datasets in practice.

% We study E , A , and D optimal subset selection together with Frank Wolfe relaxations whose duality gaps provide post hoc certificates for
% all three continuous design objectives.
% Unlike next best view methods that guide acquisition sequentially using local
% information gain, our formulation targets a global fixed budget subset whose
% informativeness can be assessed after selection.
% We show that the relaxations
% provide valid upper bounds on the discrete optimum for all three criteria; FW E recovers solutions within near optimal values 2.5\% of the true discrete E optimum, FW-D is
% exact in all tested cases, and FW A remains within 10.6\%.

% We evaluate the resulting framework on AprilGrid and ChArUco
% benchmarks using Photorealistic IsaacSim simulations, across multiple noise levels 
% and distortion regimes, and two real world different camera capture sessions, and with an
% additional Kalibr backend study vs opencv backend study.
% We find out that Information driven selection consistently outperforms geometric heuristics
% and random sampling, while training reprojection error proves to be a weak indicator of true calibration quality.
% The results also show interestingly shows that the best design criterion is nuianced and the choice of criteria is dependent on three things the budget K views. Under a fixed pose OpenCV backend, E-optimality effectively strengthens weak
% observable directions and often improves focal length recovery, whereas
% A optimality is more robust at higher noise and for joint bundle-adjustment
% backends such as Kalibr. This yields a practical guide: A optimality is the safest default, especially
% for joint bundle adjustment, high noise, or small budget settings, while
% FW E or Greedy-E are preferred for fixed pose solvers when the view budget is sufficient and corner noise is moderate.
% \end{abstract}

% -------------------------------------------------------

% Calibration is a foundational requirement for modern robotic perception
% systems.
% Accurate estimates of camera intrinsics, sensor extrinsics, temporal
% offsets, and other latent parameters are essential for downstream tasks such
% as localization, mapping, sensor fusion, and control.
% Although substantial progress has been made in calibration solvers and
% toolchains, the success of these methods still depends strongly on the
% quality of the collected calibration data.
% Poorly chosen views, trajectories, or motion patterns can leave important
% parameter directions weakly observable, resulting in unstable estimates and
% degraded downstream performance.

% In current practice, calibration data collection is often guided by simple
% heuristics such as target coverage or practitioner intuition.
% These strategies can be effective, but they do not explicitly reason about
% whether the resulting dataset is sufficiently informative for accurate
% parameter recovery.
% Classical calibration pipelines such as Kalibr are highly effective at
% estimating multi-camera and visual-inertial calibration parameters from a
% provided dataset, but they largely assume that the dataset has already been
% collected and do not directly solve the upstream problem of selecting
% certifiably informative measurements \citep{furgale2013kalibr,kalibr_toolbox}.
% Likewise, foundational camera calibration methods such as Zhang's
% planar-pattern formulation provide highly practical and influential estimation
% procedures, but not a principled notion of optimal data selection
% \citep{zhang2000flexible}.

% A related line of work comes from next-best-view and active calibration
% methods, which select sensing actions or poses to improve visibility, reduce
% uncertainty, or maximize information gain.
% Systems such as AprilCal and Calibration Wizard explicitly acknowledge that
% calibration quality depends strongly on the acquired data and therefore guide
% the user toward improved views \citep{richardson2013aprilcal,peng2019calibration}.
% Later work on efficient pose selection, next-best-view planning for
% eye-in-hand calibration, and observability-aware active trajectory optimization
% further reinforces the importance of calibration-aware data acquisition
% \citep{rojtberg2018efficient,yang2023nbv,wang2025active}.
% However, these methods are typically framed as sequential next-best-pose or
% trajectory-planning procedures rather than as certifiable global
% subset-selection problems over candidate calibration data.

% Another growing research direction is deep learning-based calibration,
% especially for targetless, online, or weakly supervised settings.
% Methods such as RegNet, CalibNet, and LCCNet reduce manual effort and can
% directly predict extrinsic corrections from multimodal data
% \citep{schneider2017regnet,iyer2018calibnet,lv2021lccnet}.
% More recent work continues to improve automation and online refinement
% \citep{liao2023survey,huang2025whatmatters}.
% At the same time, recent survey and critique papers note persistent challenges
% in robustness, generalization, and consistent evaluation across sensor types,
% scenes, and camera models \citep{liao2023survey,huang2025whatmatters}.
% As a result, classical geometric pipelines remain the dominant choice in
% precision-oriented robotic systems, even as learning-based methods continue
% to broaden the scope of what can be automated.

% In this paper, we cast calibration data collection as an information-theoretic
% experimental design problem with post-hoc certificates.
% Rather than treating calibration purely as a downstream estimation task, we
% formulate the selection of calibration views, poses, or motion segments as a
% fixed-budget subset-selection problem over candidate measurements.
% Using the marginal Fisher information matrix obtained after eliminating
% nuisance poses, we study E-, A-, and D-optimal design criteria together with
% greedy selection and Frank--Wolfe relaxations.
% The resulting framework makes it possible to assess the informativeness of a
% selected dataset directly, to certify continuous-relaxation quality across
% all three criteria through duality gaps, and to validate discrete solution
% quality against exhaustive optima on small instances.

% The main contributions of this framing are:
% \begin{enumerate}[leftmargin=1.5em]
%     \item We formulate robotic calibration data collection as a subset
%     selection problem over candidate views, poses, or motion segments using a
%     marginal calibration Fisher information matrix.
%     \item We study certifiable information-theoretic design for calibration
%     through greedy and Frank--Wolfe optimisation of E-, A-, and D-optimal
%     criteria, including post-hoc certificates for the relaxed problem across
%     all three objectives and exhaustive validation on small discrete
%     instances.
%     \item We show experimentally that information-driven selection
%     outperforms geometric heuristics, but that the preferred criterion is
%     solver- and regime-dependent, leading to a practical guide based on
%     backend structure, noise level, and view budget.
\section{Introduction}
    Calibration is a foundational requirement for modern robotic perception systems.
Accurate estimates of camera intrinsics, sensor extrinsics, temporal offsets, and
other latent parameters are essential for localization, mapping, sensor fusion,
and control. Although substantial progress has been made in calibration solvers
and toolchains, calibration accuracy depends not only on the backend estimator,
but also on whether the collected data sufficiently excites the calibration
parameters. Poorly chosen views, trajectories, or motion patterns can leave
important parameter directions weakly observable, resulting in unstable estimates
and degraded downstream performance.

In practice, calibration data collection is still often guided by target
coverage, motion excitation heuristics, or practitioner intuition. These
strategies are simple and effective in many settings, but they do not explicitly
measure whether the resulting dataset is informative for parameter recovery.
Classical calibration pipelines such as Kalibr estimate multi camera and
visual inertial calibration parameters from a provided dataset, but they largely
assume that the dataset has already been collected. Similarly, foundational
camera calibration methods such as Zhang's planar-pattern formulation provide
practical and influential estimation procedures, but do not address the upstream
problem of selecting an informative calibration dataset.

Active calibration and next best view methods address this limitation by
guiding the acquisition process. Systems such as AprilCal and Calibration Wizard
recognize that calibration quality depends on the acquired data and therefore
suggest improved views to the user. More recent methods for pose selection,
eye in hand calibration, and observability aware trajectory optimization further
reinforce the importance of calibration aware data acquisition. However, these
methods are typically framed as sequential acquisition procedures based on local
coverage, visibility, or information gain criteria. They do not generally provide
a post hoc certificate that the final fixed budget dataset is close to optimal
under a global information theoretic objective.

This motivates the central question of this paper: given a finite set of
candidate calibration views, poses, or motion segments, can we select a
fixed budget subset whose informativeness can be certified after selection?

We answer this question by casting calibration data collection as an
information theoretic experimental design problem. Inspired by OASIS, we
formulate calibration measurement selection as a subset selection problem over
candidate views, poses, or motion segments. We construct a marginal calibration
Fisher information matrix by eliminating nuisance pose variables by Shurs Compliment and study
E , A , and D optimal design criteria. These objectives capture complementary
notions of dataset quality: E optimality strengthens the weakest observable
parameter direction, A optimality reduces average parameter uncertainty, and
D optimality maximizes overall information volume.

The resulting subset selection problem is combinatorial and NP hard in general.
To make the problem practical, we combine greedy selection with Frank Wolfe
relaxations over the continuous design domain. The Frank Wolfe duality gap
provides a post hoc certificate for the relaxed information design problem,
while exhaustive search on small instances allows us to measure the quality of
rounded discrete solutions. Importantly, the certificate applies to the
information theoretic design surrogate, not to the full nonlinear calibration
pipeline. We therefore treat it as model relative evidence that a dataset is
informative, and validate its practical value through downstream calibration
experiments.

The main contributions of this paper are:
\begin{enumerate}[leftmargin=1.5em]
    \item We formulate robotic calibration data collection as a fixed budget
    subset selection problem over candidate views, poses, or motion segments
    using a marginal calibration Fisher information matrix.

    \item We study certifiable information theoretic design for calibration
    under E , A , and D optimal objectives using greedy selection and
    Frank Wolfe relaxations, including post hoc certificates for the relaxed
    objectives and exhaustive validation on small discrete instances.

    \item We evaluate the framework on AprilGrid and ChArUco calibration
    benchmarks, photorealistic Isaac Sim simulations, real world camera capture
    sessions, and OpenCV and Kalibr backend studies, showing that
    information driven selection outperforms random sampling and geometric
    heuristics.

    \item We show that the preferred design criterion is regime dependent:
    E optimality is effective in our OpenCV intrinsics benchmark when weak
    observable directions dominate and the problem remains close to the
    linearised Fisher model, while A optimality is more robust under
    higher noise, small view budgets, and joint bundle adjustment backends such
    as Kalibr.
\end{enumerate}


% -------------------------------------------------------
% \section{Related Work}

% \subsection{Classical Calibration Backends}
% Classical calibration methods estimate sensor intrinsics, extrinsics, and
% timing parameters from a collected dataset.
% Zhang's planar-pattern method remains a foundational reference for practical
% camera calibration because it combines a closed-form initialization with
% nonlinear refinement under a maximum-likelihood objective
% \citep{zhang2000flexible}.
% In robotics, Kalibr has become one of the most widely used practical backends
% for multi-camera and visual-inertial calibration, supporting spatial and
% temporal calibration together with rich camera model support
% \citep{furgale2013kalibr,kalibr_toolbox}.
% These methods are highly effective at estimation, but they primarily operate
% on data that has already been collected.
% Consequently, they do not directly answer the upstream question of how to
% collect a calibration dataset that is certifiably informative.

% \subsection{Self-Calibration, Observability, and Informative Data Selection}
% A second line of work studies calibration through the lens of observability
% and information content.
% Maye et al.\ introduced a generic self-supervised calibration framework that
% used information-theoretic measures to identify and retain novel measurement
% sequences while handling unobservable directions explicitly
% \citep{maye2013self}.
% Schneider et al.\ later showed that visual-inertial self-calibration can be
% posed as a segment-based problem over a small subset of informative motion
% segments, yielding accuracy comparable to full-batch approaches at much lower
% computational cost \citep{schneider2017informative}.
% Their subsequent observability-aware work further emphasized that calibration
% quality depends strongly on motion excitation and proposed information metrics
% to assess which trajectory segments are worth retaining for calibration
% \citep{schneider2019observability}.
% Related work in online visual-inertial self-calibration provides a more
% complete observability analysis and identifies degenerate motion primitives
% that render subsets of calibration parameters unobservable
% \citep{yang2023online}.
% This literature strongly supports the premise that calibration failure is
% often caused by weak observability and poor excitation, not merely by the
% optimizer.

% \subsection{Guided Capture, Next-Best-View, and Active Calibration}
% Another closely related direction explicitly guides the collection process.
% AprilCal introduced an assisted and repeatable camera calibration workflow
% aimed at reducing reliance on operator intuition by suggesting improved
% target placements and providing intuitive uncertainty summaries
% \citep{richardson2013aprilcal}.
% Calibration Wizard pushed this idea further by computing, from previously
% acquired images, the next pose expected to yield the largest reduction in
% intrinsic uncertainty while also accounting for corner-detection uncertainty
% \citep{peng2019calibration}.
% Additional work on efficient pose selection, next-best-view selection for
% robot eye-in-hand calibration, and observability-aware active trajectory
% optimization continues this trend
% \citep{rojtberg2018efficient,yang2023nbv,wang2025active}.
% These methods show that active calibration can significantly reduce the number
% of required views or improve parameter uncertainty.
% However, they are usually formulated as sequential guidance or local planning
% procedures and do not provide a certificate that the final selected
% calibration set is globally optimal or near-optimal under a global
% subset-selection objective.

% \subsection{Information-Theoretic Online Calibration}
% Information-theoretic calibration ideas have also been explored in online and
% multi-camera settings.
% Dexheimer et al.\ developed an information-theoretic framework for online
% multi-camera extrinsic calibration that used entropy-based keyframe selection
% and segment evaluation to keep the estimation problem tractable
% \citep{dexheimer2022information}.
% These methods are conceptually close to the present viewpoint because they
% explicitly distinguish informative from redundant data.
% Nevertheless, they typically use heuristics such as entropy or keyframe
% scoring rather than a certifiable E-optimal subset-selection formulation over
% candidate calibration measurements.

% \subsection{Certifiable Calibration Solvers}
% There is also a growing literature on certifiable calibration solvers.
% Giamou et al.\ formulated extrinsic calibration from per-sensor egomotion as
% a quadratically constrained quadratic program and solved it via semidefinite
% relaxation to obtain certifiably globally optimal estimates
% \citep{giamou2019certifiable}.
% Follow-on work extended this line to online-capable variants and related
% hand-eye settings.
% This literature is highly relevant conceptually, but its certification target
% is different from ours: these methods certify the solution of a fixed
% calibration estimation problem, whereas our goal is to certify the
% informativeness of the selected dataset itself.
% In other words, certifiable solver work addresses \emph{how to solve} the
% posed calibration problem globally, while our framing addresses \emph{how to
% collect} data that makes calibration well-conditioned in the first place.

% \subsection{Targetless Geometric Calibration}
% Targetless geometric calibration provides another important neighboring
% direction, especially for cross-sensor calibration.
% A classical example is the mutual-information-based LiDAR--camera calibration
% method of Pandey et al., which aligns sensor observations without fiducial
% targets and even reports a Cram\'er--Rao lower bound analysis
% \citep{pandey2015mutual}.
% This line of work demonstrates that calibration can be carried out directly
% from natural scenes, making in-situ recalibration practical.
% Yet these methods still optimize a registration objective over the data at
% hand; they typically do not ask whether the selected data is globally optimal
% under a calibration-specific experimental design objective.

% \subsection{Learning-Based Calibration}
% A final major direction is learning-based calibration.
% Early milestones such as RegNet cast multimodal registration as direct deep
% regression of the six-degree-of-freedom extrinsic transform
% \citep{schneider2017regnet}, while CalibNet introduced self-supervised
% geometric training without direct supervision on the calibration parameters
% \citep{iyer2018calibnet}.
% LCCNet later improved online LiDAR--camera self-calibration by exploiting
% cost-volume style feature matching and iterative refinement \citep{lv2021lccnet}.
% More recent surveys show that learning-based calibration has expanded
% substantially across standard camera calibration, distortion estimation,
% cross-view calibration, and cross-sensor calibration \citep{liao2023survey}.
% At the same time, recent critique work argues that many regression-style
% LiDAR--camera calibration networks still struggle in complex real scenes
% \citep{huang2025whatmatters}.
% For this reason, learning-based calibration is best viewed as complementary
% to, rather than a replacement for, high-accuracy geometric calibration.
% Most importantly for our setting, these methods improve automation and
% targetless operation, but they generally do not certify observability or
% the informativeness of the collected data.

% \subsection{Relation to OASIS}
% The work most closely aligned with our certifiable-design perspective is
% OASIS, which formulates sensor arrangement for SLAM as a subset selection
% problem under an E-optimality criterion and combines greedy selection with
% convex-relaxation-based post-hoc verification \citep{kaveti2023oasis}.
% OASIS shows that robotics design problems can be treated as
% information-theoretic subset selection problems with meaningful certificates.
% Our work borrows that conceptual template, but shifts the focus from sensor
% arrangement for SLAM to the selection of informative views, poses, or motion
% segments for calibration.
% Thus, unlike Kalibr-style pipelines that estimate from a fixed dataset,
% unlike AprilCal- and Calibration-Wizard-style methods that guide data
% acquisition sequentially, and unlike certifiable solver approaches that
% certify only the downstream optimizer, our goal is to certify the
% informativeness of calibration data itself.

\section{Related Work}

\subsection{Calibration Backends, Observability, and Informative Data}
Classical calibration pipelines estimate sensor intrinsics, extrinsics, and
temporal offsets from a collected dataset. Zhang's planar pattern method
remains a foundational approach for practical camera calibration, combining a
closed-form initialization with nonlinear maximum likelihood refinement
\citep{zhang2000flexible}. In robotics, Kalibr is widely used for
multi camera and visual inertial calibration, supporting spatial and temporal
calibration with flexible camera models
\citep{furgale2013kalibr,kalibr_toolbox}. These systems are effective
estimation backbends, but they primarily operate after the calibration data has
already been collected.

A complementary line of work studies calibration through observability and
information content. Maye et al.\ used information theoretic measures to retain
novel measurement sequences while accounting for unobservable directions
\citep{maye2013self}. Schneider et al.\ showed that visual inertial
self calibration can be performed using a small subset of informative motion
segments, achieving accuracy comparable to full batch estimation with lower
computational cost \citep{schneider2017informative}. Their later work further
emphasized that calibration quality depends strongly on motion excitation and
trajectory observability \citep{schneider2019observability}. Recent online
visual inertial calibration work similarly identifies degenerate motions that
make subsets of parameters unobservable \citep{yang2023online}. These works
motivate our central premise: calibration failure is often caused not only by
the optimizer, but also by insufficiently informative data.

\subsection{Guided Capture and Active Calibration}
Several methods explicitly guide calibration data collection. AprilCal proposed
an assisted and repeatable camera calibration workflow that reduces reliance on
operator intuition by suggesting improved target placements and providing
uncertainty summaries \citep{richardson2013aprilcal}. Calibration Wizard
extended this idea by recommending the next target pose expected to reduce
intrinsic uncertainty, while accounting for corner-detection uncertainty
\citep{peng2019calibration}. Other work has studied efficient pose selection,
next-best-view planning for robot eye-in-hand calibration, and
observability-aware active trajectory optimization
\citep{rojtberg2018efficient,yang2023nbv,wang2025active}.

These methods show that active calibration can reduce the number of required
views and improve parameter uncertainty. However, they are typically formulated
as sequential guidance or local planning procedures. In contrast, our goal is
to select a batch of informative calibration views, poses, or motion segments
and provide post-hoc evidence that the selected set is globally or
near-globally informative under a specified experimental-design objective.

\subsection{Certifiable Estimation and Certifiable Data Design}
Certifiable calibration solvers provide guarantees for the solution of a fixed
estimation problem. For example, Giamou et al.\ formulated extrinsic
calibration from per-sensor egomotion as a quadratically constrained quadratic
program and used semidefinite relaxation to obtain certifiably globally optimal
estimates \citep{giamou2019certifiable}. Related work has extended this idea to
online-capable and hand-eye calibration settings. This certification target is
different from ours: certifiable solvers certify the optimizer output for a
given dataset, whereas our work aims to certify the informativeness of the
dataset selected before estimation.

The closest inspiration for our design perspective is OASIS, which formulates
sensor arrangement for SLAM as an E-optimal subset-selection problem and uses
greedy selection together with convex-relaxation-based post-hoc verification
\citep{kaveti2023oasis}. OASIS demonstrates that robotics design problems can
be treated as information-theoretic subset-selection problems with meaningful
certificates. Our work adapts this viewpoint from sensor arrangement for SLAM
to calibration data design. Rather than certifying a downstream calibration
solution, we certify the selected calibration views or motion segments with
respect to an information-based design surrogate.

\subsection{Targetless and Learning-Based Calibration}
Targetless and learning-based calibration methods provide another neighboring
direction. Mutual-information-based LiDAR-camera calibration aligns
cross-sensor observations without fiducial targets and has been used for
in-situ calibration from natural scenes \citep{pandey2015mutual}. Learning-based
methods such as RegNet, CalibNet, and LCCNet cast calibration as direct
regression, self-supervised geometric alignment, or iterative feature-based
refinement \citep{schneider2017regnet,iyer2018calibnet,lv2021lccnet}. Recent
surveys show that learning-based calibration has expanded across camera,
distortion, cross-view, and cross-sensor settings \citep{liao2023survey},
although recent critique work suggests that regression-style LiDAR-camera
calibration networks can still struggle in complex real-world scenes
\citep{huang2025whatmatters}.

These methods improve automation and reduce dependence on calibration targets,
but they generally do not certify whether the collected data sufficiently
excites the calibration parameters. Our work is therefore complementary: it
focuses on information-driven and certifiable selection of calibration data,
independent of the particular downstream estimator used to refine the final
parameters.

\section{Problem Formulation}
\label{sec:problem_formulation}

We consider the problem of selecting a small but maximally informative calibration dataset from a larger pool of candidate measurements. In this paper, we instantiate the formulation for camera intrinsic calibration, where the goal is to estimate focal lengths, principal point, skew, and distortion parameters from observations of a known calibration target. However, the formulation is not limited to camera intrinsics. The same structure applies to any sensor calibration problem in which candidate measurements contribute information about calibration parameters after nuisance variables are marginalized.

Thus, the framework can be extended to stereo calibration, multi-camera calibration, camera-IMU calibration, LiDAR-camera extrinsic calibration, hand-eye calibration, and other sensor calibration problems by redefining the parameter vector, measurement model, and per-candidate information blocks. Our formulation follows the OASIS-style view of perception design as budgeted subset selection over additive Fisher information matrices, with relaxation-based post-hoc certificates for near-optimality~\cite{kaveti2023oasis}.

\subsection{Candidate Calibration Measurements}
\label{subsec:candidate_measurements}

Let
\begin{equation}
    \theta \in \mathbb{R}^{d_\theta}
\end{equation}
denote the calibration parameters of interest. For intrinsic camera calibration, $\theta$ contains the camera intrinsic and distortion parameters. For other calibration problems, $\theta$ may contain relative sensor poses, temporal offsets, hand-eye transforms, or other calibration variables.

Assume that we are given a finite candidate set
\begin{equation}
    \mathcal{V}=\{v_1,\ldots,v_N\}.
\end{equation}
Each candidate $v_i$ may correspond to a calibration image, a target pose, a robot motion segment, a multi-sensor observation, or a proposed next-best-view action.

Selecting a calibration dataset is equivalent to selecting a subset
\begin{equation}
    \mathcal{S}\subseteq\mathcal{V}.
\end{equation}
We encode this subset using a binary design vector
\begin{equation}
    s\in\{0,1\}^{N},
\end{equation}
where
\begin{equation}
    s_i =
    \begin{cases}
        1, & v_i\in\mathcal{S},\\
        0, & v_i\notin\mathcal{S}.
    \end{cases}
\end{equation}
Given a calibration budget $K$, the selected design must satisfy
\begin{equation}
    \mathbf{1}^{\top}s=K.
\end{equation}
The central design question is therefore: which $K$ candidate measurements should be selected to minimize uncertainty in $\theta$?

\subsection{Generic MLE, Optimality Conditions, and Laplace Approximation}
\label{subsec:mle_laplace}

We first introduce the estimation problem abstractly before specifying the camera measurement model. Let $X$ denote nuisance variables that are needed to explain the measurements but are not the primary calibration parameters. For intrinsic camera calibration, $X$ contains the per-image calibration target poses. For multi-sensor calibration, $X$ may contain platform poses, target poses, trajectories, correspondence variables, temporal offsets, or other latent variables.

Define the full parameter vector
\begin{equation}
    \Theta=(\theta,X).
\end{equation}
For a selected design $s$, the likelihood of the available calibration measurements is
\begin{equation}
    p(Z\mid \Theta;s).
\end{equation}
The maximum-likelihood estimator is
\begin{equation}
    \hat{\Theta}(s)
    =
    \operatorname*{arg\,max}_{\Theta}
    p(Z\mid \Theta;s).
\end{equation}
Equivalently,
\begin{equation}
    \hat{\Theta}(s)
    =
    \operatorname*{arg\,min}_{\Theta}
    \mathcal{L}(\Theta;s),
\end{equation}
where
\begin{equation}
    \mathcal{L}(\Theta;s)
    =
    -\log p(Z\mid \Theta;s)
\end{equation}
is the negative log-likelihood.

Assume that $\hat{\Theta}(s)$ is a locally isolated minimizer of $\mathcal{L}$. Then the first-order necessary condition is
\begin{equation}
    \nabla_{\Theta}\mathcal{L}(\hat{\Theta}(s);s)=0.
\end{equation}
The second-order necessary condition is
\begin{equation}
    \nabla_{\Theta}^{2}\mathcal{L}(\hat{\Theta}(s);s)\succeq 0.
\end{equation}
Equivalently, defining the Hessian
\begin{equation}
    H_{\Theta}(\hat{\Theta}(s);s)
    =
    \nabla_{\Theta}^{2}
    \mathcal{L}(\hat{\Theta}(s);s),
\end{equation}
we have
\begin{equation}
    H_{\Theta}(\hat{\Theta}(s);s)\succeq 0.
\end{equation}
For the Laplace approximation to define a nonsingular local covariance, we further require the Hessian to be positive definite on the identifiable parameter subspace:
\begin{equation}
    H_{\Theta}(\hat{\Theta}(s);s)\succ 0.
\end{equation}
If the calibration problem has gauge freedoms or unobservable directions, this condition may fail unless the gauge is fixed, nuisance variables are marginalized properly, or a weak prior or regularizer is added.

A second-order Taylor expansion around $\hat{\Theta}(s)$ gives
\begin{align}
    \mathcal{L}(\Theta;s)
    \approx\;
    &\mathcal{L}(\hat{\Theta}(s);s)
    +
    \nabla_{\Theta}\mathcal{L}(\hat{\Theta}(s);s)^\top
    (\Theta-\hat{\Theta}(s)) \nonumber \\
    &+
    \frac{1}{2}
    (\Theta-\hat{\Theta}(s))^\top
    H_{\Theta}(\hat{\Theta}(s);s)
    (\Theta-\hat{\Theta}(s)).
\end{align}
Using the first-order condition, the linear term vanishes:
\begin{equation}
    \nabla_{\Theta}\mathcal{L}(\hat{\Theta}(s);s)=0.
\end{equation}
Thus,
\begin{equation}
    \mathcal{L}(\Theta;s)
    \approx
    \mathcal{L}(\hat{\Theta}(s);s)
    +
    \frac{1}{2}
    (\Theta-\hat{\Theta}(s))^\top
    H_{\Theta}(\hat{\Theta}(s);s)
    (\Theta-\hat{\Theta}(s)).
\end{equation}
Substituting this quadratic approximation into the likelihood gives the local Laplace approximation
\begin{equation}
    p(\Theta\mid Z;s)
    \approx
    \mathcal{N}
    \left(
    \hat{\Theta}(s),
    H_{\Theta}(\hat{\Theta}(s);s)^{-1}
    \right).
\end{equation}
Therefore,
\begin{equation}
    \operatorname{Cov}(\hat{\Theta}(s))
    \approx
    H_{\Theta}(\hat{\Theta}(s);s)^{-1}.
\end{equation}

Under this local Gaussian approximation, for a unit perturbation direction $u$, the local increase in the objective is
\begin{equation}
    \mathcal{L}(\hat{\Theta}(s)+\epsilon u;s)
    -
    \mathcal{L}(\hat{\Theta}(s);s)
    \approx
    \frac{\epsilon^2}{2}
    u^\top
    H_{\Theta}(\hat{\Theta}(s);s)
    u.
\end{equation}
Thus, $u^\top H_{\Theta}(\hat{\Theta}(s);s)u$ measures how strongly the selected dataset constrains direction $u$. Under the same Gaussian approximation, larger curvature corresponds to smaller marginal variance along that direction, while weakly constrained calibration modes correspond to directions with small curvature. This motivates choosing calibration measurements that shape the Hessian, or equivalently the local information matrix, to reduce uncertainty in the calibration parameters of interest.

\subsection{Camera Intrinsic Calibration Measurement Model}
\label{subsec:camera_measurement_model}

We now specialize the generic formulation to camera intrinsic calibration. Let the calibration target contain $M$ known points
\begin{equation}
    \mathcal{P}=\{p_j\}_{j=1}^{M},
    \qquad
    p_j\in\mathbb{R}^{3},
\end{equation}
expressed in the target coordinate frame.

For candidate image $v_i$, let
\begin{equation}
    x_i\in SE(3)
\end{equation}
denote the pose of the calibration target relative to the camera. The camera observes a subset of visible target points
\begin{equation}
    \mathcal{P}_i\subseteq\mathcal{P}.
\end{equation}
For each detected point $p_j\in\mathcal{P}_i$, the image measurement is
\begin{equation}
    z_{ij}
    =
    \pi(\theta,x_i,p_j)+\epsilon_{ij},
\end{equation}
where $\pi(\cdot)$ is the camera projection model and
\begin{equation}
    \epsilon_{ij}\sim\mathcal{N}(0,\Sigma_{ij})
\end{equation}
is image-space detection noise.

The likelihood of measurements from candidate image $i$ is
\begin{equation}
    p_i(Z_i\mid \theta,x_i)
    =
    \prod_{j\in\mathcal{P}_i}
    p(z_{ij}\mid \theta,x_i,p_j).
\end{equation}
Given a selected design $s$, the joint likelihood is
\begin{equation}
    p(Z\mid \theta,X;s)
    =
    \prod_{i=1}^{N}
    p_i(Z_i\mid \theta,x_i)^{s_i},
\end{equation}
where
\begin{equation}
    X=\{x_i\}_{i=1}^{N}.
\end{equation}

Under Gaussian pixel noise, the negative log-likelihood becomes a reprojection-error least-squares problem:
\begin{equation}
    (\hat{\theta}(s),\hat{X}(s))
    =
    \operatorname*{arg\,min}_{\theta,X}
    \sum_{i=1}^{N}
    s_i\ell_i(\theta,x_i),
\end{equation}
where
\begin{equation}
    \ell_i(\theta,x_i)
    =
    \sum_{j\in\mathcal{P}_i}
    \left\|
    z_{ij}
    -
    \pi(\theta,x_i,p_j)
    \right\|_{\Sigma_{ij}^{-1}}^{2}.
\end{equation}
Thus, the selected subset $s$ directly determines the nonlinear calibration problem used to estimate the intrinsics.

\subsection{Fisher Information and Additivity Across Candidate Views}
\label{subsec:fisher_information}

Under the Laplace approximation, local calibration uncertainty is governed by the Hessian of the negative log-likelihood. To obtain an information-theoretic design criterion that is independent of a particular noise realization, we write this curvature in terms of the Fisher information matrix.

For the full calibration state $\Theta=(\theta,X)$, the Fisher information matrix induced by a selected design $s$ is
\begin{equation}
    I(\Theta;s)
    =
    \mathbb{E}_{Z}
    \left[
    -
    \nabla_{\Theta}^{2}
    \log p(Z\mid \Theta;s)
    \right].
\end{equation}
Equivalently, under standard regularity conditions,
\begin{equation}
    I(\Theta;s)
    =
    \mathbb{E}_{Z}
    \left[
    \nabla_{\Theta}\log p(Z\mid \Theta;s)
    \nabla_{\Theta}\log p(Z\mid \Theta;s)^\top
    \right].
\end{equation}

For a selected calibration design, the joint likelihood factorizes as
\begin{equation}
    p(Z\mid \Theta;s)
    =
    \prod_{i=1}^{N}
    p_i(Z_i\mid \theta,x_i)^{s_i}.
\end{equation}
Therefore,
\begin{equation}
    \log p(Z\mid \Theta;s)
    =
    \sum_{i=1}^{N}
    s_i
    \log p_i(Z_i\mid \theta,x_i).
\end{equation}
Substituting this into the Fisher information definition gives
\begin{equation}
    I(\Theta;s)
    =
    \mathbb{E}_{Z}
    \left[
    -
    \nabla_{\Theta}^{2}
    \sum_{i=1}^{N}
    s_i
    \log p_i(Z_i\mid \theta,x_i)
    \right].
\end{equation}
By linearity of expectation and differentiation,
\begin{equation}
    I(\Theta;s)
    =
    \sum_{i=1}^{N}
    s_i
    \mathbb{E}_{Z_i}
    \left[
    -
    \nabla_{\Theta}^{2}
    \log p_i(Z_i\mid \theta,x_i)
    \right].
\end{equation}
Define the Fisher information contribution of candidate view $v_i$ as
\begin{equation}
    I_i(\Theta)
    \triangleq
    \mathbb{E}_{Z_i}
    \left[
    -
    \nabla_{\Theta}^{2}
    \log p_i(Z_i\mid \theta,x_i)
    \right].
\end{equation}
Then the total information matrix decomposes additively:
\begin{equation}
    I(\Theta;s)
    =
    \sum_{i=1}^{N}
    s_i I_i(\Theta).
\end{equation}

For candidate view $v_i$, the information matrix can be partitioned with respect to the calibration parameters $\theta$ and nuisance variables $x_i$:
\begin{equation}
    I_i(\theta,x_i)
    =
    \begin{bmatrix}
        A_i & B_i\\
        B_i^\top & C_i
    \end{bmatrix}.
\end{equation}
The blocks are defined in expectation form as
\begin{equation}
    A_i
    =
    \mathbb{E}_{Z_i}
    \left[
    -
    \nabla_{\theta\theta}^{2}
    \log p_i(Z_i\mid \theta,x_i)
    \right],
\end{equation}
\begin{equation}
    B_i
    =
    \mathbb{E}_{Z_i}
    \left[
    -
    \nabla_{\theta x_i}^{2}
    \log p_i(Z_i\mid \theta,x_i)
    \right],
\end{equation}
and
\begin{equation}
    C_i
    =
    \mathbb{E}_{Z_i}
    \left[
    -
    \nabla_{x_i x_i}^{2}
    \log p_i(Z_i\mid \theta,x_i)
    \right].
\end{equation}
Equivalently, using score functions,
\begin{equation}
    A_i
    =
    \mathbb{E}_{Z_i}
    \left[
    \nabla_{\theta}
    \log p_i(Z_i\mid \theta,x_i)
    \nabla_{\theta}
    \log p_i(Z_i\mid \theta,x_i)^\top
    \right],
\end{equation}
\begin{equation}
    B_i
    =
    \mathbb{E}_{Z_i}
    \left[
    \nabla_{\theta}
    \log p_i(Z_i\mid \theta,x_i)
    \nabla_{x_i}
    \log p_i(Z_i\mid \theta,x_i)^\top
    \right],
\end{equation}
and
\begin{equation}
    C_i
    =
    \mathbb{E}_{Z_i}
    \left[
    \nabla_{x_i}
    \log p_i(Z_i\mid \theta,x_i)
    \nabla_{x_i}
    \log p_i(Z_i\mid \theta,x_i)^\top
    \right].
\end{equation}

For the camera intrinsic calibration model, where each view contains conditionally independent target-point observations,
\begin{equation}
    p_i(Z_i\mid \theta,x_i)
    =
    \prod_{j\in\mathcal{P}_i}
    p(z_{ij}\mid \theta,x_i,p_j),
\end{equation}
the per-view Fisher information further decomposes as
\begin{equation}
    I_i(\theta,x_i)
    =
    \sum_{j\in\mathcal{P}_i}
    \mathbb{E}_{z_{ij}}
    \left[
    -
    \nabla_{(\theta,x_i)}^{2}
    \log p(z_{ij}\mid \theta,x_i,p_j)
    \right].
\end{equation}

Under the Gaussian reprojection-error model,
\begin{equation}
    z_{ij}
    =
    \pi(\theta,x_i,p_j)+\epsilon_{ij},
    \qquad
    \epsilon_{ij}\sim\mathcal{N}(0,\Sigma_{ij}),
\end{equation}
this Fisher information reduces locally to the Gauss-Newton form
\begin{equation}
    I_i(\theta,x_i)
    \approx
    J_i^\top \Sigma_i^{-1}J_i,
\end{equation}
where $J_i$ is the Jacobian of the stacked reprojection residuals with respect to $(\theta,x_i)$. Thus, the expectation-based Fisher information gives the general design object, while the Gauss-Newton matrix gives the practical computable form used for calibration view selection.

\subsection{Marginal Calibration Information via Schur Complement}
\label{subsec:schur_information}

The full Fisher information matrix includes both the calibration parameters $\theta$ and nuisance variables $X$. Since the goal is to estimate $\theta$, we eliminate nuisance variables using a Schur complement.

For each candidate view, the Fisher information has the block form
\begin{equation}
    I_i(\theta,x_i)
    =
    \begin{bmatrix}
        A_i & B_i\\
        B_i^\top & C_i
    \end{bmatrix}.
\end{equation}
Using the Schur complement, the information contributed by candidate view $v_i$ after eliminating the nuisance target pose is
\begin{equation}
    Q_i
    =
    A_i
    -
    B_iC_i^{-1}B_i^\top.
\end{equation}
If $C_i$ is singular or nearly singular, one may use a generalized Schur complement with $C_i^\dagger$, or add a weak pose prior:
\begin{equation}
    Q_i
    =
    A_i
    -
    B_iC_i^{\dagger}B_i^\top.
\end{equation}

The marginal information matrix over the calibration parameters is then
\begin{equation}
    H(s)
    =
    H_0
    +
    \sum_{i=1}^{N}
    s_iQ_i.
\end{equation}
Here, $H_0\succeq 0$ represents prior information, an initial calibration dataset, or a small regularization term. If no prior information is available, one may use
\begin{equation}
    H_0=\epsilon I
\end{equation}
for numerical stability.

Under the Laplace approximation,
\begin{equation}
    \operatorname{Cov}(\hat{\theta}(s))
    \approx
    H(s)^{-1}.
\end{equation}
Thus, calibration data selection becomes the problem of choosing candidate measurements whose accumulated marginal Fisher information matrix produces small local uncertainty in $\theta$.

\subsection{E-, D-, and A-Optimal Calibration Objectives}
\label{subsec:eda_objectives}

Let
\begin{equation}
    H(s)
    =
    H_0+\sum_{i=1}^{N}s_iQ_i.
\end{equation}
We consider three standard optimal experimental design criteria.

\paragraph{E-optimality.}
E-optimality maximizes the weakest observable calibration direction:
\begin{equation}
    \Phi_E(H)
    =
    \lambda_{\min}(H).
\end{equation}
Equivalently,
\begin{equation}
    \lambda_{\min}(H)
    =
    \min_{\|u\|=1}
    u^\top Hu.
\end{equation}
Thus, E-optimality solves
\begin{equation}
    s_E^\star
    =
    \operatorname*{arg\,max}_{s\in\{0,1\}^{N}}
    \lambda_{\min}
    \left(
    H_0+\sum_{i=1}^{N}s_iQ_i
    \right)
    \quad
    \text{s.t.}
    \quad
    \mathbf{1}^{\top}s=K.
\end{equation}

\paragraph{D-optimality.}
D-optimality maximizes information volume:
\begin{equation}
    \Phi_D(H)
    =
    \log\det(H).
\end{equation}
If $H=U\Lambda U^\top$ with eigenvalues $\lambda_1,\ldots,\lambda_{d_\theta}$, then
\begin{equation}
    \log\det(H)
    =
    \sum_{k=1}^{d_\theta}
    \log \lambda_k.
\end{equation}
Since the local covariance ellipsoid volume is proportional to $1/\sqrt{\det H}$, maximizing $\log\det(H)$ minimizes the volume of the local uncertainty ellipsoid. Thus, D-optimality solves
\begin{equation}
    s_D^\star
    =
    \operatorname*{arg\,max}_{s\in\{0,1\}^{N}}
    \log\det
    \left(
    H_0+\sum_{i=1}^{N}s_iQ_i
    \right)
    \quad
    \text{s.t.}
    \quad
    \mathbf{1}^{\top}s=K.
\end{equation}

\paragraph{A-optimality.}
A-optimality minimizes the average parameter variance:
\begin{equation}
    \operatorname{tr}(H^{-1}).
\end{equation}
To write all objectives as maximization problems, define
\begin{equation}
    \Phi_A(H)
    =
    -
    \operatorname{tr}(H^{-1}).
\end{equation}
Thus, A-optimality solves
\begin{equation}
    s_A^\star
    =
    \operatorname*{arg\,max}_{s\in\{0,1\}^{N}}
    -
    \operatorname{tr}
    \left[
    \left(
    H_0+\sum_{i=1}^{N}s_iQ_i
    \right)^{-1}
    \right]
    \quad
    \text{s.t.}
    \quad
    \mathbf{1}^{\top}s=K.
\end{equation}
Equivalently, in minimization form,
\begin{equation}
    s_A^\star
    =
    \operatorname*{arg\,min}_{s\in\{0,1\}^{N}}
    \operatorname{tr}
    \left[
    \left(
    H_0+\sum_{i=1}^{N}s_iQ_i
    \right)^{-1}
    \right]
    \quad
    \text{s.t.}
    \quad
    \mathbf{1}^{\top}s=K.
\end{equation}

\subsection{Unified Budgeted Calibration Design Problem}
\label{subsec:unified_design_problem}

Let
\begin{equation}
    \Phi\in\{\Phi_E,\Phi_D,\Phi_A\}
\end{equation}
be the chosen information criterion. The unified calibration design problem is
\begin{equation}
\boxed{
    s_\Phi^\star
    =
    \operatorname*{arg\,max}_{s\in\{0,1\}^{N}}
    \Phi
    \left(
    H_0+\sum_{i=1}^{N}s_iQ_i
    \right)
    \quad
    \text{s.t.}
    \quad
    \mathbf{1}^{\top}s=K.
}
\end{equation}
This is a combinatorial subset-selection problem over $\binom{N}{K}$ possible calibration datasets and is generally intractable for large candidate pools.

\subsection{Greedy Approximation}
\label{subsec:greedy_selection}

A simple approximation strategy is greedy subset selection. Starting from
\begin{equation}
    \mathcal{S}_0=\emptyset,
\end{equation}
we iteratively add the candidate that produces the largest increase in the chosen objective. At iteration $t$, define
\begin{equation}
    H_{t-1}
    =
    H_0+\sum_{v_j\in\mathcal{S}_{t-1}}Q_j.
\end{equation}
The next candidate is selected as
\begin{equation}
    v_t^\star
    =
    \operatorname*{arg\,max}_{v_i\in\mathcal{V}\setminus\mathcal{S}_{t-1}}
    \Phi(H_{t-1}+Q_i).
\end{equation}
Then
\begin{equation}
    \mathcal{S}_t
    =
    \mathcal{S}_{t-1}\cup\{v_t^\star\}.
\end{equation}
After $K$ iterations, greedy selection returns a feasible binary design $s_g$.

For E-optimality,
\begin{equation}
    v_{t,E}^\star
    =
    \operatorname*{arg\,max}_{v_i}
    \lambda_{\min}(H_{t-1}+Q_i).
\end{equation}
For D-optimality,
\begin{equation}
    v_{t,D}^\star
    =
    \operatorname*{arg\,max}_{v_i}
    \log\det(H_{t-1}+Q_i).
\end{equation}
For A-optimality,
\begin{equation}
    v_{t,A}^\star
    =
    \operatorname*{arg\,min}_{v_i}
    \operatorname{tr}
    \left[
    (H_{t-1}+Q_i)^{-1}
    \right].
\end{equation}

\subsection{Boolean Relaxation and Certifiable Suboptimality}
\label{subsec:boolean_relaxation}

To obtain a certificate, relax the binary variable $s_i\in\{0,1\}$ to a continuous weight
\begin{equation}
    \omega_i\in[0,1].
\end{equation}
Define the relaxed feasible set
\begin{equation}
    \mathcal{P}
    =
    \left\{
    \omega\in[0,1]^N:
    \mathbf{1}^{\top}\omega=K
    \right\}.
\end{equation}
The relaxed information matrix is
\begin{equation}
    H(\omega)
    =
    H_0+\sum_{i=1}^{N}\omega_iQ_i.
\end{equation}
The relaxed design problem is
\begin{equation}
    \mu_\Phi^\star
    =
    \max_{\omega\in\mathcal{P}}
    \Phi(H(\omega)).
\end{equation}
The relaxed feasible set contains the original integer feasible set. Therefore,
\begin{equation}
    \mu_\Phi^\star
    \geq
    f_\Phi^\star,
\end{equation}
where $f_\Phi^\star$ is the optimal value of the original discrete problem. Thus, for any feasible binary design $s$,
\begin{equation}
    f_\Phi^\star-f_\Phi(s)
    \leq
    \mu_\Phi^\star-f_\Phi(s).
\end{equation}
The quantity
\begin{equation}
    \Delta_\Phi(s)
    =
    \mu_\Phi^\star-f_\Phi(s)
\end{equation}
is a valid post-hoc upper bound on the suboptimality of $s$. A relative certificate can be written as
\begin{equation}
    \gamma_\Phi(s)
    =
    \frac{
    \mu_\Phi^\star-f_\Phi(s)
    }{
    |f_\Phi(s)|
    }.
\end{equation}

For E-optimality,
\begin{equation}
    \gamma_E(s)
    =
    \frac{
    \mu_E^\star-\lambda_{\min}(H(s))
    }{
    \lambda_{\min}(H(s))
    }.
\end{equation}
For D-optimality,
\begin{equation}
    \gamma_D(s)
    =
    \frac{
    \mu_D^\star-\log\det(H(s))
    }{
    |\log\det(H(s))|
    }.
\end{equation}
For A-optimality, it is often clearer to report the minimization form. Define
\begin{equation}
    \rho_A^\star
    =
    \min_{\omega\in\mathcal{P}}
    \operatorname{tr}(H(\omega)^{-1}).
\end{equation}
Then
\begin{equation}
    \rho_A^\star
    \leq
    a_A^\star
    \leq
    \operatorname{tr}(H(s)^{-1}),
\end{equation}
where $a_A^\star$ is the optimal value of the original discrete A-optimal minimization problem. The A-optimal suboptimality certificate is
\begin{equation}
    \Delta_A(s)
    =
    \operatorname{tr}(H(s)^{-1})
    -
    \rho_A^\star.
\end{equation}
A relative A-optimal certificate is
\begin{equation}
    \gamma_A(s)
    =
    \frac{
    \operatorname{tr}(H(s)^{-1})-\rho_A^\star
    }{
    \operatorname{tr}(H(s)^{-1})
    }.
\end{equation}

\subsection{Concavity of the Relaxed Objectives}
\label{subsec:concavity}

The map
\begin{equation}
    \omega
    \mapsto
    H(\omega)
    =
    H_0+\sum_{i=1}^{N}\omega_iQ_i
\end{equation}
is affine. The E-optimal objective $H\mapsto\lambda_{\min}(H)$ is concave over symmetric matrices. The D-optimal objective $H\mapsto\log\det(H)$ is concave over positive definite matrices. The A-optimal maximization objective $H\mapsto-\operatorname{tr}(H^{-1})$ is concave because $H\mapsto\operatorname{tr}(H^{-1})$ is convex over positive definite matrices. Therefore, the relaxed E-, D-, and A-optimal design problems are convex optimization problems in the sense of maximizing a concave function over a convex feasible set.

\subsection{Frank-Wolfe Optimization}
\label{subsec:frank_wolfe}

The relaxed problem can be solved efficiently using Frank-Wolfe. At iteration $t$, define
\begin{equation}
    H_t
    =
    H(\omega_t)
    =
    H_0+\sum_{i=1}^{N}\omega_{t,i}Q_i.
\end{equation}
The Frank-Wolfe update is
\begin{equation}
    \omega_{t+1}
    =
    (1-\eta_t)\omega_t+\eta_t v_t,
\end{equation}
where
\begin{equation}
    v_t
    =
    \operatorname*{arg\,max}_{v\in\mathcal{P}}
    \langle \nabla f_\Phi(\omega_t),v\rangle.
\end{equation}
Because $\mathcal{P}$ is a capped simplex, the linear oracle has a simple closed form: assign weight $1$ to the $K$ largest gradient entries and $0$ to the rest.

For E-optimality, let
\begin{equation}
    u_t=u_{\min}(H_t).
\end{equation}
If the smallest eigenvalue is simple, then
\begin{equation}
    \frac{\partial f_E}{\partial \omega_i}
    =
    u_t^\top Q_i u_t.
\end{equation}
If the minimum eigenvalue has multiplicity greater than one, any convex combination of projectors onto the minimum eigenspace gives a valid subgradient.

For D-optimality,
\begin{equation}
    \frac{\partial f_D}{\partial \omega_i}
    =
    \operatorname{tr}(H_t^{-1}Q_i).
\end{equation}
For A-optimality, with
\begin{equation}
    f_A(\omega)
    =
    -
    \operatorname{tr}(H(\omega)^{-1}),
\end{equation}
the gradient is
\begin{equation}
    \frac{\partial f_A}{\partial \omega_i}
    =
    \operatorname{tr}
    \left(
    H_t^{-1}Q_iH_t^{-1}
    \right).
\end{equation}

The Frank-Wolfe duality gap is
\begin{equation}
    g_t
    =
    \max_{v\in\mathcal{P}}
    \left\langle
    \nabla f_\Phi(\omega_t),
    v-\omega_t
    \right\rangle.
\end{equation}
For concave maximization, this gap upper-bounds the relaxed suboptimality:
\begin{equation}
    0
    \leq
    \mu_\Phi^\star-f_\Phi(\omega_t)
    \leq
    g_t.
\end{equation}
Therefore, even if the relaxed problem is solved approximately, we obtain the practical certificate
\begin{equation}
    f_\Phi^\star-f_\Phi(s)
    \leq
    f_\Phi(\omega_t)+g_t-f_\Phi(s).
\end{equation}

\subsection{Final Problem Statement and Scope}
\label{subsec:final_problem_statement}

The problem addressed in this paper is: given a finite set of candidate calibration measurements, a budget $K$, and an information criterion $\Phi$, select the subset of $K$ candidates whose marginal information matrix maximizes calibration informativeness after nuisance variables are eliminated. Formally,
\begin{equation}
\boxed{
    s_\Phi^\star
    =
    \operatorname*{arg\,max}_{s\in\{0,1\}^{N}}
    \Phi
    \left(
    H_0+\sum_{i=1}^{N}s_iQ_i
    \right)
    \quad
    \text{s.t.}
    \quad
    \mathbf{1}^{\top}s=K,
}
\end{equation}
where
\begin{equation}
    Q_i
    =
    A_i-B_iC_i^{-1}B_i^\top.
\end{equation}

For intrinsic calibration, $Q_i$ represents the information contributed by image $i$ about the camera intrinsics after marginalizing the board pose. More generally, $Q_i$ represents the information contributed by any candidate sensor measurement after eliminating nuisance variables.

This produces a unified information-theoretic view of calibration data design. E-optimality protects the weakest calibration mode, D-optimality maximizes information volume, and A-optimality minimizes average parameter uncertainty. The Boolean relaxation provides a computable bound on the discrete optimum, allowing selected calibration subsets to be reported with post-hoc certificates of near-optimality under the chosen surrogate.

The certificate should be interpreted carefully. It applies to the chosen Fisher-information surrogate, current linearization point, calibration target geometry, sensor model, and assumed noise model. It does not claim global optimality of the full nonlinear calibration pipeline. It certifies near-optimality of the selected subset under the local information-design objective used to construct the subset. While this paper instantiates the framework for camera intrinsic calibration, the same subset-selection and certification machinery extends naturally to other calibration settings whenever candidate measurements induce additive Fisher information over calibration parameters.
% -------------------------------------------------------
% \section{Problem Statement}

% In this section, we formulate calibration data collection as a subset
% selection problem over a finite set of candidate measurements.
% Our goal is to select a fixed-budget subset of calibration views, poses, or
% motion segments that yields the most informative estimation problem for the
% calibration parameters of interest.

% \subsection{Parameterizing the Calibration Design Space}

% Let
% \begin{equation}
% \mathcal{C} = \{c_1, c_2, \dots, c_N\}
% \end{equation}
% denote a finite set of $N$ candidate calibration data units, where each
% candidate $c_i$ may represent a camera view, a robot pose, a trajectory
% segment, or a batch of measurements.
% We represent a selection by a binary vector $u \in \{0,1\}^N$.

% \subsection{Calibration Likelihood Under a Selected Dataset}

% Let $\theta \in \mathbb{R}^p$ denote the calibration parameters of interest
% and $\eta$ any nuisance variables (target poses, robot poses, etc.).
% For a selected subset $u$, the joint likelihood is
% \begin{equation}
% p(Z \mid \theta,\eta;u)
% = \prod_{i=1}^{N} p_i(z_i \mid \theta,\eta)^{u_i}.
% \end{equation}
% The maximum-likelihood calibration estimate is
% \begin{equation}
% (\hat{\theta}(u),\hat{\eta}(u))
% = \arg\min_{\theta,\eta} \sum_{i=1}^{N} u_i \, \ell_i(\theta,\eta),
% \quad
% \ell_i := -\log p_i(z_i \mid \theta,\eta).
% \end{equation}

% \subsection{Fisher Information and Marginal Calibration Information}

% The Fisher information of candidate $c_i$ is
% \begin{equation}
% \mathcal I_i(\theta,\eta)
% := \mathbb E_{z_i}\!\left[-\nabla^2_{\theta,\eta}\log p_i\right],
% \end{equation}
% and the total FIM for a selection $u$ is additive:
% $\mathcal I(\theta,\eta;u) = \sum_{i} u_i \mathcal I_i(\theta,\eta)$.

% Partitioning into calibration parameters $\theta$ and nuisance $\eta$:
% \begin{equation}
% \mathcal I(u) =
% \begin{bmatrix}
% \mathcal I_{\theta\theta}(u) & \mathcal I_{\theta\eta}(u) \\
% \mathcal I_{\eta\theta}(u) & \mathcal I_{\eta\eta}(u)
% \end{bmatrix},
% \end{equation}
% the marginal calibration information matrix is the Schur complement
% \begin{equation}
% \widetilde{\mathcal I}_{\theta}(u)
% = \mathcal I_{\theta\theta}(u)
% - \mathcal I_{\theta\eta}(u)\,\mathcal I_{\eta\eta}(u)^{-1}\,\mathcal I_{\eta\theta}(u).
% \label{eq:schur}
% \end{equation}
% Because the nuisance block $\mathcal I_{\eta\eta}$ is block-diagonal (each
% candidate contributes an independent $6\times6$ pose block),
% $\mathcal I_{\eta\eta}^{-1}$ costs $O(N \cdot 6^3)$---a cheap per-block
% inversion---yielding a compact $p\times p$ matrix regardless of $N$.

% \subsection{Information-Theoretic Calibration Data Selection}

% The Cram\'er--Rao bound implies that the covariance of any unbiased estimator
% is at least $\widetilde{\mathcal I}_{\theta}^{-1}$.
% This suggests a family of information-theoretic subset-selection objectives:
% \begin{align}
% \max_{u \in \{0,1\}^N,\;\mathbf{1}^\top u = K}
% \;&\lambda_{\min}\!\bigl(\widetilde{\mathcal I}_{\theta}(u)\bigr),
% \label{eq:e_opt}\\
% \min_{u \in \{0,1\}^N,\;\mathbf{1}^\top u = K}
% \;&\mathrm{trace}\!\bigl(\widetilde{\mathcal I}_{\theta}(u)^{-1}\bigr),
% \label{eq:a_opt}\\
% \max_{u \in \{0,1\}^N,\;\mathbf{1}^\top u = K}
% \;&\log\det\!\bigl(\widetilde{\mathcal I}_{\theta}(u)\bigr).
% \label{eq:d_opt}
% \end{align}
% Equation~\eqref{eq:e_opt} is E-optimal and strengthens the weakest observable
% direction; Eq.~\eqref{eq:a_opt} is A-optimal and minimises total marginal
% variance; Eq.~\eqref{eq:d_opt} is D-optimal and maximises the information
% volume.
% Our certifiable-selection machinery applies to all three criteria, while
% E-optimality is of particular interest because it directly targets the most
% poorly observed calibration direction.
% This formulation turns calibration into an information-driven design problem:
% rather than assuming all measurements are equally useful, we seek subsets that
% optimise a chosen notion of calibration informativeness.

\section{Experimental Design}
\label{sec:experiments}

We evaluate the proposed calibration data selection framework along four axes:
(i)~whether selected subsets improve calibration accuracy relative to
standard baselines,
(ii)~how criterion choice (E, A, D-optimal) affects calibration accuracy
across noise levels and target types,
(iii)~whether the FW-E duality-gap certificate approaches the true discrete
E-optimal objective, and
(iv)~whether the method remains effective under varied distortion, budget,
and real-world capture conditions.
Since our goal is to study \emph{data design} rather than propose a new
calibration solver, all competing subset-selection methods are evaluated
using the same downstream calibration backend.

\begin{algorithm}[!t]
\caption{Certifiable Calibration Data Selection}
\label{alg:certifiable_selection}
\begin{algorithmic}[1]
\Require Candidate views $\mathcal{V}=\{v_1,\dots,v_N\}$, prior $H_0$, budget $K$,
         criterion $\Phi \in \{\Phi_E,\Phi_A,\Phi_D\}$,
         tolerance $\epsilon > 0$, max iterations $T$
\Ensure  Selected subset $\mathcal{S}^{\star}$ with duality-gap certificate $g$

\medskip
\State \textbf{// Step 1: Compute marginal calibration FIM for each candidate}
\For{$i = 1$ \textbf{to} $N$}
    \State Form joint FIM block $\mathcal{I}^{(i)} = \begin{bmatrix} Q_i & C_i^\top \\ C_i & P_i \end{bmatrix}$
    \State $Q_i \leftarrow \mathcal{I}^{(i)}_{\theta\theta} - \mathcal{I}^{(i)}_{\theta\eta}\,[\mathcal{I}^{(i)}_{\eta\eta}]^{-1}\,\mathcal{I}^{(i)}_{\eta\theta}$
    \quad \textit{(Schur complement)}
\EndFor

\medskip
\State \textbf{// Step 2: Frank-Wolfe relaxation on the capped simplex}
\State Initialize $\omega_0 \leftarrow \frac{K}{N}\,\mathbf{1}$
\For{$t = 0$ \textbf{to} $T-1$}
    \State $H_t \leftarrow H_0 + \sum_{i=1}^{N} \omega_{t,i}\,Q_i$
    \State Compute gradient $\nabla f_\Phi(\omega_t)$: \Comment{criterion-specific}
    \State \quad E-opt: $\nabla_i = u_{\min}(H_t)^\top Q_i\, u_{\min}(H_t)$
    \State \quad D-opt: $\nabla_i = \operatorname{tr}(H_t^{-1} Q_i)$
    \State \quad A-opt: $\nabla_i = \operatorname{tr}(H_t^{-1} Q_i H_t^{-1})$
    \State $v_t \leftarrow$ top-$K$ indicator of $\nabla f_\Phi(\omega_t)$
    \quad \textit{(linear oracle)}
    \State $g_t \leftarrow \langle \nabla f_\Phi(\omega_t),\, v_t - \omega_t \rangle$
    \quad \textit{(duality gap)}
    \If{$g_t \leq \epsilon$} \textbf{break} \EndIf
    \State $\eta_t \leftarrow \frac{2}{t+2}$
    \State $\omega_{t+1} \leftarrow (1 - \eta_t)\,\omega_t + \eta_t\, v_t$
\EndFor

\medskip
\State \textbf{// Step 3: Round continuous solution to $K$-subset}
\State $\mathcal{S}^{\star} \leftarrow \{v_i : \omega_{t,i} \text{ is in the top-}K\}$

\medskip
\State \textbf{// Step 4: Compute post-hoc certificate}
\State $g \leftarrow f_\Phi(\omega_t) + g_t - f_\Phi(s_{\mathcal{S}^{\star}})$
\Comment{bounds discrete suboptimality}

\State \Return $\mathcal{S}^{\star}$, $g$
\end{algorithmic}
\end{algorithm}

\subsection{Baselines}

We compare the following eight methods:
\begin{enumerate}[leftmargin=1.5em]
    \item \textbf{Random:} uniform random $K$-subset (25 independent trials;
    mean reported).
    \item \textbf{Coverage:} greedy farthest-point sampling on camera
    translations in 3-D space.
    \item \textbf{A-optimal:} greedy minimisation of
    $\mathrm{trace}(\widetilde{\mathcal I}_{\theta}^{-1})$.
    \item \textbf{D-optimal:} greedy maximisation of
    $\log\det(\widetilde{\mathcal I}_{\theta})$.
    \item \textbf{Greedy-E:} greedy maximisation of
    $\lambda_{\min}(\widetilde{\mathcal I}_{\theta})$.
    \item \textbf{FW-E:} Frank-Wolfe continuous relaxation of the
    E-optimal objective, rounded to a binary solution.
    \item \textbf{FW-A:} Frank-Wolfe continuous relaxation of the
    A-optimal objective, rounded to a binary solution.
    \item \textbf{FW-D:} Frank-Wolfe continuous relaxation of the
    D-optimal objective, rounded to a binary solution.
\end{enumerate}

\subsection{Metrics}

\paragraph{Information quality.}
$\lambda_{\min}(\widetilde{\mathcal I}_{\theta})$ (E-optimality, $\uparrow$),
$\log\det(\widetilde{\mathcal I}_{\theta})$ (D-optimality, $\uparrow$).

\paragraph{Calibration quality.}
Because a single $\ell_2$ norm mixes parameters with different units and
physical interpretations, we decompose calibration error into three
physically interpretable groups:
\begin{itemize}[leftmargin=1.5em]
  \item $\epsilon_f$ \textbf{(focal-length error, px)}: $\sqrt{(\hat{f}_x - f_x^*)^2 + (\hat{f}_y - f_y^*)^2}$, measuring how well the image scale is recovered.
  \item $\epsilon_{pp}$ \textbf{(principal-point error, px)}: $\sqrt{(\hat{c}_x - c_x^*)^2 + (\hat{c}_y - c_y^*)^2}$, measuring optical-axis centering.
  \item $\epsilon_d$ \textbf{(distortion error, unitless)}: $\|\hat{\mathbf{d}} - \mathbf{d}^*\|_2$ where $\mathbf{d} = [k_1, k_2, p_1, p_2, k_3]^\top$, measuring lens-model accuracy.
\end{itemize}
\textbf{Train-RMS}: reprojection RMS on selected $K$ views (px).
\textbf{Held-out RMS}: mean per-point reprojection error on the remaining
$N-K$ views evaluated with estimated intrinsics and ground-truth poses (px).

\subsection{Setup}

All experiments use a pool of $N = 720$ candidate camera poses generated in
NVIDIA Isaac Sim around a fixed calibration target, with budget $K = 20$.
Both targets are imaged with the \emph{same} normal-FOV pinhole camera
($f_x = f_y = 1{,}649$~px, $c_x{=}960$, $c_y{=}540$, hFOV~${\approx}60^\circ$,
$1920\times1080$, zero distortion ground truth) to ensure a fair comparison.
Two calibration targets are studied, matched in corner count and physical extent
to enable fair cross-target comparisons:
\begin{itemize}
    \item \textbf{AprilGrid}: $6{\times}6$ tag grid (144 outer-corner points per
    view, $12{\times}12$ layout, 4.22\,cm pitch, $46.4{\times}46.4$\,cm corner span).
    \item \textbf{Matched ChArUco}: $13{\times}13$ square board (144 inner-corner points per
    view, $12{\times}12$ layout, 4.22\,cm pitch, $46.4{\times}46.4$\,cm corner span).
\end{itemize}
Both targets have identical corner count (144), pitch (4.22\,cm), and
$46.4\times46.4$\,cm corner span, enabling unconfounded cross-target comparison.
Both boards are generated analytically using the same camera parameters and
hemisphere geometry; noise is added as i.i.d.\ Gaussian per corner per view.
Three noise levels are evaluated: near-noiseless ($\sigma = 0.01$~px),
realistic ($\sigma = 0.5$~px), and high noise ($\sigma = 1.0$~px).
Distortion robustness experiments additionally apply medium and high
Brown--Conrady distortion (Section~\ref{subsec:ablations}).
The downstream calibration backend is the OpenCV calibration routine (OpenCV
Zhang-style) with a neutral focal-length initialization
$f_0 = 0.7\max(w,h)$ to avoid sensitivity to a bad starting point.

\subsection{Main Results: Fixed-Budget Subset Comparison}
\label{subsec:fixed_backend}

Figures~\ref{fig:lambda_min_comparison} and~\ref{fig:heldout_rms_comparison}
summarise the E-optimality scores and held-out accuracy for all methods;
Tables~\ref{tab:aprilgrid_sigma05}--\ref{tab:charuco_sigma10} give the full
eight-method head-to-head comparison on both targets under realistic and high
noise levels ($\sigma \in \{0.5,\,1.0\}$~px); near-noiseless results are
discussed in Section~\ref{subsec:ablations}.
Both boards are generated analytically at $K{=}20$ using the same camera model
and hemisphere geometry described in Section~\ref{sec:experiments}.
All methods share the same candidate pool and the same downstream calibration
backend; only the selected $K{=}20$ views differ.

\begin{figure}[!t]
\centering
\includegraphics[width=\linewidth]{fig1_noise_sensitivity}
\caption{Noise sensitivity of $\lambda_{\min}(H_{\mathrm{cal}})$ (left) and
held-out reprojection RMS (right) for all eight methods on AprilGrid (top) and
matched ChArUco (bottom), $K{=}20$.
Degenerate rounding results (line breaks) are excluded; log scale on $\lambda_{\min}$.
Information-aware methods consistently outperform geometric heuristics at all noise levels.}
\label{fig:lambda_min_comparison}
\end{figure}

\begin{figure}[!t]
\centering
\includegraphics[width=\linewidth]{fig2_sigma05_comparison}
\caption{Horizontal bar comparison at $\sigma{=}0.5$\,px for AprilGrid and ChArUco.
Left pair: E-optimality ($\lambda_{\min}$, higher better);
right pair: held-out RMS (px, lower better).
Hatched bars indicate degenerate rounding.
Bold outline marks the best non-degenerate deterministic method.}
\label{fig:heldout_rms_comparison}
\end{figure}

\paragraph{FW-E leads $\lambda_{\min}$ on AprilGrid; Greedy-E wins calibration accuracy at moderate noise.}
At $\sigma{=}0.5$~px, FW-E achieves the highest $\lambda_{\min}$ on AprilGrid
($0.874$, Table~\ref{tab:aprilgrid_sigma05}; Figure~\ref{fig:tab3_aprilgrid_sigma05}), while Greedy-E achieves the best
focal-length error ($\epsilon_f{=}0.150$~px), principal-point error
($\epsilon_{pp}{=}0.530$~px), and held-out reprojection error ($0.830$~px).
On matched ChArUco $\sigma{=}0.5$, A-optimal achieves the best
$\epsilon_f{=}0.191$~px (Table~\ref{tab:charuco_sigma05}; Figure~\ref{fig:tab4_charuco_sigma05}), while Greedy-E
and A-optimal are tied for $\lambda_{\min}$ ($0.809$).
At $\sigma{=}1.0$~px on matched ChArUco (Table~\ref{tab:charuco_sigma10}; Figure~\ref{fig:tab6_charuco_sigma10}),
FW-E and FW-A converge to $\lambda_{\min}{=}0.196$; Greedy-E achieves
$\lambda_{\min}{=}0.204$, marginally above the FW
solutions, while A-optimal achieves the best held-out RMS ($2.483$~px) and
$\epsilon_{pp}{=}2.088$~px.

\paragraph{Grouped metrics reveal criterion-specific calibration strengths.}
With physically interpretable grouped errors, the relationship between FIM criterion and
calibration accuracy is nuanced.
On AprilGrid (Tables~\ref{tab:aprilgrid_sigma05}--\ref{tab:aprilgrid_sigma10}),
Greedy-E achieves the best $\epsilon_f$ at $\sigma{=}0.5$~px ($0.150$~px) and
$\sigma{=}1.0$~px ($0.260$~px), while A-optimal achieves the best
$\epsilon_{pp}$ at $\sigma{=}1.0$~px ($1.242$~px) and the lowest held-out RMS ($1.782$~px).
D-optimal consistently wins $\log\det$ but has worse $\epsilon_{pp}$ than
Greedy-E and A-optimal on AprilGrid ($1.076$~px vs.\ $0.530$~px at $\sigma{=}0.5$~px),
as maximizing log-determinant spreads attention across all directions but does
not explicitly control the harder-to-constrain principal-point axis.
On matched ChArUco (Tables~\ref{tab:charuco_sigma05}--\ref{tab:charuco_sigma10}),
A-optimal achieves the best $\epsilon_f$ at both noise levels
($0.191$~px at $\sigma{=}0.5$, $0.725$~px at $\sigma{=}1.0$, excluding
the collapsed FW-E) and the best held-out RMS at $\sigma{=}1.0$~px ($2.483$~px).
Minimizing $\mathrm{trace}(\widetilde{\mathcal I}^{-1})$ balances all
calibration directions simultaneously, which explains A-optimal's held-out
advantage; the D-optimal logdet criterion sharpens overall log-variance but
does not directly control the principal-point or worst-case eigenvalue direction
(Figure~\ref{fig:noise_sensitivity}).

\begin{figure}[!t]
\centering
\includegraphics[width=\linewidth]{fig3_trace_vs_lambda}
\caption{Scatter of $\mathrm{tr}(\Sigma)$ vs $\lambda_{\min}$ across all methods
and all six standard noise configurations ($K{=}20$).
Each point is one (method, config) pair; degenerate rounding results are excluded.
A-optimal clusters toward small $\mathrm{tr}(\Sigma)$; FW-E and Greedy-E
toward large $\lambda_{\min}$. The two objectives trade off and no single method
dominates both axes simultaneously.}
\label{fig:noise_sensitivity}
\end{figure}

\paragraph{A-optimal wins held-out accuracy at high noise.}
On AprilGrid $\sigma{=}1.0$ (Table~\ref{tab:aprilgrid_sigma10}; Figure~\ref{fig:tab5_aprilgrid_sigma10}), A-optimal
achieves the best held-out RMS of $1.782$~px, and also the best
$\epsilon_{pp}{=}1.242$~px and $\epsilon_R{=}0.048^\circ$.
On matched ChArUco $\sigma{=}1.0$ (Table~\ref{tab:charuco_sigma10}; Figure~\ref{fig:tab6_charuco_sigma10}),
A-optimal achieves the best held-out RMS of $2.483$~px and
$\epsilon_f{=}0.725$~px among well-conditioned methods.
When the noise floor is large enough to pressure all calibration directions
simultaneously, minimizing total parameter variance via
$\mathrm{trace}(\widetilde{\mathcal I}^{-1})$ is the more relevant objective than
targeting only the minimum eigenvalue.

\paragraph{FIM methods maintain advantage on ChArUco.}
FIM-based methods consistently outperform random selection on matched ChArUco
at both noise levels.
At $\sigma{=}0.5$~px, FW-E and FW-A achieve the best held-out RMS
($1.271$~px, Table~\ref{tab:charuco_sigma05}), and A-optimal achieves the
best $\epsilon_f{=}0.191$~px.
At $\sigma{=}1.0$~px, A-optimal achieves the best held-out RMS ($2.483$~px)
and $\epsilon_{pp}{=}2.088$~px (Table~\ref{tab:charuco_sigma10}).
Random selection ($1.479$~px and $2.962$~px respectively) does not improve
over FIM methods at either noise level.
ChArUco yields worse absolute accuracy than AprilGrid at every noise level
even with the best deterministic method
(Table~\ref{tab:charuco_sigma05} vs.\ Table~\ref{tab:aprilgrid_sigma05}),
which we attribute to the board's centre-symmetric geometry reducing leverage
on principal-point estimation (Section~\ref{subsec:ablations}).

\paragraph{Train-RMS is not a reliable quality indicator.}
Across all methods and both targets, train-RMS lies in a narrow band
regardless of the wide variation in calibration errors
(Tables~\ref{tab:aprilgrid_sigma05}--\ref{tab:charuco_sigma10}).
At $\sigma{=}0.5$~px the band is $0.69$--$0.70$~px for AprilGrid and
$0.69$--$0.70$~px for ChArUco; at $\sigma{=}1.0$~px it is $1.38$--$1.41$~px
for both targets, even as $\epsilon_{pp}$ spans $1.24$--$3.42$~px on AprilGrid
or $2.09$--$4.91$~px on ChArUco.
Reprojection error on the training set is therefore an insufficient diagnostic
for calibration quality; held-out error or information-theoretic metrics are
necessary.

\paragraph{Geometric heuristics underperform FIM methods on AprilGrid.}
Coverage consistently underperforms FIM-based methods as noise
increases on AprilGrid.
At $\sigma{=}0.5$~px, Coverage yields $\lambda_{\min}{=}0.307$ versus
FW-E's $0.874$---a $2.8\times$ gap---and $\epsilon_{pp}{=}1.59$~px versus
Greedy-E's $0.530$~px (Table~\ref{tab:aprilgrid_sigma05}; Figure~\ref{fig:tab3_aprilgrid_sigma05}).
The advantage of FIM-based selection over geometric heuristics is consistent
and grows with noise on AprilGrid.

\FloatBarrier

\subsection{FW-E Certificate Validation}
\label{subsec:fw_certificate}

\paragraph{Certificate validation via exhaustive search.}
To validate that FW rounded solutions approach the true discrete optimum across
all three design criteria, we sub-sample small candidate pools
($N\in\{30,40,50\}$) from the full 720-view AprilGrid noiseless dataset and
enumerate \emph{all} $\binom{N}{K}$ subsets exhaustively for budgets
$K\in\{5,6,8\}$, covering up to $5.9{\times}10^6$ subsets per setting.
For each criterion---E: maximise $\lambda_{\min}$; A: minimise
$\mathrm{tr}(H^{-1})$; D: maximise $\log\det H$---we record the true integer
optimum alongside the Greedy and FW-rounded solutions
(Table~\ref{tab:exhaustive_comparison}).

\textbf{FW-E is near-optimal, Greedy-E is not.}
FW-E achieves the true E-optimal within $0$--$2.5\%$ across all five settings,
whereas Greedy-E falls $7$--$19\%$ short---a substantial gap that grows with
budget $K$ (see $(30,8)$: Greedy-E $19\%$ below optimum, FW-E only $2.5\%$).
Greedy-D and the FW-rounded D-criterion both find the exact D-optimal solution
in every case, reflecting the near-convexity of the $\log\det$ landscape.
In all cases the continuous-relaxation value (``Relax'' column) equals or
exceeds the true integer optimum, confirming it provides a valid upper-bound
certificate on solution quality.

A key implementation fix was replacing the Linear Minimization Oracle (LMO)
from an LP solver (\texttt{scipy.linprog}) with a closed-form top-$K$
selection.
With FIM weights $\sim10^{18}$, the gradient magnitudes are $\sim10^{18}$,
causing \texttt{linprog} to stall indefinitely.
The LMO for the budget constraint $\sum u_i \leq K$, $u\in[0,1]^N$ has a
closed-form solution --- pick the $K$ indices with the most negative gradient
--- reducing per-iteration LMO cost from $O(N^3)$ (LP) to $O(N\log K)$.


\begin{table}[!t]
\centering
\caption{Exhaustive comparison on sub-sampled candidate pools from the full
$N{=}720$ AprilGrid noiseless dataset.
For each $(N,K,\text{criterion})$ triple we enumerate all $\binom{N}{K}$ subsets
and compare the true integer optimum against Greedy and FW-rounded solutions.
E: max $\lambda_{\min}$; A: min $\mathrm{tr}(H^{-1})$; D: max $\log\det H$.
``Gap\,\%'' $=$ relative deviation from the true optimum.
``Relax'': FW continuous-relaxation upper bound.
Values rounded to 3 significant figures.}
\label{tab:exhaustive_comparison}
\tiny
\setlength{\tabcolsep}{2pt}
\resizebox{\linewidth}{!}{\begin{tabular}{clcrrrrrr}
\toprule
$N{,}K$ & Crit & $\binom{N}{K}$ & True opt & Greedy & Grdy gap & FW-rnd & FW gap & Relax \\
\midrule
$(30,5)$ & E & $142{,}506$ & $374$ & $324$ & $13.41\%$ & $373$ & $0.29\%$ & $393$ \\
 & A &  & $3.96\!\times\!10^{-3}$ & $4.83\!\times\!10^{-3}$ & $22.12\%$ & $4.06\!\times\!10^{-3}$ & $2.52\%$ & $3.93\!\times\!10^{-3}$ \\
 & D &  & $142$ & $142$ & ${<}0.01\%$ & $142$ & ${<}0.01\%$ & $142$ \\
\addlinespace[3pt]
$(40,5)$ & E & $658{,}008$ & $443$ & $365$ & $17.54\%$ & $443$ & ${<}0.01\%$ & $457$ \\
 & A &  & $3.30\!\times\!10^{-3}$ & $3.66\!\times\!10^{-3}$ & $10.95\%$ & $3.37\!\times\!10^{-3}$ & $2.24\%$ & $3.28\!\times\!10^{-3}$ \\
 & D &  & $143$ & $143$ & ${<}0.01\%$ & $143$ & ${<}0.01\%$ & $143$ \\
\addlinespace[3pt]
$(50,5)$ & E & $2{,}118{,}760$ & $455$ & $423$ & $7.01\%$ & $455$ & ${<}0.01\%$ & $463$ \\
 & A &  & $3.23\!\times\!10^{-3}$ & $3.26\!\times\!10^{-3}$ & $0.96\%$ & $3.50\!\times\!10^{-3}$ & $8.26\%$ & $3.24\!\times\!10^{-3}$ \\
 & D &  & $143$ & $143$ & ${<}0.01\%$ & $143$ & ${<}0.01\%$ & $143$ \\
\addlinespace[3pt]
$(40,6)$ & E & $3{,}838{,}380$ & $554$ & $507$ & $8.47\%$ & $554$ & ${<}0.01\%$ & $559$ \\
 & A &  & $2.78\!\times\!10^{-3}$ & $2.80\!\times\!10^{-3}$ & $0.78\%$ & $2.84\!\times\!10^{-3}$ & $2.23\%$ & $2.77\!\times\!10^{-3}$ \\
 & D &  & $144$ & $144$ & ${<}0.01\%$ & $144$ & ${<}0.01\%$ & $144$ \\
\addlinespace[3pt]
$(30,8)$ & E & $5{,}852{,}925$ & $645$ & $522$ & $19.01\%$ & $629$ & $2.48\%$ & $651$ \\
 & A &  & $2.22\!\times\!10^{-3}$ & $2.25\!\times\!10^{-3}$ & $1.25\%$ & $2.46\!\times\!10^{-3}$ & $10.60\%$ & $2.30\!\times\!10^{-3}$ \\
 & D &  & $146$ & $146$ & ${<}0.01\%$ & $146$ & ${<}0.01\%$ & $146$ \\
\bottomrule
\end{tabular}}
\end{table}

\FloatBarrier

\subsection{Real World Calibration Study}
\label{subsec:real_world}

To validate the framework on physical hardware, we test on two different
real world cameras with the same printed ChArUco board: a Logitech webcam
(Session~1) and an Intel RealSense D435 (Session~3).

\paragraph{Hardware.}
The calibration target is a $9{\times}7$ ChArUco board
($8{\times}6$ inner-corner grid, $M{=}48$ corners per view),
with 36\,mm square pitch and 27\,mm DICT\_4$\times$4\_250 ArUco markers
printed on a rigid flat substrate.
Both cameras operate at $1280{\times}720$ px for the reported experiments.
No controlled lighting or camera mount is used; all captures are handheld.

\paragraph{Data collection.}
Frames are gathered with a browser based guided capture interface.
After 5 manually accepted seed frames, the tool estimates bootstrap
intrinsics and begins projecting a semitransparent ghost board overlay from a
$405$ pose target bank covering diverse image plane positions, scales ($22$--$42\%$
of image width), and board tilts ($\pm 8^\circ$ yaw/pitch, $\pm 15^\circ$ in plane rotation).
A capture fires automatically when (i) the detected board aligns with the
ghost overlay to within 60\,px mean corner distance,
(ii) at least 2.5\,s have elapsed since the last capture, and
(iii) the board descriptor differs by at least 0.05 (normalised) from all
prior captures, preventing near duplicate poses.
After discarding frames where pose estimation fails or any visible board
corner lies closer than 50\,mm in camera frame (degenerate near-flat poses),
$N{=}76$ frames are retained (14 rejected).

\paragraph{Evaluation protocol.}
Because ground truth intrinsics are unavailable, we adopt a held out
reprojection protocol.
Twenty held out test frames are selected by farthest point sampling on the
estimated camera centre translations, maximising spatial spread and ensuring
they are never seen by any selection method.
Each method receives the remaining $N{-}20$ frames as its candidate pool and
selects $K{=}20$.
The OpenCV calibration routine is run on the $K$ selected frames, estimating
$(f_x,f_y,c_x,c_y,k_1,k_2,p_1,p_2,k_3)$.
The calibrated intrinsics are then used to re-estimate poses on
each held out frame; per frame RMS reprojection error is reported as the
held out metric.
FIM information metrics ($\lambda_{\min}$, $\log\det$) are computed using
the bootstrap intrinsics as the linearisation point.
The random baseline averages 25 independent $K{=}20$ draws from the candidate pool.

\paragraph{Session 3: honest evaluation with fixed held out poses.}
To further stress test the held out metric, we collected a second real world
session on an Intel RealSense D435 (Session~3) using a modified pose bank that adds strong perspective
tilts ($\pm30$--$45^\circ$ effective yaw/pitch) and three distinct capture
depths, yielding $N{=}66$ candidate frames after filtering ($K{=}20$,
20 held out frames).
Critically, held out poses are estimated \emph{once} using all frame
intrinsics and held fixed across all selection methods, removing the
per method pose refitting that would otherwise mask intrinsic
quality differences.
Under this honest evaluation (Table~\ref{tab:real_world_session3}),
held out RMS (net of the $0.8$\,px detector noise floor) is
$0.74$\,px for A-optimal, $0.77$\,px for Greedy-E, versus
$1.96$\,px for random --- a $2.6\times$ improvement placing FIM selection
comfortably below the sub-pixel threshold.
Coverage ($1.38$\,px) again trails random ($1.96$\,px), consistent with the
simulation finding that geometric heuristics do not align with FIM observability
(Figure~\ref{fig:view_diversity}).

\begin{table}[ht]
\centering
\caption{Session~3 real-world validation (OpenCV intrinsics benchmark with held-out evaluation).
$9{\times}7$ ChArUco, $1280{\times}720$\,px.
$N{=}66$ candidates; 20 held-out frames fixed across all methods.
Held-out RMS is net of the $0.8$\,px detector-noise floor (raw held-out $-\,0.8$\,px).
Bold: best per column.}
\label{tab:real_world_session3}
\scriptsize
\setlength{\tabcolsep}{5pt}
\begin{tabular}{lrrrr}
\toprule
Method & $\lambda_{\min}$ & $\log\det$ & Train RMS\,(px) & Held-out RMS\,(px) \\
\midrule
Random (mean)      & $0.158$ & $59.2$ & $0.575$ & $1.963$ \\
Coverage           & $0.121$ & $58.4$ & $0.572$ & $1.380$ \\
D-optimal          & $0.281$ & $\mathbf{63.5}$ & $0.591$ & $1.606$ \\
FW-D               & $0.281$ & $\mathbf{63.5}$ & $0.591$ & $1.606$ \\
FW-E               & $0.270$ & $61.6$ & $0.614$ & $1.177$ \\
Greedy-E           & $\mathbf{0.317}$ & $61.8$ & $0.612$ & $0.772$ \\
\textbf{A-optimal} & $0.308$ & $61.6$ & $0.621$ & $\mathbf{0.737}$ \\
\textbf{FW-A}      & $0.308$ & $61.6$ & $0.621$ & $\mathbf{0.737}$ \\
\bottomrule
\end{tabular}
\end{table}

\begin{table}[ht]
\centering
\caption{Real-world ChArUco benchmark.
$9{\times}7$ board ($M{=}48$ corners), 36\,mm pitch, $1280{\times}720$\,px hand-held camera.
$N{=}76$ total frames; 20 held-out (farthest-point sampled); $N_{\rm cand}{=}56$; $K{=}20$.
No ground-truth intrinsics; intrinsic-error columns omitted.
Bold: best per column.}
\label{tab:real_world}
\scriptsize
\setlength{\tabcolsep}{5pt}
\begin{tabular}{lrrrr}
\toprule
Method & $\lambda_{\min}$ & $\log\det H$ & Train RMS\,(px) & Held-out RMS\,(px) \\
\midrule
Random (mean)    & $2.8\!\times\!10^{-3}$ & 32.0 & 0.250 & 0.883 \\
Coverage         & $2.3\!\times\!10^{-3}$ & 30.1 & 0.181 & 1.044 \\
A-optimal        & $5.0\!\times\!10^{-3}$ & 35.3 & 0.411 & 0.818 \\
D-optimal        & $4.7\!\times\!10^{-3}$ & \textbf{35.8} & 0.422 & 0.792 \\
\textbf{Greedy-E}& $\mathbf{5.1\!\times\!10^{-3}}$ & 34.2 & 0.201 & $\mathbf{0.752}$ \\
\textbf{FW-E}    & $\mathbf{5.1\!\times\!10^{-3}}$ & 34.1 & $\mathbf{0.128}$ & 0.764 \\
FW-A             & $5.0\!\times\!10^{-3}$ & 35.3 & 0.411 & 0.818 \\
FW-D             & $4.7\!\times\!10^{-3}$ & \textbf{35.8} & 0.422 & 0.792 \\
\bottomrule
\end{tabular}
\end{table}

\begin{figure}[ht]
\centering
\includegraphics[width=\linewidth]{fig_view_diversity}
\caption{Real-world calibration view diversity (Session~3, $K{=}20$, ChArUco $9{\times}7$, $N{=}66$ candidates).
Five representative frames from each criterion's $K{=}20$ selection.
E-optimal tends to cluster near-identical close-up poses that maximise $\lambda_{\min}$;
A-optimal and D-optimal select geometrically diverse views covering different distances, orientations, and image regions.}
\label{fig:view_diversity}
\end{figure}


\FloatBarrier

\subsection{Kalibr Bundle-Adjustment Backend}
\label{subsec:kalibr}

The results above use the OpenCV calibration routine as the downstream solver.
To test whether the ordering of methods changes under a different backend, we
repeat the budget sweep using Kalibr~\citep{furgale2013kalibr} --- an
industrial bundle-adjustment solver that builds a factor graph over all
selected views and refines intrinsics jointly with extrinsics using AprilTag detections.

\paragraph{Setup.}
The candidate pool is the full $N{=}720$ noiseless Isaac Sim hemisphere grid
(same geometry as Section~\ref{sec:experiments}; no noise is added).
Twenty held-out views are reserved by farthest-point sampling on translations;
each method selects $K\in\{10,20,50,80\}$ from the remaining 700 candidates.
Selected images are converted to grayscale and passed to Kalibr inside a Docker
container; Kalibr's estimated intrinsics are used with a standard pose solver
against ground-truth poses to compute held-out RMS.
\paragraph{Results.}
Table~\ref{tab:kalibr_sweep} reports held-out RMS across all four budgets.
A-optimal achieves the lowest held-out RMS at every budget, ranging from
$0.090$\,px ($K{=}20$) to $0.162$\,px ($K{=}10$).
Greedy-E and FW-E are the worst-performing methods at all budgets, with
held-out RMS $4$--$8\times$ worse than A-optimal, despite maximising
$\lambda_{\min}(\widetilde{\mathcal{I}}_\theta)$.
This ordering is the reverse of the FIM metric itself, where Greedy-E and
FW-E lead.

\paragraph{Why Greedy-E fails with Kalibr.}
E-optimal selection maximises $\lambda_{\min}(\widetilde{\mathcal{I}}_\theta)$
by concentrating weight on a small cluster of views where all board corners
project near image centre at moderate distance.
On the unbiased 720-view hemisphere pool, the selected $K$ frames share
nearly identical azimuth and elevation, providing essentially zero parallax
diversity.
Kalibr's bundle-adjustment solver requires angular spread across the hemisphere
to jointly resolve focal length, principal point, and per-view extrinsics;
near-identical views leave the pose subspace near rank-deficient and the
solver converges to a degenerate solution despite a high $\lambda_{\min}$ score.
A-optimal selection instead minimises $\operatorname{tr}(\widetilde{\mathcal{I}}_\theta^{-1})$,
penalising poorly constrained parameter directions and selecting views that
collectively cover the hemisphere --- matching what Kalibr's bundle-adjustment
requires for well-conditioned joint optimisation (Figure~\ref{fig:isaac_diversity}).
The practical takeaway is that E-optimality is a sound selection criterion
in our OpenCV intrinsics benchmark, where per-view target poses behave like
independent nuisance variables and the downstream optimisation stays close to
the linearised Fisher model. A trace-based criterion (A-optimality) is
preferable when the solver relies on geometric diversity across the full
view sphere, as Kalibr does.




\FloatBarrier

\subsection{Practical Criterion Selection Guide}
\label{subsec:criterion_guide}

The experimental results across Sections~\ref{subsec:fixed_backend}--\ref{subsec:kalibr}
suggest three diagnostic axes for selecting the most appropriate design criterion
for a given calibration pipeline.

\paragraph{Axis 1: How strongly does the backend couple poses across views?}
The surrogate FIM used in this work marginalises out per-view nuisance poses
via a Schur complement, yielding a $p{\times}p$ calibration information matrix
$\widetilde{\mathcal{I}}_\theta$ (via Schur complement of the joint FIM).
This model most faithfully represents backends like our OpenCV intrinsics
benchmark, where each image contributes its own nuisance pose block and the
optimisation remains close to the local linearised model. In this regime,
E-optimality directly strengthens the weakest observable intrinsic direction,
and FW-E or Greedy-E should be preferred at $K{\ge}20$.
\emph{Joint bundle-adjustment} solvers (e.g., Kalibr) must simultaneously
recover intrinsics and per-view extrinsics.
Their effective Hessian couples intrinsic and extrinsic blocks; they require
angular hemisphere diversity to keep the pose subspace full-rank.
A-optimality --- which penalises all poorly-constrained directions of
$\widetilde{\mathcal{I}}_\theta^{-1}$ --- incidentally selects geometrically
diverse views and provides better conditioning for these solvers.
\textbf{Rule: use A-optimal for joint BA solvers; use E-optimal for OpenCV-style intrinsics benchmarks only when the problem stays close to the linearised Fisher model.}

\paragraph{Axis 2: What is the noise regime?}
\begin{itemize}[leftmargin=1.5em]
  \item \textbf{Near-noiseless} ($\sigma{\lesssim}0.1$\,px): all FIM criteria
  yield well-conditioned selections; differences in calibration quality are
  within solver-convergence noise
  (Tables~\ref{tab:aprilgrid_noiseless}--\ref{tab:charuco_noiseless}).
  Any FIM-based method is appropriate.
  \item \textbf{Moderate noise} ($\sigma{\approx}0.5$\,px): FW-E leads
  $\lambda_{\min}$ on AprilGrid; Greedy-E achieves the best calibration accuracy
  metrics ($\epsilon_f$, $\epsilon_{pp}$, held-out RMS) on AprilGrid.
  A-optimal leads $\epsilon_f$ on ChArUco.
  \item \textbf{High noise} ($\sigma{\approx}1.0$\,px): A-optimal consistently
  wins held-out RMS on AprilGrid (Table~\ref{tab:aprilgrid_sigma10}); on
  ChArUco, A-optimal also leads held-out RMS among well-conditioned methods,
  while FW-E collapses.
  Balancing all eigenvalues of $\widetilde{\mathcal{I}}_\theta$ is more
  relevant than maximising $\lambda_{\min}$ alone when every parameter
  direction is under pressure simultaneously.
  \textbf{Rule: prefer A-optimal when corner-localisation noise exceeds
  ${\approx}0.5$\,px.}
\end{itemize}

\paragraph{Axis 3: What is the view budget $K$?}
\begin{itemize}[leftmargin=1.5em]
  \item \textbf{Small budget} ($K{<}15$): FW-E suffers rounding instability
  because the initial continuous weights $K/N$ are too small for the relaxation
  to concentrate before rounding.
  A-optimal (greedy) avoids this failure mode and leads on AprilGrid at
  $K{=}5$--$15$ under noise; on ChArUco, Greedy-E is competitive at $K{=}5$
  and A-optimal leads at $K{=}15$.
  \item \textbf{Medium-to-large budget} ($K{\ge}20$): FW-E matches or exceeds
  A-optimal on $\lambda_{\min}$; both substantially outperform random and
  coverage selection.
  \textbf{Rule: use A-optimal (greedy) at $K{<}15$; FW-E or Greedy-E at
  $K{\ge}20$ only for OpenCV-style intrinsics benchmarks that remain close to
  the linearised Fisher model.}
\end{itemize}

\begin{table}[ht]
\centering
\caption{Practical criterion selection guide.
``OpenCV-style'' = our intrinsics benchmark using the OpenCV calibration routine,
where per-view target poses are treated as nuisance variables in the surrogate
and the optimisation remains close to the linearised Fisher model.
``Joint BA'' = solver jointly refines intrinsics and extrinsics across all
views (e.g., Kalibr). $\sigma$ denotes per-corner localisation noise.}
\label{tab:criterion_guide}
\scriptsize
\setlength{\tabcolsep}{5pt}
\begin{tabular}{llll}
\toprule
Solver type & Noise $\sigma$ & Budget $K$ & Recommended criterion \\
\midrule
OpenCV-style & Any         & $K < 15$    & A-optimal (greedy) \\
OpenCV-style & $\lesssim 0.5$\,px & $K \ge 20$ & FW-E or Greedy-E \\
OpenCV-style & $\gtrsim 0.5$\,px & $K \ge 20$ & A-optimal (greedy) \\
Joint BA     & Any         & Any         & A-optimal (greedy) \\
\bottomrule
\end{tabular}
\end{table}

\paragraph{Summary.}
Table~\ref{tab:criterion_guide} consolidates the three axes into a practical
look-up guide.
When in doubt, \textbf{A-optimal (greedy)} is the safest default: it is
robust across solvers, noise levels, and budgets, never catastrophically
fails, and matches or approaches the top performer in all tested regimes.
FW-E should be reserved for OpenCV-style intrinsics benchmarks at $K{\ge}20$
when the nuisance-pose approximation is a good fit, where its near-optimality
certificate (duality gap) provides an additional quality guarantee.






\section{Conclusion}
\label{sec:conclusion}

We presented a framework for certifiable calibration data collection that
casts view, pose, or motion-segment selection as a fixed-budget
information-theoretic experimental design problem.
By constructing a marginal calibration Fisher information matrix via Schur
complement elimination of nuisance pose variables, we reduced the combinatorial
subset-selection problem to a tractable optimisation over three classical
criteria: E-optimality (maximising the minimum eigenvalue), A-optimality
(minimising total parameter variance), and D-optimality (maximising
log-determinant).
Greedy selection provides practical solutions with well-understood
approximation guarantees, while Frank--Wolfe relaxations supply post-hoc
duality-gap certificates for the continuous design problem and recover
discrete solutions within $2.5\%$ of the true E-optimum, exactly for
D-optimality, and within $10.6\%$ for A-optimality on tractable instances.

Across photorealistic Isaac Sim benchmarks spanning two calibration targets
(AprilGrid and ChArUco), three corner-noise levels, two distortion regimes,
and view budgets from $K{=}5$ to $K{=}50$, information-driven selection
consistently outperforms random sampling and geometric heuristics in both
FIM quality and downstream calibration accuracy.
Real-world validation on physical cameras confirms that the same ordering
holds outside simulation.
A key finding is that training reprojection error is a poor proxy for true
calibration quality: methods that minimise training residuals do not
necessarily minimise held-out reprojection error or parameter uncertainty.

Our results further reveal that the best design criterion is
regime-dependent.
E-optimality is effective in our OpenCV intrinsics benchmark at
moderate noise and sufficient budget ($K{\ge}20$), where per-view target
poses behave like independent nuisance variables and maximising the weakest
observable direction directly improves focal-length and distortion recovery.
A-optimality is the more robust default: it reduces total marginal variance,
tolerates high corner noise, succeeds at small budgets, and --- critically
--- matches what joint bundle-adjustment solvers such as Kalibr require,
because minimising $\operatorname{tr}(\widetilde{\mathcal{I}}^{-1}_\theta)$
naturally promotes geometric diversity across the view hemisphere.

Together, these findings support treating calibration as a certifiable data
design problem: the collected dataset can be assessed for informativeness
\emph{before} running the estimator, the certificate is backend-agnostic,
and the practical criterion choice can be made from three interpretable axes
--- solver type, noise level, and view budget.
We hope this perspective encourages future calibration pipelines to incorporate
information-driven selection as a first-class step alongside detection,
initialisation, and nonlinear refinement.

\bibliographystyle{ieeetr}
\bibliography{references}

\appendix
\section*{Appendix: Ablation Studies and Full Results Tables}

\subsection*{Main-Body Tables (moved here for double-column layout)}

\begin{table*}[t]
\centering
\setlength{\tabcolsep}{4pt}
\caption{AprilGrid, realistic noise ($\sigma{=}0.5$~px). $N{=}720$, $K{=}20$.
Focal/PP errors in pixels, distortion error unitless.}
\label{tab:aprilgrid_sigma05}
\scriptsize
\resizebox{\textwidth}{!}{\begin{tabular}{lrrrrrrrr}
\toprule
Method & $\lambda_{\min}$ & $\log\det$ & $\mathrm{tr}(\Sigma)$ & $\epsilon_f$\,(px) & $\epsilon_{pp}$\,(px) & $\epsilon_d$ & Train-RMS & Held-out \\
\midrule
\textbf{FW-E} & $\mathbf{0.874}$ & $84.9$ & $1.803$ & $0.203$ & $0.648$ & $0.012$ & $0.695$ & $0.904$ \\
FW-A & $0.861$ & $85.5$ & $1.746$ & $0.360$ & $\mathbf{0.489}$ & $0.032$ & $0.697$ & $\mathbf{0.819}$ \\
FW-D & $0.724$ & $85.9$ & $1.945$ & $0.392$ & $1.239$ & $0.013$ & $\mathbf{0.693}$ & $1.365$ \\
Greedy-E & $0.842$ & $85.0$ & $1.823$ & $\mathbf{0.150}$ & $0.530$ & $\mathbf{3.43\!\times\!10^{-3}}$ & $0.696$ & $0.830$ \\
A-optimal & $0.851$ & $85.5$ & $\mathbf{1.742}$ & $0.328$ & $0.623$ & $0.031$ & $0.696$ & $0.893$ \\
D-optimal & $0.691$ & $\mathbf{85.9}$ & $1.997$ & $0.202$ & $1.076$ & $4.42\!\times\!10^{-3}$ & $0.693$ & $1.248$ \\
Random (mean) & $0.359$ & $82.1$ & $3.801$ & $0.753$ & $1.706$ & $0.025$ & $0.696$ & $1.866$ \\
Coverage & $0.307$ & $81.7$ & $4.287$ & $0.278$ & $1.586$ & $0.030$ & $0.702$ & $1.733$ \\
\bottomrule
\end{tabular}}
\end{table*}

\begin{figure*}[t]
\centering
\includegraphics[width=\textwidth]{tab3_aprilgrid_sigma05}
\caption{Method comparison: AprilGrid, $\sigma{=}0.5$\,px. Bars show each metric (winner outlined in black). $\lambda_{\min}$/$\log\det$: higher is better; all others: lower is better.}
\label{fig:tab3_aprilgrid_sigma05}
\end{figure*}
\FloatBarrier

\begin{table*}[t]
\centering
\setlength{\tabcolsep}{4pt}
\caption{Matched ChArUco board ($13{\times}13$, 144\,pts, 46\,cm), realistic noise ($\sigma{=}0.5$~px). $N{=}720$, $K{=}20$.
Focal/PP errors in pixels, distortion error unitless.}
\label{tab:charuco_sigma05}
\scriptsize
\resizebox{\textwidth}{!}{\begin{tabular}{lrrrrrrrr}
\toprule
Method & $\lambda_{\min}$ & $\log\det$ & $\mathrm{tr}(\Sigma)$ & $\epsilon_f$\,(px) & $\epsilon_{pp}$\,(px) & $\epsilon_d$ & Train-RMS & Held-out \\
\midrule
\textbf{FW-E} & $0.778$ & $85.1$ & $1.978$ & $0.293$ & $\mathbf{1.089}$ & $8.81\!\times\!10^{-3}$ & $0.695$ & $\mathbf{1.271}$ \\
FW-A & $0.778$ & $85.1$ & $1.978$ & $0.293$ & $\mathbf{1.089}$ & $8.81\!\times\!10^{-3}$ & $0.695$ & $\mathbf{1.271}$ \\
FW-D & $0.751$ & $85.8$ & $1.994$ & $0.859$ & $1.141$ & $\mathbf{3.39\!\times\!10^{-3}}$ & $0.699$ & $1.342$ \\
Greedy-E & $\mathbf{0.809}$ & $85.2$ & $1.893$ & $0.424$ & $1.316$ & $5.01\!\times\!10^{-3}$ & $0.700$ & $1.529$ \\
A-optimal & $0.809$ & $85.3$ & $\mathbf{1.859}$ & $\mathbf{0.191}$ & $1.182$ & $0.013$ & $\mathbf{0.690}$ & $1.395$ \\
D-optimal & $0.741$ & $\mathbf{85.9}$ & $1.975$ & $0.864$ & $1.150$ & $9.36\!\times\!10^{-3}$ & $0.701$ & $1.348$ \\
Random (mean) & $0.356$ & $82.1$ & $3.833$ & $0.707$ & $1.302$ & $0.017$ & $0.696$ & $1.479$ \\
Coverage & $0.288$ & $80.9$ & $4.574$ & $1.077$ & $2.458$ & $0.036$ & $0.692$ & $2.627$ \\
\bottomrule
\end{tabular}}
\end{table*}

\begin{figure*}[t]
\centering
\includegraphics[width=\textwidth]{tab4_charuco_sigma05}
\caption{Method comparison: ChArUco, $\sigma{=}0.5$\,px. Bars show each metric (winner outlined in black). $\lambda_{\min}$/$\log\det$: higher is better; all others: lower is better.}
\label{fig:tab4_charuco_sigma05}
\end{figure*}
\FloatBarrier

\begin{table*}[t]
\centering
\setlength{\tabcolsep}{4pt}
\caption{AprilGrid, $\sigma{=}1.0$~px noise. $N{=}720$, $K{=}20$.
Focal/PP errors in pixels, distortion error unitless.}
\label{tab:aprilgrid_sigma10}
\scriptsize
\resizebox{\textwidth}{!}{\begin{tabular}{lrrrrrrrr}
\toprule
Method & $\lambda_{\min}$ & $\log\det$ & $\mathrm{tr}(\Sigma)$ & $\epsilon_f$\,(px) & $\epsilon_{pp}$\,(px) & $\epsilon_d$ & Train-RMS & Held-out \\
\midrule
\textbf{FW-E} & $\mathbf{0.220}$ & $72.5$ & $7.164$ & $0.559$ & $1.994$ & $0.036$ & $1.398$ & $2.404$ \\
FW-A & $0.217$ & $73.0$ & $6.937$ & $0.698$ & $1.447$ & $0.074$ & $1.395$ & $1.955$ \\
FW-D & $0.183$ & $73.4$ & $7.720$ & $0.783$ & $2.480$ & $0.027$ & $\mathbf{1.386}$ & $2.734$ \\
Greedy-E & $0.217$ & $72.7$ & $7.116$ & $\mathbf{0.260}$ & $1.476$ & $0.053$ & $1.409$ & $2.002$ \\
A-optimal & $0.215$ & $73.1$ & $\mathbf{6.922}$ & $0.656$ & $\mathbf{1.242}$ & $0.063$ & $1.392$ & $\mathbf{1.782}$ \\
D-optimal & $0.174$ & $\mathbf{73.4}$ & $7.921$ & $0.405$ & $2.157$ & $\mathbf{8.82\!\times\!10^{-3}}$ & $1.386$ & $2.499$ \\
Random (mean) & $0.091$ & $69.7$ & $14.9$ & $1.511$ & $3.417$ & $0.050$ & $1.392$ & $3.736$ \\
Coverage & $0.078$ & $69.2$ & $16.8$ & $0.551$ & $3.232$ & $0.062$ & $1.405$ & $3.528$ \\
\bottomrule
\end{tabular}}
\end{table*}

\begin{figure*}[t]
\centering
\includegraphics[width=\textwidth]{tab5_aprilgrid_sigma10}
\caption{Method comparison: AprilGrid, $\sigma{=}1.0$\,px. Bars show each metric (winner outlined in black). $\lambda_{\min}$/$\log\det$: higher is better; all others: lower is better.}
\label{fig:tab5_aprilgrid_sigma10}
\end{figure*}
\FloatBarrier

\begin{table*}[t]
\centering
\setlength{\tabcolsep}{4pt}
\caption{Matched ChArUco board ($13{\times}13$, 144\,pts, 46\,cm), $\sigma{=}1.0$~px noise. $N{=}720$, $K{=}20$.
Focal/PP errors in pixels, distortion error unitless.}
\label{tab:charuco_sigma10}
\scriptsize
\resizebox{\textwidth}{!}{\begin{tabular}{lrrrrrrrr}
\toprule
Method & $\lambda_{\min}$ & $\log\det$ & $\mathrm{tr}(\Sigma)$ & $\epsilon_f$\,(px) & $\epsilon_{pp}$\,(px) & $\epsilon_d$ & Train-RMS & Held-out \\
\midrule
\textbf{FW-E} & $0.196$ & $72.6$ & $7.853$ & $\mathbf{0.590}$ & $2.206$ & $0.019$ & $1.389$ & $2.569$ \\
FW-A & $0.196$ & $72.6$ & $7.853$ & $\mathbf{0.590}$ & $2.206$ & $0.019$ & $1.389$ & $2.569$ \\
FW-D & $0.189$ & $73.4$ & $7.920$ & $1.719$ & $2.295$ & $\mathbf{6.47\!\times\!10^{-3}}$ & $1.398$ & $2.699$ \\
Greedy-E & $\mathbf{0.204}$ & $72.8$ & $7.529$ & $0.944$ & $3.243$ & $0.033$ & $1.382$ & $3.668$ \\
A-optimal & $0.204$ & $72.9$ & $\mathbf{7.390}$ & $0.725$ & $\mathbf{2.088}$ & $0.010$ & $\mathbf{1.381}$ & $\mathbf{2.483}$ \\
D-optimal & $0.187$ & $\mathbf{73.4}$ & $7.840$ & $1.731$ & $2.309$ & $0.019$ & $1.402$ & $2.705$ \\
Random (mean) & $0.091$ & $69.7$ & $15.1$ & $1.415$ & $2.607$ & $0.034$ & $1.391$ & $2.962$ \\
Coverage & $0.074$ & $68.5$ & $17.9$ & $2.141$ & $4.913$ & $0.072$ & $1.385$ & $5.256$ \\
\bottomrule
\end{tabular}}
\end{table*}

\begin{figure*}[t]
\centering
\includegraphics[width=\textwidth]{tab6_charuco_sigma10}
\caption{Method comparison: ChArUco, $\sigma{=}1.0$\,px. Bars show each metric (winner outlined in black). $\lambda_{\min}$/$\log\det$: higher is better; all others: lower is better.}
\label{fig:tab6_charuco_sigma10}
\end{figure*}
\FloatBarrier

\begin{table}[ht]
\centering
\caption{Kalibr backend: held-out RMS\,(px) vs.\ budget $K$.
$N{=}720$ Isaac Sim AprilGrid, noiseless, 20 held-out views (FPS on translations).
Bold: best per column.}
\label{tab:kalibr_sweep}
\scriptsize
\setlength{\tabcolsep}{3.5pt}
\begin{tabular}{lrrrr}
\toprule
Method & $K{=}10$ & $K{=}20$ & $K{=}50$ & $K{=}80$ \\
\midrule
A-optimal  & $0.162$ & \textbf{0.090} & \textbf{0.099} & \textbf{0.100} \\
FW-D       & \textbf{0.111} & $0.214$ & $0.182$ & $0.190$ \\
Coverage   & $0.148$ & $0.190$ & $0.182$ & $0.170$ \\
FW-A       & $0.276$ & $0.276$ & $0.312$ & $0.243$ \\
D-optimal  & $0.310$ & $0.345$ & $0.184$ & $0.190$ \\
Random     & $0.692$ & $0.209$ & $0.192$ & $0.209$ \\
Greedy-E   & $0.610$ & $0.732$ & $0.631$ & $0.586$ \\
FW-E       & $0.772$ & $0.761$ & $0.205$ & $0.213$ \\
\bottomrule
\end{tabular}
\end{table}

\renewcommand{\thesubsection}{6.0}
\subsection{Ablation Studies}
\label{subsec:ablations}

\begin{figure*}[t]
\centering
\includegraphics[width=\textwidth]{fig_isaac_diversity_8methods}
\caption{Isaac Sim view diversity: 10 representative frames per method ($K{=}50$, $\sigma{=}0.5$\,px).
Greedy-E and FW-E cluster near-identical frontal close-up views that maximise $\lambda_{\min}$
but sacrifice geometric spread.
FW-A, FW-D, A-optimal, and D-optimal select views covering a wide range of distances,
tilt angles, and board positions, providing the diversity needed for well-conditioned
joint bundle-adjustment solvers such as Kalibr.}
\label{fig:isaac_diversity}
\end{figure*}

\paragraph{Noiseless baseline.}
Near-noiseless ($\sigma{\lesssim}0.01$\,px), all FIM methods yield
well-conditioned subsets
(Tables~\ref{tab:aprilgrid_noiseless} and~\ref{tab:charuco_noiseless}; Figures~\ref{fig:tab1_aprilgrid_noiseless} and~\ref{fig:tab2_charuco_noiseless}).
On AprilGrid, FW-E achieves the highest $\lambda_{\min}$ ($2.18\!\times\!10^3$)
and A-optimal achieves the best held-out RMS ($1.52\!\times\!10^{-2}$~px).
On matched ChArUco, Greedy-E leads $\lambda_{\min}$ ($2.04\!\times\!10^3$)
with the best held-out RMS ($2.25\!\times\!10^{-2}$~px).
Differences between FIM methods at noiseless conditions are within
solver-convergence noise.

\paragraph{E-optimal vs.\ A-optimal vs.\ D-optimal.}
No single criterion dominates across all regimes
(Tables~\ref{tab:aprilgrid_noiseless}, \ref{tab:charuco_noiseless},
\ref{tab:aprilgrid_sigma05}, and~\ref{tab:charuco_sigma10}):
Figure~\ref{fig:isaac_diversity} further shows that this criterion dependence is geometric as well as numerical: Greedy-E and FW-E often cluster similar frontal close-up views, whereas A-optimal and D-optimal select more diverse distances, tilts, and board positions.
\begin{itemize}
    \item \textbf{Near-noiseless:} FW-E wins $\lambda_{\min}$ on AprilGrid
    ($2.18\!\times\!10^3$); Greedy-E wins $\lambda_{\min}$ on ChArUco
    ($2.04\!\times\!10^3$). A-optimal achieves the best
    held-out RMS on AprilGrid ($1.52\!\times\!10^{-2}$~px). All methods
    are well-conditioned; differences are within solver-convergence noise.
    \item \textbf{$\sigma{=}0.5$~px:} FW-E wins $\lambda_{\min}$ on AprilGrid
    ($0.874$); Greedy-E and A-optimal are tied on matched ChArUco ($0.809$).
    Greedy-E wins $\epsilon_f$, $\epsilon_{pp}$, and held-out RMS on AprilGrid
    ($0.150$~px, $0.530$~px, $0.830$~px); A-optimal wins $\epsilon_f$ on ChArUco
    ($0.191$~px).
    \item \textbf{$\sigma{=}1.0$~px:} A-optimal wins held-out RMS on
    AprilGrid ($1.782$~px) and ChArUco ($2.483$~px, excluding collapsed FW-E).
    Greedy-E wins $\epsilon_f$ on AprilGrid ($0.260$~px); A-optimal wins
    $\epsilon_f$ on ChArUco ($0.725$~px). FW-E collapses on ChArUco
    ($\lambda_{\min}{=}0.0184$); Greedy-E and A-optimal achieve
    $\lambda_{\min}{=}0.204$.
\end{itemize}

\paragraph{Target type.}
Both targets are matched in corner count, pitch, and physical extent
(Section~\ref{sec:experiments}), so cross-target comparisons isolate
target geometry (inner-corner vs.\ outer-corner) from board size effects.
AprilGrid consistently achieves lower held-out RMS than matched ChArUco with
the best deterministic method at each noise level: near-noiseless
($1.52\!\times\!10^{-2}$~px vs.\ $1.80\!\times\!10^{-2}$~px),
$\sigma{=}0.5$~px ($0.830$~px vs.\ $1.271$~px),
and $\sigma{=}1.0$~px ($1.782$~px vs.\ $2.483$~px).
The outer-corner layout of AprilGrid tags is more precisely localizable in the
projection model, yielding a better-conditioned FIM.
Conversely, ChArUco's inner-corner geometry supports more uniform image-plane
sampling; random selection exploits this more effectively than FIM-based
methods on held-out accuracy at both noise levels
(Tables~\ref{tab:charuco_sigma05}--\ref{tab:charuco_sigma10}).

\paragraph{Subset budget $K$.}
Tables~\ref{tab:ksweep_aprilgrid_noiseless}, \ref{tab:ksweep_aprilgrid_sigma05},
and~\ref{tab:ksweep_charuco_sigma05} report $\lambda_{\min}$ across
$K\in\{5,10,15,20,30,50,100,200\}$ on three representative configurations
($N{=}720$ throughout), and Figure~\ref{fig:ksweep_mineig} visualises the same
trend.
In several K-sweep settings, the FW-rounded A- or D-optimal solutions coincide
exactly with the corresponding greedy discrete solutions, so some appendix rows
are intentionally numerically identical rather than copy errors.

\noindent\textbf{Diminishing returns.}
The marginal gain in $\lambda_{\min}$ per additional pose decreases as $K$
grows.
On ChArUco $\sigma{=}0.5$, A-optimal delivers $\lambda_{\min}{=}0.190$ at
$K{=}5$ but only ${\approx}0.010$ per additional pose in the
$K{=}100\!\to\!200$ range---an $18\times$ drop in per-pose efficiency.
The practical ``knee'' lies around $K{=}20$--$30$: beyond 30 frames, each
additional view adds less than $2\%$ of the cumulative gain already achieved.

\noindent\textbf{FW-E rounding instability at small $K$.}
FW-E fails at small $K$ on all three configurations.
On AprilGrid $\sigma{=}0.5$, FW-E collapses at $K{=}5$ ($1.11\!\times\!10^{-3}$)
and $K{=}10$ ($0.073$) before recovering at $K{\ge}15$.
On ChArUco $\sigma{=}0.5$, FW-E collapses at $K{=}5$ ($1.77\!\times\!10^{-3}$)
and $K{=}15$ ($2.34\!\times\!10^{-3}$), but achieves the best $\lambda_{\min}$
at $K{=}10$ ($0.429$) before the $K{=}15$ re-collapse.
On near-noiseless AprilGrid, FW-E collapses at $K{=}5$ ($1.16\!\times\!10^{-3}$)
but recovers at $K{\ge}10$ and leads from $K{\ge}15$.
In all cases the mechanism is the same: at small $K$, each candidate receives
initial weight $K/N{\approx}0.007$--$0.021$, the continuous relaxation remains
near-flat, and top-$K$ rounding assembles a poor discrete set.

\noindent\textbf{A-optimal leads at small-to-medium $K$ under noise.}
On AprilGrid $\sigma{=}0.5$, A-optimal leads all methods at $K{=}5$--$15$
($0.222$, $0.437$, $0.652$) and is surpassed only by FW-E at $K{\ge}20$.
On ChArUco $\sigma{=}0.5$, Greedy-E leads at $K{=}5$ ($0.197$ vs.\
A-optimal $0.190$), and A-optimal leads at $K{=}15$ ($0.618$ vs.\ $0.613$).
By balancing all eigenvalues of $H_{\mathrm{cal}}$, A-optimal avoids the
early pose-clustering that limits Greedy-E and the near-flat-relaxation
failure that prevents FW-E from converging at small $K$.
On noiseless AprilGrid, A-optimal leads at $K{=}5$ ($451$ vs.\ Greedy-E $428$)
but is overtaken by Greedy-E at $K{=}10$ ($1.04\!\times\!10^{3}$
vs.\ $1.01\!\times\!10^{3}$) and by FW-E from $K{\ge}15$.

\noindent\textbf{FW-E and Greedy-E lead at large budgets.}
At $K{\ge}20$ on both noisy configurations, FW-E and Greedy-E converge to
nearly identical $\lambda_{\min}$ (e.g., $0.874$ vs.\ $0.842$ at $K{=}20$ on
AprilGrid $\sigma{=}0.5$, within $3.7\%$) and both substantially outperform
D-optimal ($0.691$) and random ($0.346$) at the same budget.
At large budgets, D-optimal remains well above random under noise ---
$\lambda_{\min}{=}3.10$ vs.\ random $2.00$ at $K{=}100$ on AprilGrid $\sigma{=}0.5$ ---
and does \emph{not} regress below the baseline at any tested $K$.

\begin{figure}[t]
\centering
\includegraphics[width=\linewidth]{fig_ksweep_mineig}
\caption{$\lambda_{\min}(H_\mathrm{cal})$ vs.\ frame budget $K$ for the three
configurations in Tables~\ref{tab:ksweep_aprilgrid_noiseless}, \ref{tab:ksweep_aprilgrid_sigma05}, and~\ref{tab:ksweep_charuco_sigma05} (log scale).
FW-E leads consistently from $K{\ge}20$ on AprilGrid; A-Opt leads at
$K{\le}15$ under noise. Degenerate FW-E points at small $K$ are visible as
low outliers.}
\label{fig:ksweep_mineig}
\end{figure}

\paragraph{Full calibration metrics vs.\ budget $K$.}
Tables~\ref{tab:ksweep_aprilgrid_noiseless}--\ref{tab:ksweep_charuco_sigma10} extend the $\lambda_{\min}$ sweep above with
$\lambda_{\min}$, held-out RMS, and intrinsic error trends across a wide budget range.
The per-method held-out set grows as $K$ decreases, providing a stringent
out-of-sample test without any fixed split.
Figures~\ref{fig:ksweep_heldout} and~\ref{fig:ksweep_calib_errors} summarise
the held-out RMS and intrinsic-parameter errors visually.

\begin{figure}[t]
\centering
\includegraphics[width=\linewidth]{fig_ksweep_heldout}
\caption{Held-out reprojection RMS (px) vs.\ frame budget $K$.
All methods converge by $K{=}20$--$30$; the sharpest drop occurs between
$K{=}5$ and $K{=}15$. A-Opt achieves the lowest held-out error at small
budgets under noise; Coverage is worst at every budget.}
\label{fig:ksweep_heldout}
\end{figure}

\begin{figure}[t]
\centering
\includegraphics[width=\linewidth]{fig_ksweep_calib_errors}
\caption{Focal-length error $\epsilon_f$ (top) and principal-point error $\epsilon_{pp}$
(bottom) vs.\ budget $K$ at $\sigma{=}0.5$\,px. Log scale. Errors decrease
rapidly from $K{=}5$ to $K{=}20$, then level off. $\epsilon_{pp}$ shows more
method variation and does not converge as cleanly as $\epsilon_f$.}
\label{fig:ksweep_calib_errors}
\end{figure}

\paragraph{FW-E convergence and runtime.}
FW-E runs the full 300-iteration cap under the AND stopping criterion
($|\Delta f|/|f| < 10^{-4}$ \emph{and} duality gap ${<}10^{-4}$),
ensuring sufficient exploration of the relaxation before rounding
(Table~\ref{tab:fw_convergence}).
The duality gap $\delta$ decreases as $O(1/k)$ and requires ${\gg}300$
iterations to reach the $10^{-4}$ threshold from typical starting values
$\delta{\approx}3$--$9$; the 300-step cap is therefore the effective terminator.
At $\sigma{=}1.0$~px on matched ChArUco, FW-E collapses to
$\lambda_{\min}{=}0.0184$ and held-out RMS~$=10.5$~px despite completing
300 iterations (Table~\ref{tab:fw_convergence}).
The near-flat FIM landscape at high noise prevents the continuous relaxation
from concentrating before rounding; the short runtime ($1.20$~s) reflects
early closure of the relative-change condition before concentration.
At $\sigma{=}1.0$~px on AprilGrid, FW-E recovers successfully
($\lambda_{\min}{=}0.220$, runtime~$1.62$~s), so the failure is specific to
ChArUco at the highest tested noise level.
Wall-clock time is $1.2$--$2.5$~s for FW-E at 300 iterations across all six
standard configurations, dominated by block-extraction pre-processing.

\paragraph{Distortion robustness.}
To test whether method rankings hold under lens distortion, we generated twelve
additional datasets using the same hemisphere geometry ($N{=}720$, $K{=}20$)
but with Brown--Conrady distortion coefficients applied analytically at two levels:
medium distortion ($k_1{=}{-}0.15$, $k_2{=}0.05$, $p_1{=}0.001$, $p_2{=}0.001$)
and high distortion ($k_1{=}{-}0.35$, $k_2{=}0.15$, $p_1{=}0.002$, $p_2{=}0.002$),
each for both targets at all three noise levels ($\sigma{\approx}0$, $0.5$, $1.0$\,px).
$\lambda_{\min}(H_{\mathrm{cal}})$ results appear in
Tables~\ref{tab:aprilgrid_medium_noiseless}--\ref{tab:charuco_high_sigma10},
while Figures~\ref{fig:distortion_line}, \ref{fig:distortion_bar},
and~\ref{fig:distortion_rank_heatmap} summarise the trend, per-method
comparisons, and rank stability visually.

\begin{figure*}[t]
\centering
\includegraphics[width=\textwidth]{fig_distortion_line}
\caption{Distortion robustness: $\lambda_{\min}$ and held-out RMS vs.\ distortion
level (standard / medium / high) at $\sigma{=}0.5$\,px, $K{=}20$.
Left: AprilGrid; right: ChArUco.
Rankings are stable across distortion levels; $\lambda_{\min}$ rises with distortion
as additional lens-model parameters improve FIM conditioning.}
\label{fig:distortion_line}
\end{figure*}

\begin{figure*}[t]
\centering
\includegraphics[width=\textwidth]{fig_distortion_bar}
\caption{Grouped bar chart of $\lambda_{\min}$ (top) and held-out RMS (bottom)
per method across the three distortion levels at $\sigma{=}0.5$\,px.
FW-E and Greedy-E bars grow with distortion; the FW-E / Greedy-E gap narrows.}
\label{fig:distortion_bar}
\end{figure*}

\begin{figure*}[t]
\centering
\includegraphics[width=\textwidth]{fig_distortion_rank_heatmap}
\caption{Rank stability heatmap across all 18 distortion configurations
(2 targets $\times$ 3 distortion levels $\times$ 3 noise levels).
Rows: methods; columns: configurations; colour: rank ($1{=}$best, green;
$9{=}$worst, red).
FW-E and Greedy-E hold ranks 1--2 on $\lambda_{\min}$ across all 18 configs.}
\label{fig:distortion_rank_heatmap}
\end{figure*}

\paragraph{Medium distortion leaves rankings intact.}
Under medium distortion, FW-E achieves the highest $\lambda_{\min}$ on both
targets at noiseless and $\sigma{=}0.5$\,px conditions
($2.41\!\times\!10^{3}$ and $0.965$ on AprilGrid;
$2.23\!\times\!10^{3}$ and $0.899$ on ChArUco), with Greedy-E a close second
at every level (Tables~\ref{tab:aprilgrid_medium_noiseless}--\ref{tab:charuco_medium_sigma10};
Figures~\ref{fig:tab9_aprilgrid_medium_noiseless}--\ref{fig:tab14_charuco_medium_sigma10}).
The ${\approx}2.5\times$ FW-E/Random gap in $\lambda_{\min}$ is preserved across
all noise levels.

\paragraph{High distortion: FW-E leads noiseless; Greedy-E leads under noise.}
Under high distortion at noiseless conditions, FW-E achieves the highest
$\lambda_{\min}$ on both targets
($3.53\!\times\!10^{3}$ vs.\ Greedy-E $3.48\!\times\!10^{3}$ on AprilGrid;
$3.34\!\times\!10^{3}$ vs.\ $3.31\!\times\!10^{3}$ on ChArUco,
Tables~\ref{tab:aprilgrid_high_noiseless}--\ref{tab:charuco_high_sigma10};
Figures~\ref{fig:tab15_aprilgrid_high_noiseless}--\ref{fig:tab20_charuco_high_sigma10}).
Under noise ($\sigma{=}0.5$\,px), Greedy-E leads on AprilGrid
($1.393$ vs.\ FW-E $0.348$), indicating that FW-E rounding becomes unreliable
under high distortion combined with moderate noise.
Greedy-E and A-optimal remain stable across all high-distortion configurations.

\paragraph{Session 1 real-world.}
Table~\ref{tab:real_world} reports results across all eight methods on
$N_{\rm cand}{=}56$ candidate frames from a Logitech webcam capture, with a
fixed 20-frame held-out test set.
FIM-based methods outperform geometric heuristics on held-out reprojection
error, confirming the simulation findings at $\sigma{\approx}1$\,px
.
Greedy-E achieves the best held-out RMS ($0.752$\,px), closely followed by
FW-E ($0.764$\,px); both achieve the highest $\lambda_{\min}$
($5.1\!\times\!10^{-3}$).
Coverage is the worst ($1.044$\,px), a $39\%$ gap over Greedy-E.
Random ($0.883$\,px) falls between these groups, consistent with the
ChArUco pattern observed in simulation.
D-optimal achieves the best $\log\det H$ ($35.8$) and second-best held-out
RMS ($0.792$\,px), slightly behind the E-optimal variants.
Train-RMS spans a narrow band ($1.13$--$1.43$\,px) that does not reflect
the wider variation in held-out accuracy, replicating the simulation
finding that reprojection error on the training set is an unreliable
quality indicator (Section~\ref{subsec:fixed_backend}).

\paragraph{Appendix full-results atlas.}
For completeness, each appendix table is paired with a matching bar-plot
summary figure: noiseless standard-noise results use
Figures~\ref{fig:tab1_aprilgrid_noiseless} and~\ref{fig:tab2_charuco_noiseless},
medium-distortion results use
Figures~\ref{fig:tab9_aprilgrid_medium_noiseless}--\ref{fig:tab14_charuco_medium_sigma10},
high-distortion results use
Figures~\ref{fig:tab15_aprilgrid_high_noiseless}--\ref{fig:tab20_charuco_high_sigma10},
and K-sweep summaries use
Figures~\ref{fig:tab21_ksweep_aprilgrid_noiseless}--\ref{fig:tab26_ksweep_charuco_sigma10}.


\begin{table*}[t]
\centering
\setlength{\tabcolsep}{4pt}
\caption{AprilGrid, noiseless ($\sigma{\approx}0$\,px). $N{=}720$, $K{=}20$.
Focal/PP errors in pixels, distortion error unitless.
Random: mean over 25 trials.}
\label{tab:aprilgrid_noiseless}
\scriptsize
\resizebox{\textwidth}{!}{\begin{tabular}{lrrrrrrrr}
\toprule
Method & $\lambda_{\min}$ & $\log\det$ & $\mathrm{tr}(\Sigma)$ & $\epsilon_f$\,(px) & $\epsilon_{pp}$\,(px) & $\epsilon_d$ & Train-RMS & Held-out \\
\midrule
\textbf{FW-E} & $\mathbf{2176}$ & $155$ & $7.25\!\times\!10^{-4}$ & $\mathbf{3.77\!\times\!10^{-3}}$ & $0.023$ & $\mathbf{2.62\!\times\!10^{-4}}$ & $0.014$ & $0.027$ \\
FW-A & $2146$ & $156$ & $\mathbf{7.00\!\times\!10^{-4}}$ & $7.26\!\times\!10^{-3}$ & $9.86\!\times\!10^{-3}$ & $6.35\!\times\!10^{-4}$ & $0.014$ & $0.016$ \\
FW-D & $1804$ & $156$ & $7.80\!\times\!10^{-4}$ & $7.89\!\times\!10^{-3}$ & $0.025$ & $2.70\!\times\!10^{-4}$ & $\mathbf{0.014}$ & $0.027$ \\
Greedy-E & $2093$ & $155$ & $7.29\!\times\!10^{-4}$ & $5.07\!\times\!10^{-3}$ & $0.015$ & $6.00\!\times\!10^{-4}$ & $0.014$ & $0.020$ \\
A-optimal & $2112$ & $156$ & $7.08\!\times\!10^{-4}$ & $7.11\!\times\!10^{-3}$ & $\mathbf{8.06\!\times\!10^{-3}}$ & $7.09\!\times\!10^{-4}$ & $0.014$ & $\mathbf{0.015}$ \\
D-optimal & $1721$ & $\mathbf{156}$ & $8.01\!\times\!10^{-4}$ & $0.011$ & $0.029$ & $3.58\!\times\!10^{-4}$ & $0.014$ & $0.031$ \\
Random (mean) & $892$ & $153$ & $1.53\!\times\!10^{-3}$ & $0.015$ & $0.034$ & $4.95\!\times\!10^{-4}$ & $0.014$ & $0.037$ \\
Coverage & $761$ & $152$ & $1.73\!\times\!10^{-3}$ & $5.67\!\times\!10^{-3}$ & $0.031$ & $5.71\!\times\!10^{-4}$ & $0.014$ & $0.034$ \\
\bottomrule
\end{tabular}}
\end{table*}

\begin{figure*}[t]
\centering
\includegraphics[width=\textwidth]{tab1_aprilgrid_noiseless}
\caption{Method comparison: AprilGrid, $\sigma{\approx}0$. Bars show each metric (winner outlined in black). $\lambda_{\min}$/$\log\det$: higher is better; all others: lower is better.}
\label{fig:tab1_aprilgrid_noiseless}
\end{figure*}
\FloatBarrier


\begin{table*}[t]
\centering
\setlength{\tabcolsep}{4pt}
\caption{Matched ChArUco board ($13{\times}13$, 144\,pts, 46\,cm), near-noiseless ($\sigma{=}0.01$\,px). $N{=}720$, $K{=}20$.
Focal/PP errors in pixels, distortion error unitless.}
\label{tab:charuco_noiseless}
\scriptsize
\resizebox{\textwidth}{!}{\begin{tabular}{lrrrrrrrr}
\toprule
Method & $\lambda_{\min}$ & $\log\det$ & $\mathrm{tr}(\Sigma)$ & $\epsilon_f$\,(px) & $\epsilon_{pp}$\,(px) & $\epsilon_d$ & Train-RMS & Held-out \\
\midrule
\textbf{FW-E} & $1939$ & $155$ & $7.93\!\times\!10^{-4}$ & $\mathbf{5.88\!\times\!10^{-3}}$ & $0.022$ & $1.63\!\times\!10^{-4}$ & $0.014$ & $0.025$ \\
FW-A & $1939$ & $155$ & $7.93\!\times\!10^{-4}$ & $\mathbf{5.88\!\times\!10^{-3}}$ & $0.022$ & $1.63\!\times\!10^{-4}$ & $0.014$ & $0.025$ \\
FW-D & $1873$ & $156$ & $7.99\!\times\!10^{-4}$ & $0.017$ & $0.023$ & $7.24\!\times\!10^{-5}$ & $0.014$ & $0.027$ \\
Greedy-E & $\mathbf{2037}$ & $156$ & $\mathbf{7.50\!\times\!10^{-4}}$ & $8.43\!\times\!10^{-3}$ & $\mathbf{0.018}$ & $\mathbf{2.51\!\times\!10^{-5}}$ & $0.014$ & $\mathbf{0.022}$ \\
A-optimal & $2008$ & $156$ & $7.51\!\times\!10^{-4}$ & $8.02\!\times\!10^{-3}$ & $0.029$ & $1.12\!\times\!10^{-4}$ & $0.014$ & $0.031$ \\
D-optimal & $1842$ & $\mathbf{156}$ & $8.04\!\times\!10^{-4}$ & $0.016$ & $0.027$ & $7.70\!\times\!10^{-5}$ & $0.014$ & $0.031$ \\
Random (mean) & $885$ & $153$ & $1.54\!\times\!10^{-3}$ & $0.014$ & $0.026$ & $3.42\!\times\!10^{-4}$ & $0.014$ & $0.030$ \\
Coverage & $715$ & $151$ & $1.84\!\times\!10^{-3}$ & $0.022$ & $0.049$ & $7.13\!\times\!10^{-4}$ & $\mathbf{0.014}$ & $0.053$ \\
\bottomrule
\end{tabular}}
\end{table*}

\begin{figure*}[t]
\centering
\includegraphics[width=\textwidth]{tab2_charuco_noiseless}
\caption{Method comparison: ChArUco, $\sigma{\approx}0$. Bars show each metric (winner outlined in black). $\lambda_{\min}$/$\log\det$: higher is better; all others: lower is better.}
\label{fig:tab2_charuco_noiseless}
\end{figure*}
\FloatBarrier


\begin{table*}[t]
\centering
\caption{Frank-Wolfe solver convergence summary ($N{=}720$, $K{=}20$).
Matched ChArUco uses $\sigma{=}0.01$\,px for the near-noiseless row.
Under the AND stopping criterion ($|\Delta f|/|f| < 10^{-4}$ and duality gap $<10^{-4}$),
all solvers complete the full 300-iteration budget.
\emph{Steps}: FW iterations until both $|\Delta f|/|f| < 10^{-4}$ and duality gap $<10^{-4}$, capped at 300.}
\label{tab:fw_convergence}
\scriptsize
\setlength{\tabcolsep}{3.5pt}
\begin{tabular}{lllrrrr}
\toprule
Method & Target & $\sigma$ & Steps & Time (s) & Rounded $\lambda_{\min}$ & Held-out RMS\\
\midrule
FW-E & ChArUco (matched) & 0.01 & 300 & 2.56 & $1.94\!\times\!10^{3}$  & $2.52\!\times\!10^{-2}$ \\
FW-E & ChArUco (matched) & 0.5  & 300 & 2.40 & $7.78\!\times\!10^{-1}$ & $1.271$ \\
FW-E & ChArUco (matched) & 1.0  & 300 & 2.54 & $1.96\!\times\!10^{-1}$ & $2.569$ \\
FW-E & AprilGrid         & 0    & 300 & 2.50 & $2.18\!\times\!10^{3}$  & $2.70\!\times\!10^{-2}$ \\
FW-E & AprilGrid         & 0.5  & 300 & 2.37 & $8.74\!\times\!10^{-1}$ & $0.904$ \\
FW-E & AprilGrid         & 1.0  & 300 & 2.43 & $2.20\!\times\!10^{-1}$ & $2.404$ \\
\midrule
FW-A & ChArUco (matched) & 0.01 & 300 & 1.63 & $1.94\!\times\!10^{3}$  & $2.52\!\times\!10^{-2}$ \\
FW-A & ChArUco (matched) & 0.5  & 300 & 2.46 & $7.78\!\times\!10^{-1}$ & $1.271$ \\
FW-A & ChArUco (matched) & 1.0  & 300 & 2.14 & $1.96\!\times\!10^{-1}$ & $2.569$ \\
FW-A & AprilGrid         & 0    & 300 & 1.34 & $2.15\!\times\!10^{3}$  & $1.65\!\times\!10^{-2}$ \\
FW-A & AprilGrid         & 0.5  & 300 & 2.20 & $8.61\!\times\!10^{-1}$ & $0.819$ \\
FW-A & AprilGrid         & 1.0  & 300 & 2.66 & $2.17\!\times\!10^{-1}$ & $1.955$ \\
\midrule
FW-D & ChArUco (matched) & 0.01 & 300 & 2.20 & $1.87\!\times\!10^{3}$  & $2.66\!\times\!10^{-2}$ \\
FW-D & ChArUco (matched) & 0.5  & 300 & 2.21 & $7.51\!\times\!10^{-1}$ & $1.342$ \\
FW-D & ChArUco (matched) & 1.0  & 300 & 2.19 & $1.89\!\times\!10^{-1}$ & $2.699$ \\
FW-D & AprilGrid         & 0    & 300 & 2.33 & $1.80\!\times\!10^{3}$  & $2.74\!\times\!10^{-2}$ \\
FW-D & AprilGrid         & 0.5  & 300 & 2.48 & $7.24\!\times\!10^{-1}$ & $1.365$ \\
FW-D & AprilGrid         & 1.0  & 300 & 2.25 & $1.83\!\times\!10^{-1}$ & $2.734$ \\
\bottomrule
\end{tabular}
\end{table*}

\begin{table*}[t]
\centering
\setlength{\tabcolsep}{4pt}
\caption{Medium distortion ($k_1{=}{-}0.15$, $k_2{=}0.05$, $p_1{=}p_2{=}0.001$), AprilGrid, noiseless ($\sigma{\approx}0$\,px). $N{=}720$, $K{=}20$. Focal/PP errors in pixels, distortion error unitless. Random: mean over 25 trials.}
\label{tab:aprilgrid_medium_noiseless}
\scriptsize
\resizebox{\textwidth}{!}{\begin{tabular}{lrrrrrrrr}
\toprule
Method & $\lambda_{\min}$ & $\log\det$ & $\mathrm{tr}(\Sigma)$ & $\epsilon_f$\,(px) & $\epsilon_{pp}$\,(px) & $\epsilon_d$ & Train-RMS & Held-out \\
\midrule
\textbf{FW-E} & $2406$ & $158$ & $6.57\!\times\!10^{-4}$ & $0.018$ & $0.031$ & $4.22\!\times\!10^{-5}$ & $0.014$ & $0.033$ \\
FW-A & $2363$ & $158$ & $6.56\!\times\!10^{-4}$ & $4.92\!\times\!10^{-3}$ & $0.021$ & $1.88\!\times\!10^{-4}$ & $0.014$ & $0.024$ \\
FW-D & $1770$ & $158$ & $7.83\!\times\!10^{-4}$ & $0.010$ & $0.010$ & $\mathbf{3.92\!\times\!10^{-5}}$ & $0.014$ & $0.016$ \\
Greedy-E & $\mathbf{2406}$ & $158$ & $6.59\!\times\!10^{-4}$ & $0.018$ & $0.030$ & $1.04\!\times\!10^{-4}$ & $0.014$ & $0.032$ \\
A-optimal & $2365$ & $158$ & $\mathbf{6.50\!\times\!10^{-4}}$ & $\mathbf{1.27\!\times\!10^{-3}}$ & $0.013$ & $2.12\!\times\!10^{-4}$ & $0.014$ & $0.018$ \\
D-optimal & $1782$ & $\mathbf{158}$ & $7.76\!\times\!10^{-4}$ & $0.011$ & $8.94\!\times\!10^{-3}$ & $1.43\!\times\!10^{-4}$ & $0.014$ & $\mathbf{0.015}$ \\
Random (mean) & $937$ & $155$ & $1.46\!\times\!10^{-3}$ & $0.015$ & $0.034$ & $3.32\!\times\!10^{-4}$ & $0.014$ & $0.036$ \\
Coverage & $777$ & $155$ & $1.68\!\times\!10^{-3}$ & $2.27\!\times\!10^{-3}$ & $0.011$ & $5.33\!\times\!10^{-4}$ & $\mathbf{0.014}$ & $0.017$ \\
\bottomrule
\end{tabular}}
\end{table*}

\begin{figure*}[t]
\centering
\includegraphics[width=\textwidth]{tab9_aprilgrid_medium_noiseless}
\caption{Method comparison: AprilGrid, med.\ dist., $\sigma{\approx}0$. Bars show each metric (winner outlined in black). $\lambda_{\min}$/$\log\det$: higher is better; all others: lower is better.}
\label{fig:tab9_aprilgrid_medium_noiseless}
\end{figure*}
\FloatBarrier


\begin{table*}[t]
\centering
\setlength{\tabcolsep}{4pt}
\caption{Medium distortion ($k_1{=}{-}0.15$, $k_2{=}0.05$, $p_1{=}p_2{=}0.001$), Matched ChArUco, near-noiseless ($\sigma{=}0.01$\,px). $N{=}720$, $K{=}20$. Focal/PP errors in pixels, distortion error unitless. Random: mean over 25 trials.}
\label{tab:charuco_medium_noiseless}
\scriptsize
\resizebox{\textwidth}{!}{\begin{tabular}{lrrrrrrrr}
\toprule
Method & $\lambda_{\min}$ & $\log\det$ & $\mathrm{tr}(\Sigma)$ & $\epsilon_f$\,(px) & $\epsilon_{pp}$\,(px) & $\epsilon_d$ & Train-RMS & Held-out \\
\midrule
\textbf{FW-E} & $\mathbf{2230}$ & $158$ & $\mathbf{6.90\!\times\!10^{-4}}$ & $7.58\!\times\!10^{-3}$ & $0.016$ & $1.48\!\times\!10^{-4}$ & $\mathbf{0.014}$ & $0.019$ \\
FW-A & $\mathbf{2230}$ & $158$ & $\mathbf{6.90\!\times\!10^{-4}}$ & $7.58\!\times\!10^{-3}$ & $0.016$ & $1.48\!\times\!10^{-4}$ & $\mathbf{0.014}$ & $0.019$ \\
FW-D & $1821$ & $159$ & $7.63\!\times\!10^{-4}$ & $5.85\!\times\!10^{-3}$ & $\mathbf{7.78\!\times\!10^{-3}}$ & $1.01\!\times\!10^{-4}$ & $0.014$ & $\mathbf{0.014}$ \\
Greedy-E & $2165$ & $158$ & $7.03\!\times\!10^{-4}$ & $8.30\!\times\!10^{-3}$ & $0.022$ & $2.45\!\times\!10^{-4}$ & $0.014$ & $0.024$ \\
A-optimal & $2134$ & $158$ & $6.98\!\times\!10^{-4}$ & $0.010$ & $0.024$ & $3.82\!\times\!10^{-4}$ & $0.014$ & $0.027$ \\
D-optimal & $1674$ & $\mathbf{159}$ & $8.14\!\times\!10^{-4}$ & $5.82\!\times\!10^{-3}$ & $0.013$ & $\mathbf{2.68\!\times\!10^{-5}}$ & $0.014$ & $0.018$ \\
Random (mean) & $947$ & $155$ & $1.44\!\times\!10^{-3}$ & $0.014$ & $0.027$ & $2.83\!\times\!10^{-4}$ & $0.014$ & $0.030$ \\
Coverage & $732$ & $154$ & $1.78\!\times\!10^{-3}$ & $\mathbf{5.47\!\times\!10^{-3}}$ & $0.029$ & $7.62\!\times\!10^{-5}$ & $0.014$ & $0.032$ \\
\bottomrule
\end{tabular}}
\end{table*}

\begin{figure*}[t]
\centering
\includegraphics[width=\textwidth]{tab10_charuco_medium_noiseless}
\caption{Method comparison: ChArUco, med.\ dist., $\sigma{\approx}0$. Bars show each metric (winner outlined in black). $\lambda_{\min}$/$\log\det$: higher is better; all others: lower is better.}
\label{fig:tab10_charuco_medium_noiseless}
\end{figure*}
\FloatBarrier


\begin{table*}[t]
\centering
\setlength{\tabcolsep}{4pt}
\caption{Medium distortion ($k_1{=}{-}0.15$, $k_2{=}0.05$, $p_1{=}p_2{=}0.001$), AprilGrid, $\sigma{=}0.5$\,px noise. $N{=}720$, $K{=}20$. Focal/PP errors in pixels, distortion error unitless. Random: mean over 25 trials.}
\label{tab:aprilgrid_medium_sigma05}
\scriptsize
\resizebox{\textwidth}{!}{\begin{tabular}{lrrrrrrrr}
\toprule
Method & $\lambda_{\min}$ & $\log\det$ & $\mathrm{tr}(\Sigma)$ & $\epsilon_f$\,(px) & $\epsilon_{pp}$\,(px) & $\epsilon_d$ & Train-RMS & Held-out \\
\midrule
\textbf{FW-E} & $\mathbf{0.965}$ & $87.2$ & $1.639$ & $0.897$ & $1.577$ & $2.43\!\times\!10^{-3}$ & $0.693$ & $1.645$ \\
FW-A & $0.948$ & $87.3$ & $\mathbf{1.622}$ & $\mathbf{0.060}$ & $0.642$ & $0.011$ & $0.696$ & $0.884$ \\
FW-D & $0.710$ & $88.0$ & $1.953$ & $0.520$ & $0.508$ & $2.01\!\times\!10^{-3}$ & $0.703$ & $0.802$ \\
Greedy-E & $0.926$ & $87.1$ & $1.690$ & $0.899$ & $1.496$ & $\mathbf{1.08\!\times\!10^{-3}}$ & $0.697$ & $1.571$ \\
A-optimal & $0.905$ & $87.4$ & $1.667$ & $0.071$ & $0.706$ & $7.05\!\times\!10^{-3}$ & $0.698$ & $0.927$ \\
D-optimal & $0.715$ & $\mathbf{88.1}$ & $1.935$ & $0.572$ & $0.449$ & $7.18\!\times\!10^{-3}$ & $0.704$ & $\mathbf{0.742}$ \\
Random (mean) & $0.377$ & $84.6$ & $3.624$ & $0.763$ & $1.702$ & $0.017$ & $0.695$ & $1.815$ \\
Coverage & $0.313$ & $84.2$ & $4.168$ & $0.116$ & $0.593$ & $0.027$ & $\mathbf{0.692}$ & $0.861$ \\
\bottomrule
\end{tabular}}
\end{table*}

\begin{figure*}[t]
\centering
\includegraphics[width=\textwidth]{tab11_aprilgrid_medium_sigma05}
\caption{Method comparison: AprilGrid, med.\ dist., $\sigma{=}0.5$\,px. Bars show each metric (winner outlined in black). $\lambda_{\min}$/$\log\det$: higher is better; all others: lower is better.}
\label{fig:tab11_aprilgrid_medium_sigma05}
\end{figure*}
\FloatBarrier


\begin{table*}[t]
\centering
\setlength{\tabcolsep}{4pt}
\caption{Medium distortion ($k_1{=}{-}0.15$, $k_2{=}0.05$, $p_1{=}p_2{=}0.001$), Matched ChArUco, $\sigma{=}0.5$\,px noise. $N{=}720$, $K{=}20$. Focal/PP errors in pixels, distortion error unitless. Random: mean over 25 trials.}
\label{tab:charuco_medium_sigma05}
\scriptsize
\resizebox{\textwidth}{!}{\begin{tabular}{lrrrrrrrr}
\toprule
Method & $\lambda_{\min}$ & $\log\det$ & $\mathrm{tr}(\Sigma)$ & $\epsilon_f$\,(px) & $\epsilon_{pp}$\,(px) & $\epsilon_d$ & Train-RMS & Held-out \\
\midrule
\textbf{FW-E} & $\mathbf{0.899}$ & $87.8$ & $1.728$ & $0.449$ & $\mathbf{0.431}$ & $3.07\!\times\!10^{-3}$ & $\mathbf{0.677}$ & $\mathbf{0.752}$ \\
FW-A & $0.878$ & $87.6$ & $\mathbf{1.723}$ & $0.251$ & $0.697$ & $2.47\!\times\!10^{-3}$ & $0.683$ & $0.924$ \\
FW-D & $0.594$ & $88.7$ & $2.229$ & $0.278$ & $0.472$ & $6.82\!\times\!10^{-3}$ & $0.688$ & $0.772$ \\
Greedy-E & $0.865$ & $87.5$ & $1.770$ & $0.300$ & $0.741$ & $0.015$ & $0.687$ & $0.940$ \\
A-optimal & $0.861$ & $87.7$ & $1.728$ & $\mathbf{0.147}$ & $0.648$ & $6.37\!\times\!10^{-3}$ & $0.686$ & $0.891$ \\
D-optimal & $0.681$ & $\mathbf{88.7}$ & $2.014$ & $0.283$ & $0.839$ & $9.92\!\times\!10^{-3}$ & $0.691$ & $1.005$ \\
Random (mean) & $0.381$ & $84.8$ & $3.583$ & $0.700$ & $1.362$ & $0.014$ & $0.695$ & $1.527$ \\
Coverage & $0.295$ & $83.3$ & $4.420$ & $0.279$ & $1.431$ & $3.76\!\times\!10^{-3}$ & $0.695$ & $1.597$ \\
\bottomrule
\end{tabular}}
\end{table*}

\begin{figure*}[t]
\centering
\includegraphics[width=\textwidth]{tab12_charuco_medium_sigma05}
\caption{Method comparison: ChArUco, med.\ dist., $\sigma{=}0.5$\,px. Bars show each metric (winner outlined in black). $\lambda_{\min}$/$\log\det$: higher is better; all others: lower is better.}
\label{fig:tab12_charuco_medium_sigma05}
\end{figure*}
\FloatBarrier


\begin{table*}[t]
\centering
\setlength{\tabcolsep}{4pt}
\caption{Medium distortion ($k_1{=}{-}0.15$, $k_2{=}0.05$, $p_1{=}p_2{=}0.001$), AprilGrid, $\sigma{=}1.0$\,px noise. $N{=}720$, $K{=}20$. Focal/PP errors in pixels, distortion error unitless. Random: mean over 25 trials.}
\label{tab:aprilgrid_medium_sigma10}
\scriptsize
\resizebox{\textwidth}{!}{\begin{tabular}{lrrrrrrrr}
\toprule
Method & $\lambda_{\min}$ & $\log\det$ & $\mathrm{tr}(\Sigma)$ & $\epsilon_f$\,(px) & $\epsilon_{pp}$\,(px) & $\epsilon_d$ & Train-RMS & Held-out \\
\midrule
\textbf{FW-E} & $\mathbf{0.241}$ & $74.7$ & $6.545$ & $1.641$ & $2.912$ & $0.015$ & $1.391$ & $3.068$ \\
FW-A & $0.239$ & $74.8$ & $\mathbf{6.449}$ & $\mathbf{0.114}$ & $1.285$ & $0.021$ & $1.393$ & $1.769$ \\
FW-D & $0.179$ & $75.6$ & $7.752$ & $1.033$ & $1.019$ & $4.22\!\times\!10^{-3}$ & $1.405$ & $1.606$ \\
Greedy-E & $0.240$ & $74.8$ & $6.510$ & $1.357$ & $1.862$ & $0.020$ & $1.394$ & $2.166$ \\
A-optimal & $0.228$ & $74.9$ & $6.628$ & $0.188$ & $1.255$ & $\mathbf{4.21\!\times\!10^{-3}}$ & $1.403$ & $1.742$ \\
D-optimal & $0.175$ & $\mathbf{75.6}$ & $7.847$ & $0.996$ & $1.267$ & $4.73\!\times\!10^{-3}$ & $1.404$ & $1.766$ \\
Random (mean) & $0.096$ & $72.2$ & $14.3$ & $1.526$ & $3.411$ & $0.033$ & $1.390$ & $3.636$ \\
Coverage & $0.080$ & $71.7$ & $16.4$ & $0.226$ & $1.237$ & $0.055$ & $\mathbf{1.385}$ & $1.758$ \\
\bottomrule
\end{tabular}}
\end{table*}

\begin{figure*}[t]
\centering
\includegraphics[width=\textwidth]{tab13_aprilgrid_medium_sigma10}
\caption{Method comparison: AprilGrid, med.\ dist., $\sigma{=}1.0$\,px. Bars show each metric (winner outlined in black). $\lambda_{\min}$/$\log\det$: higher is better; all others: lower is better.}
\label{fig:tab13_aprilgrid_medium_sigma10}
\end{figure*}
\FloatBarrier


\begin{table*}[t]
\centering
\setlength{\tabcolsep}{4pt}
\caption{Medium distortion ($k_1{=}{-}0.15$, $k_2{=}0.05$, $p_1{=}p_2{=}0.001$), Matched ChArUco, $\sigma{=}1.0$\,px noise. $N{=}720$, $K{=}20$. Focal/PP errors in pixels, distortion error unitless. Random: mean over 25 trials.}
\label{tab:charuco_medium_sigma10}
\scriptsize
\resizebox{\textwidth}{!}{\begin{tabular}{lrrrrrrrr}
\toprule
Method & $\lambda_{\min}$ & $\log\det$ & $\mathrm{tr}(\Sigma)$ & $\epsilon_f$\,(px) & $\epsilon_{pp}$\,(px) & $\epsilon_d$ & Train-RMS & Held-out \\
\midrule
\textbf{FW-E} & $0.221$ & $75.1$ & $7.069$ & $0.978$ & $1.210$ & $0.013$ & $1.365$ & $1.722$ \\
FW-A & $\mathbf{0.222}$ & $75.2$ & $\mathbf{6.867}$ & $0.465$ & $0.881$ & $6.33\!\times\!10^{-3}$ & $\mathbf{1.362}$ & $1.499$ \\
FW-D & $0.169$ & $76.2$ & $8.075$ & $0.699$ & $\mathbf{0.384}$ & $\mathbf{3.20\!\times\!10^{-3}}$ & $1.382$ & $\mathbf{1.316}$ \\
Greedy-E & $0.217$ & $75.0$ & $7.058$ & $0.766$ & $1.676$ & $0.033$ & $1.372$ & $2.038$ \\
A-optimal & $0.216$ & $75.3$ & $6.922$ & $\mathbf{0.285}$ & $1.035$ & $3.97\!\times\!10^{-3}$ & $1.376$ & $1.574$ \\
D-optimal & $0.172$ & $\mathbf{76.2}$ & $8.008$ & $0.398$ & $1.554$ & $4.58\!\times\!10^{-3}$ & $1.387$ & $1.928$ \\
Random (mean) & $0.097$ & $72.3$ & $14.1$ & $1.408$ & $2.740$ & $0.029$ & $1.391$ & $3.070$ \\
Coverage & $0.075$ & $70.8$ & $17.4$ & $0.570$ & $2.848$ & $7.43\!\times\!10^{-3}$ & $1.390$ & $3.182$ \\
\bottomrule
\end{tabular}}
\end{table*}

\begin{figure*}[t]
\centering
\includegraphics[width=\textwidth]{tab14_charuco_medium_sigma10}
\caption{Method comparison: ChArUco, med.\ dist., $\sigma{=}1.0$\,px. Bars show each metric (winner outlined in black). $\lambda_{\min}$/$\log\det$: higher is better; all others: lower is better.}
\label{fig:tab14_charuco_medium_sigma10}
\end{figure*}
\FloatBarrier


\begin{table*}[t]
\centering
\setlength{\tabcolsep}{4pt}
\caption{High distortion ($k_1{=}{-}0.35$, $k_2{=}0.15$, $p_1{=}p_2{=}0.002$), AprilGrid, noiseless ($\sigma{\approx}0$\,px). $N{=}720$, $K{=}20$. Focal/PP errors in pixels, distortion error unitless. Random: mean over 25 trials.}
\label{tab:aprilgrid_high_noiseless}
\scriptsize
\resizebox{\textwidth}{!}{\begin{tabular}{lrrrrrrrr}
\toprule
Method & $\lambda_{\min}$ & $\log\det$ & $\mathrm{tr}(\Sigma)$ & $\epsilon_f$\,(px) & $\epsilon_{pp}$\,(px) & $\epsilon_d$ & Train-RMS & Held-out \\
\midrule
\textbf{FW-E} & $\mathbf{3526}$ & $161$ & $5.21\!\times\!10^{-4}$ & $0.014$ & $0.019$ & $1.34\!\times\!10^{-4}$ & $\mathbf{0.014}$ & $0.021$ \\
FW-A & $3484$ & $162$ & $\mathbf{5.12\!\times\!10^{-4}}$ & $0.016$ & $0.027$ & $1.33\!\times\!10^{-4}$ & $0.014$ & $0.028$ \\
FW-D & $2489$ & $\mathbf{163}$ & $6.13\!\times\!10^{-4}$ & $3.61\!\times\!10^{-3}$ & $0.020$ & $1.69\!\times\!10^{-4}$ & $0.014$ & $0.022$ \\
Greedy-E & $3477$ & $162$ & $5.19\!\times\!10^{-4}$ & $6.18\!\times\!10^{-3}$ & $7.01\!\times\!10^{-3}$ & $8.83\!\times\!10^{-5}$ & $0.014$ & $0.014$ \\
A-optimal & $3403$ & $162$ & $5.15\!\times\!10^{-4}$ & $7.05\!\times\!10^{-3}$ & $0.019$ & $4.73\!\times\!10^{-5}$ & $0.014$ & $0.022$ \\
D-optimal & $2428$ & $163$ & $6.20\!\times\!10^{-4}$ & $\mathbf{3.27\!\times\!10^{-3}}$ & $0.019$ & $1.40\!\times\!10^{-4}$ & $0.014$ & $0.022$ \\
Random (mean) & $1352$ & $159$ & $1.07\!\times\!10^{-3}$ & $0.013$ & $0.022$ & $1.42\!\times\!10^{-4}$ & $0.014$ & $0.024$ \\
Coverage & $988$ & $157$ & $1.40\!\times\!10^{-3}$ & $4.93\!\times\!10^{-3}$ & $0.056$ & $1.18\!\times\!10^{-4}$ & $0.014$ & $0.056$ \\
\bottomrule
\end{tabular}}
\end{table*}

\begin{figure*}[t]
\centering
\includegraphics[width=\textwidth]{tab15_aprilgrid_high_noiseless}
\caption{Method comparison: AprilGrid, high dist., $\sigma{\approx}0$. Bars show each metric (winner outlined in black). $\lambda_{\min}$/$\log\det$: higher is better; all others: lower is better.}
\label{fig:tab15_aprilgrid_high_noiseless}
\end{figure*}
\FloatBarrier


\begin{table*}[t]
\centering
\setlength{\tabcolsep}{4pt}
\caption{High distortion ($k_1{=}{-}0.35$, $k_2{=}0.15$, $p_1{=}p_2{=}0.002$), Matched ChArUco, near-noiseless ($\sigma{=}0.01$\,px). $N{=}720$, $K{=}20$. Focal/PP errors in pixels, distortion error unitless. Random: mean over 25 trials.}
\label{tab:charuco_high_noiseless}
\scriptsize
\resizebox{\textwidth}{!}{\begin{tabular}{lrrrrrrrr}
\toprule
Method & $\lambda_{\min}$ & $\log\det$ & $\mathrm{tr}(\Sigma)$ & $\epsilon_f$\,(px) & $\epsilon_{pp}$\,(px) & $\epsilon_d$ & Train-RMS & Held-out \\
\midrule
\textbf{FW-E} & $\mathbf{3342}$ & $160$ & $5.38\!\times\!10^{-4}$ & $0.013$ & $0.024$ & $1.59\!\times\!10^{-4}$ & $0.014$ & $0.026$ \\
FW-A & $3205$ & $162$ & $5.30\!\times\!10^{-4}$ & $5.26\!\times\!10^{-3}$ & $7.08\!\times\!10^{-3}$ & $9.31\!\times\!10^{-5}$ & $0.014$ & $0.014$ \\
FW-D & $2838$ & $163$ & $5.50\!\times\!10^{-4}$ & $\mathbf{3.87\!\times\!10^{-3}}$ & $0.023$ & $8.69\!\times\!10^{-5}$ & $0.014$ & $0.025$ \\
Greedy-E & $3305$ & $161$ & $5.42\!\times\!10^{-4}$ & $0.013$ & $0.026$ & $1.06\!\times\!10^{-4}$ & $0.014$ & $0.027$ \\
A-optimal & $3229$ & $162$ & $\mathbf{5.17\!\times\!10^{-4}}$ & $5.46\!\times\!10^{-3}$ & $6.51\!\times\!10^{-3}$ & $\mathbf{1.48\!\times\!10^{-5}}$ & $0.014$ & $\mathbf{0.014}$ \\
D-optimal & $2624$ & $\mathbf{163}$ & $5.80\!\times\!10^{-4}$ & $4.52\!\times\!10^{-3}$ & $0.024$ & $9.54\!\times\!10^{-5}$ & $0.014$ & $0.027$ \\
Random (mean) & $1416$ & $159$ & $1.03\!\times\!10^{-3}$ & $0.012$ & $0.027$ & $1.94\!\times\!10^{-4}$ & $0.014$ & $0.029$ \\
Coverage & $948$ & $158$ & $1.46\!\times\!10^{-3}$ & $0.016$ & $0.055$ & $1.00\!\times\!10^{-4}$ & $0.014$ & $0.054$ \\
\bottomrule
\end{tabular}}
\end{table*}

\begin{figure*}[t]
\centering
\includegraphics[width=\textwidth]{tab16_charuco_high_noiseless}
\caption{Method comparison: ChArUco, high dist., $\sigma{\approx}0$. Bars show each metric (winner outlined in black). $\lambda_{\min}$/$\log\det$: higher is better; all others: lower is better.}
\label{fig:tab16_charuco_high_noiseless}
\end{figure*}
\FloatBarrier


\begin{table*}[t]
\centering
\setlength{\tabcolsep}{4pt}
\caption{High distortion ($k_1{=}{-}0.35$, $k_2{=}0.15$, $p_1{=}p_2{=}0.002$), AprilGrid, $\sigma{=}0.5$\,px noise. $N{=}720$, $K{=}20$. Focal/PP errors in pixels, distortion error unitless. Random: mean over 25 trials.}
\label{tab:aprilgrid_high_sigma05}
\scriptsize
\resizebox{\textwidth}{!}{\begin{tabular}{lrrrrrrrr}
\toprule
Method & $\lambda_{\min}$ & $\log\det$ & $\mathrm{tr}(\Sigma)$ & $\epsilon_f$\,(px) & $\epsilon_{pp}$\,(px) & $\epsilon_d$ & Train-RMS & Held-out \\
\midrule
\textbf{FW-E} & $\mathbf{1.413}$ & $91.0$ & $1.301$ & $0.687$ & $0.926$ & $6.76\!\times\!10^{-3}$ & $\mathbf{0.679}$ & $1.033$ \\
FW-A & $1.357$ & $91.7$ & $\mathbf{1.275}$ & $0.977$ & $1.649$ & $8.17\!\times\!10^{-3}$ & $0.692$ & $1.638$ \\
FW-D & $1.025$ & $92.7$ & $1.500$ & $0.343$ & $0.539$ & $0.013$ & $0.703$ & $0.805$ \\
Greedy-E & $1.393$ & $91.2$ & $1.296$ & $0.310$ & $0.342$ & $4.44\!\times\!10^{-3}$ & $0.682$ & $0.697$ \\
A-optimal & $1.363$ & $91.7$ & $1.285$ & $0.351$ & $0.954$ & $2.38\!\times\!10^{-3}$ & $0.688$ & $1.079$ \\
D-optimal & $1.008$ & $\mathbf{92.7}$ & $1.527$ & $0.328$ & $0.672$ & $0.013$ & $0.703$ & $0.886$ \\
Random (mean) & $0.543$ & $88.9$ & $2.669$ & $0.625$ & $1.080$ & $7.12\!\times\!10^{-3}$ & $0.697$ & $1.215$ \\
Coverage & $0.398$ & $86.9$ & $3.479$ & $0.241$ & $2.820$ & $5.54\!\times\!10^{-3}$ & $0.710$ & $2.824$ \\
\bottomrule
\end{tabular}}
\end{table*}

\begin{figure*}[t]
\centering
\includegraphics[width=\textwidth]{tab17_aprilgrid_high_sigma05}
\caption{Method comparison: AprilGrid, high dist., $\sigma{=}0.5$\,px. Bars show each metric (winner outlined in black). $\lambda_{\min}$/$\log\det$: higher is better; all others: lower is better.}
\label{fig:tab17_aprilgrid_high_sigma05}
\end{figure*}
\FloatBarrier


\begin{table*}[t]
\centering
\setlength{\tabcolsep}{4pt}
\caption{High distortion ($k_1{=}{-}0.35$, $k_2{=}0.15$, $p_1{=}p_2{=}0.002$), Matched ChArUco, $\sigma{=}0.5$\,px noise. $N{=}720$, $K{=}20$. Focal/PP errors in pixels, distortion error unitless. Random: mean over 25 trials.}
\label{tab:charuco_high_sigma05}
\scriptsize
\resizebox{\textwidth}{!}{\begin{tabular}{lrrrrrrrr}
\toprule
Method & $\lambda_{\min}$ & $\log\det$ & $\mathrm{tr}(\Sigma)$ & $\epsilon_f$\,(px) & $\epsilon_{pp}$\,(px) & $\epsilon_d$ & Train-RMS & Held-out \\
\midrule
\textbf{FW-E} & $\mathbf{1.339}$ & $89.9$ & $1.343$ & $0.671$ & $1.212$ & $7.82\!\times\!10^{-3}$ & $0.696$ & $1.284$ \\
FW-A & $1.285$ & $92.2$ & $\mathbf{1.269}$ & $0.267$ & $0.587$ & $5.57\!\times\!10^{-3}$ & $0.709$ & $0.822$ \\
FW-D & $1.137$ & $92.7$ & $1.372$ & $\mathbf{0.197}$ & $1.146$ & $4.22\!\times\!10^{-3}$ & $0.704$ & $1.263$ \\
Greedy-E & $1.324$ & $90.5$ & $1.354$ & $0.676$ & $1.172$ & $2.62\!\times\!10^{-3}$ & $0.688$ & $1.247$ \\
A-optimal & $1.282$ & $92.0$ & $1.285$ & $0.435$ & $0.721$ & $\mathbf{7.14\!\times\!10^{-4}}$ & $0.696$ & $0.881$ \\
D-optimal & $1.052$ & $\mathbf{92.8}$ & $1.449$ & $0.232$ & $1.229$ & $4.70\!\times\!10^{-3}$ & $0.706$ & $1.336$ \\
Random (mean) & $0.568$ & $89.0$ & $2.562$ & $0.620$ & $1.343$ & $9.69\!\times\!10^{-3}$ & $0.698$ & $1.431$ \\
Coverage & $0.382$ & $87.6$ & $3.621$ & $0.784$ & $2.750$ & $5.02\!\times\!10^{-3}$ & $0.697$ & $2.666$ \\
\bottomrule
\end{tabular}}
\end{table*}

\begin{figure*}[t]
\centering
\includegraphics[width=\textwidth]{tab18_charuco_high_sigma05}
\caption{Method comparison: ChArUco, high dist., $\sigma{=}0.5$\,px. Bars show each metric (winner outlined in black). $\lambda_{\min}$/$\log\det$: higher is better; all others: lower is better.}
\label{fig:tab18_charuco_high_sigma05}
\end{figure*}
\FloatBarrier


\begin{table*}[t]
\centering
\setlength{\tabcolsep}{4pt}
\caption{High distortion ($k_1{=}{-}0.35$, $k_2{=}0.15$, $p_1{=}p_2{=}0.002$), AprilGrid, $\sigma{=}1.0$\,px noise. $N{=}720$, $K{=}20$. Focal/PP errors in pixels, distortion error unitless. Random: mean over 25 trials.}
\label{tab:aprilgrid_high_sigma10}
\scriptsize
\resizebox{\textwidth}{!}{\begin{tabular}{lrrrrrrrr}
\toprule
Method & $\lambda_{\min}$ & $\log\det$ & $\mathrm{tr}(\Sigma)$ & $\epsilon_f$\,(px) & $\epsilon_{pp}$\,(px) & $\epsilon_d$ & Train-RMS & Held-out \\
\midrule
\textbf{FW-E} & $\mathbf{0.354}$ & $78.5$ & $5.192$ & $1.372$ & $1.843$ & $0.014$ & $\mathbf{1.358}$ & $2.059$ \\
FW-A & $0.342$ & $78.8$ & $\mathbf{5.098}$ & $1.952$ & $3.400$ & $8.92\!\times\!10^{-3}$ & $1.381$ & $3.362$ \\
FW-D & $0.257$ & $80.2$ & $6.004$ & $0.686$ & $1.075$ & $0.025$ & $1.407$ & $1.608$ \\
Greedy-E & $0.349$ & $78.6$ & $5.174$ & $0.619$ & $0.676$ & $8.83\!\times\!10^{-3}$ & $1.365$ & $1.391$ \\
A-optimal & $0.340$ & $79.0$ & $5.106$ & $1.547$ & $3.147$ & $5.75\!\times\!10^{-3}$ & $1.378$ & $3.163$ \\
D-optimal & $0.252$ & $\mathbf{80.2}$ & $6.110$ & $0.656$ & $1.333$ & $0.027$ & $1.405$ & $1.765$ \\
Random (mean) & $0.137$ & $76.4$ & $10.6$ & $1.250$ & $2.166$ & $0.014$ & $1.394$ & $2.436$ \\
Coverage & $0.101$ & $74.4$ & $13.7$ & $0.494$ & $5.647$ & $0.010$ & $1.419$ & $5.654$ \\
\bottomrule
\end{tabular}}
\end{table*}

\begin{figure*}[t]
\centering
\includegraphics[width=\textwidth]{tab19_aprilgrid_high_sigma10}
\caption{Method comparison: AprilGrid, high dist., $\sigma{=}1.0$\,px. Bars show each metric (winner outlined in black). $\lambda_{\min}$/$\log\det$: higher is better; all others: lower is better.}
\label{fig:tab19_aprilgrid_high_sigma10}
\end{figure*}
\FloatBarrier


\begin{table*}[t]
\centering
\setlength{\tabcolsep}{4pt}
\caption{High distortion ($k_1{=}{-}0.35$, $k_2{=}0.15$, $p_1{=}p_2{=}0.002$), Matched ChArUco, $\sigma{=}1.0$\,px noise. $N{=}720$, $K{=}20$. Focal/PP errors in pixels, distortion error unitless. Random: mean over 25 trials.}
\label{tab:charuco_high_sigma10}
\scriptsize
\resizebox{\textwidth}{!}{\begin{tabular}{lrrrrrrrr}
\toprule
Method & $\lambda_{\min}$ & $\log\det$ & $\mathrm{tr}(\Sigma)$ & $\epsilon_f$\,(px) & $\epsilon_{pp}$\,(px) & $\epsilon_d$ & Train-RMS & Held-out \\
\midrule
\textbf{FW-E} & $\mathbf{0.336}$ & $77.5$ & $5.357$ & $1.344$ & $2.432$ & $0.015$ & $1.393$ & $2.576$ \\
FW-A & $0.323$ & $79.7$ & $\mathbf{5.060}$ & $0.533$ & $1.174$ & $0.011$ & $1.417$ & $1.643$ \\
FW-D & $0.265$ & $80.3$ & $5.764$ & $0.455$ & $1.687$ & $\mathbf{3.54\!\times\!10^{-3}}$ & $1.410$ & $2.007$ \\
Greedy-E & $0.332$ & $78.0$ & $5.399$ & $1.145$ & $2.486$ & $3.74\!\times\!10^{-3}$ & $1.382$ & $2.635$ \\
A-optimal & $0.322$ & $79.5$ & $5.121$ & $1.018$ & $1.763$ & $4.88\!\times\!10^{-3}$ & $1.395$ & $1.980$ \\
D-optimal & $0.265$ & $\mathbf{80.3}$ & $5.766$ & $\mathbf{0.435}$ & $1.933$ & $4.26\!\times\!10^{-3}$ & $1.412$ & $2.205$ \\
Random (mean) & $0.144$ & $76.5$ & $10.2$ & $1.238$ & $2.679$ & $0.019$ & $1.396$ & $2.856$ \\
Coverage & $0.097$ & $75.1$ & $14.3$ & $1.561$ & $5.478$ & $0.010$ & $1.395$ & $5.312$ \\
\bottomrule
\end{tabular}}
\end{table*}

\begin{figure*}[t]
\centering
\includegraphics[width=\textwidth]{tab20_charuco_high_sigma10}
\caption{Method comparison: ChArUco, high dist., $\sigma{=}1.0$\,px. Bars show each metric (winner outlined in black). $\lambda_{\min}$/$\log\det$: higher is better; all others: lower is better.}
\label{fig:tab20_charuco_high_sigma10}
\end{figure*}
\FloatBarrier



\begin{table*}[ht]
\centering
\setlength{\tabcolsep}{4pt}
\caption{AprilGrid, noiseless ($\sigma{\approx}0$\,px). $N{=}720$, $K{=}20$. Focal/PP errors in pixels, distortion error unitless. Held-out: mean reprojection RMS on all $N{-}K$ unselected frames per method. Random: mean over 25 trials. Bold: best per column per $K$.}
\label{tab:ksweep_aprilgrid_noiseless}
\scriptsize
\renewcommand{\arraystretch}{0.80}
\begin{tabular}{lrrrrrr}
\toprule
Method & $\epsilon_f$\,(px) & $\epsilon_{pp}$\,(px) & $\epsilon_d$ & $\|\Delta\theta\|$ & Train-RMS & Held-out \\
\midrule
\multicolumn{7}{l}{\textit{$K{=}5$}} \\[1pt]
FW-E & $2.55\!\times\!10^{1}$ & $2.47\!\times\!10^{1}$ & $1.82\!\times\!10^{-4}$ & $35.54$ & $1.39\!\times\!10^{-2}$ & $2.37\!\times\!10^{1}$ \\
FW-A & $2.47\!\times\!10^{-2}$ & $2.19\!\times\!10^{-2}$ & $2.42\!\times\!10^{-4}$ & $3.30\!\times\!10^{-2}$ & $1.42\!\times\!10^{-2}$ & $2.59\!\times\!10^{-2}$ \\
FW-D & $7.06\!\times\!10^{-3}$ & $1.29\!\times\!10^{-2}$ & $1.04\!\times\!10^{-4}$ & $1.47\!\times\!10^{-2}$ & $1.40\!\times\!10^{-2}$ & $1.91\!\times\!10^{-2}$ \\
Greedy-E & $3.89\!\times\!10^{-2}$ & $6.88\!\times\!10^{-2}$ & $7.09\!\times\!10^{-4}$ & $7.90\!\times\!10^{-2}$ & $1.45\!\times\!10^{-2}$ & $6.89\!\times\!10^{-2}$ \\
\textbf{A-optimal} & $\mathbf{4.69\!\times\!10^{-3}}$ & $\mathbf{4.37\!\times\!10^{-3}}$ & $\mathbf{9.21\!\times\!10^{-5}}$ & $\mathbf{6.41\!\times\!10^{-3}}$ & $\mathbf{1.38\!\times\!10^{-2}}$ & $\mathbf{1.42\!\times\!10^{-2}}$ \\
D-optimal & $1.22\!\times\!10^{-2}$ & $3.16\!\times\!10^{-2}$ & $8.35\!\times\!10^{-4}$ & $3.39\!\times\!10^{-2}$ & $1.43\!\times\!10^{-2}$ & $3.40\!\times\!10^{-2}$ \\
Random & $1.14\!\times\!10^{0}$ & $1.81\!\times\!10^{0}$ & $6.71\!\times\!10^{-4}$ & $2.137$ & $1.39\!\times\!10^{-2}$ & $1.82\!\times\!10^{0}$ \\
Coverage & $2.60\!\times\!10^{0}$ & $7.61\!\times\!10^{0}$ & $1.07\!\times\!10^{-3}$ & $8.036$ & $1.41\!\times\!10^{-2}$ & $7.53\!\times\!10^{0}$ \\
\addlinespace[1pt]
\multicolumn{7}{l}{\textit{$K{=}10$}} \\[1pt]
FW-E & $1.81\!\times\!10^{-2}$ & $1.62\!\times\!10^{-2}$ & $4.59\!\times\!10^{-4}$ & $2.43\!\times\!10^{-2}$ & $1.40\!\times\!10^{-2}$ & $2.01\!\times\!10^{-2}$ \\
FW-A & $2.27\!\times\!10^{-2}$ & $2.27\!\times\!10^{-2}$ & $6.84\!\times\!10^{-4}$ & $3.21\!\times\!10^{-2}$ & $1.40\!\times\!10^{-2}$ & $2.40\!\times\!10^{-2}$ \\
\textbf{FW-D} & $\mathbf{6.34\!\times\!10^{-3}}$ & $1.58\!\times\!10^{-2}$ & $1.89\!\times\!10^{-4}$ & $1.70\!\times\!10^{-2}$ & $1.40\!\times\!10^{-2}$ & $2.03\!\times\!10^{-2}$ \\
Greedy-E & $1.73\!\times\!10^{-2}$ & $3.98\!\times\!10^{-2}$ & $8.20\!\times\!10^{-4}$ & $4.34\!\times\!10^{-2}$ & $1.42\!\times\!10^{-2}$ & $4.17\!\times\!10^{-2}$ \\
\textbf{A-optimal} & $6.55\!\times\!10^{-3}$ & $\mathbf{1.43\!\times\!10^{-2}}$ & $2.94\!\times\!10^{-4}$ & $\mathbf{1.57\!\times\!10^{-2}}$ & $1.40\!\times\!10^{-2}$ & $\mathbf{1.88\!\times\!10^{-2}}$ \\
\textbf{D-optimal} & $9.66\!\times\!10^{-3}$ & $1.57\!\times\!10^{-2}$ & $\mathbf{7.21\!\times\!10^{-5}}$ & $1.84\!\times\!10^{-2}$ & $1.40\!\times\!10^{-2}$ & $1.93\!\times\!10^{-2}$ \\
\textbf{Random} & $4.47\!\times\!10^{-2}$ & $8.83\!\times\!10^{-2}$ & $5.96\!\times\!10^{-4}$ & $9.90\!\times\!10^{-2}$ & $\mathbf{1.39\!\times\!10^{-2}}$ & $9.12\!\times\!10^{-2}$ \\
Coverage & $2.09\!\times\!10^{-2}$ & $7.23\!\times\!10^{-2}$ & $7.57\!\times\!10^{-4}$ & $7.53\!\times\!10^{-2}$ & $1.39\!\times\!10^{-2}$ & $7.32\!\times\!10^{-2}$ \\
\addlinespace[1pt]
\multicolumn{7}{l}{\textit{$K{=}15$}} \\[1pt]
FW-E & $4.56\!\times\!10^{-3}$ & $2.16\!\times\!10^{-2}$ & $3.09\!\times\!10^{-4}$ & $2.21\!\times\!10^{-2}$ & $1.40\!\times\!10^{-2}$ & $2.52\!\times\!10^{-2}$ \\
\textbf{FW-A} & $8.47\!\times\!10^{-3}$ & $\mathbf{8.80\!\times\!10^{-3}}$ & $7.44\!\times\!10^{-4}$ & $\mathbf{1.22\!\times\!10^{-2}}$ & $1.40\!\times\!10^{-2}$ & $\mathbf{1.46\!\times\!10^{-2}}$ \\
\textbf{FW-D} & $1.08\!\times\!10^{-2}$ & $1.97\!\times\!10^{-2}$ & $1.62\!\times\!10^{-4}$ & $2.25\!\times\!10^{-2}$ & $\mathbf{1.38\!\times\!10^{-2}}$ & $2.37\!\times\!10^{-2}$ \\
Greedy-E & $8.21\!\times\!10^{-3}$ & $1.68\!\times\!10^{-2}$ & $6.96\!\times\!10^{-4}$ & $1.87\!\times\!10^{-2}$ & $1.42\!\times\!10^{-2}$ & $2.06\!\times\!10^{-2}$ \\
A-optimal & $9.62\!\times\!10^{-3}$ & $1.80\!\times\!10^{-2}$ & $3.77\!\times\!10^{-4}$ & $2.04\!\times\!10^{-2}$ & $1.39\!\times\!10^{-2}$ & $2.07\!\times\!10^{-2}$ \\
\textbf{D-optimal} & $\mathbf{4.04\!\times\!10^{-3}}$ & $2.27\!\times\!10^{-2}$ & $\mathbf{1.29\!\times\!10^{-4}}$ & $2.31\!\times\!10^{-2}$ & $1.38\!\times\!10^{-2}$ & $2.68\!\times\!10^{-2}$ \\
Random & $2.07\!\times\!10^{-2}$ & $4.50\!\times\!10^{-2}$ & $4.72\!\times\!10^{-4}$ & $4.95\!\times\!10^{-2}$ & $1.40\!\times\!10^{-2}$ & $4.74\!\times\!10^{-2}$ \\
Coverage & $1.65\!\times\!10^{-2}$ & $5.95\!\times\!10^{-2}$ & $4.43\!\times\!10^{-4}$ & $6.18\!\times\!10^{-2}$ & $1.39\!\times\!10^{-2}$ & $6.17\!\times\!10^{-2}$ \\
\addlinespace[1pt]
\multicolumn{7}{l}{\textit{$K{=}20$}} \\[1pt]
\textbf{FW-E} & $\mathbf{3.77\!\times\!10^{-3}}$ & $2.26\!\times\!10^{-2}$ & $\mathbf{2.62\!\times\!10^{-4}}$ & $2.30\!\times\!10^{-2}$ & $1.39\!\times\!10^{-2}$ & $2.70\!\times\!10^{-2}$ \\
FW-A & $7.26\!\times\!10^{-3}$ & $9.86\!\times\!10^{-3}$ & $6.35\!\times\!10^{-4}$ & $1.23\!\times\!10^{-2}$ & $1.39\!\times\!10^{-2}$ & $1.64\!\times\!10^{-2}$ \\
\textbf{FW-D} & $7.89\!\times\!10^{-3}$ & $2.49\!\times\!10^{-2}$ & $2.70\!\times\!10^{-4}$ & $2.61\!\times\!10^{-2}$ & $\mathbf{1.39\!\times\!10^{-2}}$ & $2.74\!\times\!10^{-2}$ \\
Greedy-E & $5.07\!\times\!10^{-3}$ & $1.51\!\times\!10^{-2}$ & $6.00\!\times\!10^{-4}$ & $1.59\!\times\!10^{-2}$ & $1.41\!\times\!10^{-2}$ & $1.98\!\times\!10^{-2}$ \\
\textbf{A-optimal} & $7.11\!\times\!10^{-3}$ & $\mathbf{8.06\!\times\!10^{-3}}$ & $7.09\!\times\!10^{-4}$ & $\mathbf{1.08\!\times\!10^{-2}}$ & $1.39\!\times\!10^{-2}$ & $\mathbf{1.52\!\times\!10^{-2}}$ \\
D-optimal & $1.05\!\times\!10^{-2}$ & $2.87\!\times\!10^{-2}$ & $3.58\!\times\!10^{-4}$ & $3.05\!\times\!10^{-2}$ & $1.39\!\times\!10^{-2}$ & $3.10\!\times\!10^{-2}$ \\
Random & $1.50\!\times\!10^{-2}$ & $3.40\!\times\!10^{-2}$ & $4.95\!\times\!10^{-4}$ & $3.72\!\times\!10^{-2}$ & $1.39\!\times\!10^{-2}$ & $3.73\!\times\!10^{-2}$ \\
Coverage & $5.67\!\times\!10^{-3}$ & $3.13\!\times\!10^{-2}$ & $5.71\!\times\!10^{-4}$ & $3.18\!\times\!10^{-2}$ & $1.40\!\times\!10^{-2}$ & $3.42\!\times\!10^{-2}$ \\
\addlinespace[1pt]
\multicolumn{7}{l}{\textit{$K{=}30$}} \\[1pt]
FW-E & $3.58\!\times\!10^{-3}$ & $8.07\!\times\!10^{-3}$ & $2.94\!\times\!10^{-4}$ & $8.83\!\times\!10^{-3}$ & $1.39\!\times\!10^{-2}$ & $1.51\!\times\!10^{-2}$ \\
FW-A & $1.17\!\times\!10^{-2}$ & $1.94\!\times\!10^{-2}$ & $6.22\!\times\!10^{-4}$ & $2.27\!\times\!10^{-2}$ & $1.39\!\times\!10^{-2}$ & $2.44\!\times\!10^{-2}$ \\
\textbf{FW-D} & $5.16\!\times\!10^{-3}$ & $\mathbf{4.77\!\times\!10^{-3}}$ & $3.46\!\times\!10^{-4}$ & $7.04\!\times\!10^{-3}$ & $1.39\!\times\!10^{-2}$ & $\mathbf{1.31\!\times\!10^{-2}}$ \\
Greedy-E & $5.28\!\times\!10^{-3}$ & $9.73\!\times\!10^{-3}$ & $4.42\!\times\!10^{-4}$ & $1.11\!\times\!10^{-2}$ & $1.40\!\times\!10^{-2}$ & $1.60\!\times\!10^{-2}$ \\
\textbf{A-optimal} & $9.66\!\times\!10^{-3}$ & $1.72\!\times\!10^{-2}$ & $6.88\!\times\!10^{-4}$ & $1.97\!\times\!10^{-2}$ & $\mathbf{1.38\!\times\!10^{-2}}$ & $2.23\!\times\!10^{-2}$ \\
\textbf{D-optimal} & $\mathbf{2.78\!\times\!10^{-3}}$ & $5.04\!\times\!10^{-3}$ & $\mathbf{2.71\!\times\!10^{-4}}$ & $\mathbf{5.76\!\times\!10^{-3}}$ & $1.39\!\times\!10^{-2}$ & $1.32\!\times\!10^{-2}$ \\
Random & $1.38\!\times\!10^{-2}$ & $3.10\!\times\!10^{-2}$ & $3.08\!\times\!10^{-4}$ & $3.40\!\times\!10^{-2}$ & $1.39\!\times\!10^{-2}$ & $3.37\!\times\!10^{-2}$ \\
Coverage & $9.80\!\times\!10^{-3}$ & $3.53\!\times\!10^{-2}$ & $6.49\!\times\!10^{-4}$ & $3.67\!\times\!10^{-2}$ & $1.40\!\times\!10^{-2}$ & $3.69\!\times\!10^{-2}$ \\
\addlinespace[1pt]
\multicolumn{7}{l}{\textit{$K{=}40$}} \\[1pt]
FW-E & $1.05\!\times\!10^{-2}$ & $1.13\!\times\!10^{-2}$ & $4.22\!\times\!10^{-4}$ & $1.55\!\times\!10^{-2}$ & $1.39\!\times\!10^{-2}$ & $1.72\!\times\!10^{-2}$ \\
FW-A & $8.64\!\times\!10^{-3}$ & $1.92\!\times\!10^{-2}$ & $4.40\!\times\!10^{-4}$ & $2.11\!\times\!10^{-2}$ & $1.38\!\times\!10^{-2}$ & $2.41\!\times\!10^{-2}$ \\
\textbf{FW-D} & $\mathbf{2.04\!\times\!10^{-3}}$ & $5.01\!\times\!10^{-3}$ & $4.29\!\times\!10^{-4}$ & $5.43\!\times\!10^{-3}$ & $1.38\!\times\!10^{-2}$ & $1.35\!\times\!10^{-2}$ \\
Greedy-E & $7.99\!\times\!10^{-3}$ & $1.27\!\times\!10^{-2}$ & $4.99\!\times\!10^{-4}$ & $1.50\!\times\!10^{-2}$ & $1.39\!\times\!10^{-2}$ & $1.81\!\times\!10^{-2}$ \\
A-optimal & $4.90\!\times\!10^{-3}$ & $2.38\!\times\!10^{-2}$ & $5.37\!\times\!10^{-4}$ & $2.43\!\times\!10^{-2}$ & $1.38\!\times\!10^{-2}$ & $2.83\!\times\!10^{-2}$ \\
\textbf{D-optimal} & $2.27\!\times\!10^{-3}$ & $\mathbf{4.09\!\times\!10^{-3}}$ & $5.01\!\times\!10^{-4}$ & $\mathbf{4.71\!\times\!10^{-3}}$ & $\mathbf{1.38\!\times\!10^{-2}}$ & $\mathbf{1.32\!\times\!10^{-2}}$ \\
\textbf{Random} & $1.07\!\times\!10^{-2}$ & $2.21\!\times\!10^{-2}$ & $\mathbf{3.45\!\times\!10^{-4}}$ & $2.46\!\times\!10^{-2}$ & $1.40\!\times\!10^{-2}$ & $2.64\!\times\!10^{-2}$ \\
Coverage & $1.32\!\times\!10^{-2}$ & $3.08\!\times\!10^{-2}$ & $6.62\!\times\!10^{-4}$ & $3.35\!\times\!10^{-2}$ & $1.40\!\times\!10^{-2}$ & $3.24\!\times\!10^{-2}$ \\
\addlinespace[1pt]
\multicolumn{7}{l}{\textit{$K{=}50$}} \\[1pt]
\textbf{FW-E} & $6.49\!\times\!10^{-3}$ & $1.36\!\times\!10^{-2}$ & $\mathbf{1.80\!\times\!10^{-4}}$ & $1.51\!\times\!10^{-2}$ & $1.39\!\times\!10^{-2}$ & $1.86\!\times\!10^{-2}$ \\
\textbf{FW-A} & $7.71\!\times\!10^{-3}$ & $1.31\!\times\!10^{-2}$ & $3.51\!\times\!10^{-4}$ & $1.52\!\times\!10^{-2}$ & $\mathbf{1.38\!\times\!10^{-2}}$ & $1.84\!\times\!10^{-2}$ \\
\textbf{FW-D} & $\mathbf{3.53\!\times\!10^{-3}}$ & $\mathbf{6.94\!\times\!10^{-3}}$ & $4.51\!\times\!10^{-4}$ & $\mathbf{7.80\!\times\!10^{-3}}$ & $1.38\!\times\!10^{-2}$ & $\mathbf{1.43\!\times\!10^{-2}}$ \\
Greedy-E & $7.91\!\times\!10^{-3}$ & $1.23\!\times\!10^{-2}$ & $2.26\!\times\!10^{-4}$ & $1.47\!\times\!10^{-2}$ & $1.39\!\times\!10^{-2}$ & $1.78\!\times\!10^{-2}$ \\
A-optimal & $7.16\!\times\!10^{-3}$ & $1.97\!\times\!10^{-2}$ & $6.91\!\times\!10^{-4}$ & $2.10\!\times\!10^{-2}$ & $1.39\!\times\!10^{-2}$ & $2.38\!\times\!10^{-2}$ \\
\textbf{D-optimal} & $\mathbf{3.53\!\times\!10^{-3}}$ & $\mathbf{6.94\!\times\!10^{-3}}$ & $4.51\!\times\!10^{-4}$ & $\mathbf{7.80\!\times\!10^{-3}}$ & $1.38\!\times\!10^{-2}$ & $\mathbf{1.43\!\times\!10^{-2}}$ \\
Random & $8.92\!\times\!10^{-3}$ & $2.07\!\times\!10^{-2}$ & $2.37\!\times\!10^{-4}$ & $2.25\!\times\!10^{-2}$ & $1.39\!\times\!10^{-2}$ & $2.48\!\times\!10^{-2}$ \\
Coverage & $1.66\!\times\!10^{-2}$ & $3.56\!\times\!10^{-2}$ & $7.09\!\times\!10^{-4}$ & $3.93\!\times\!10^{-2}$ & $1.39\!\times\!10^{-2}$ & $3.67\!\times\!10^{-2}$ \\
\bottomrule
\end{tabular}
\end{table*}

\begin{figure*}[t]
\centering
\includegraphics[width=\textwidth]{tab21_ksweep_aprilgrid_noiseless}
\caption{Method comparison: K-sweep: AprilGrid, $\sigma{\approx}0$. Bars show each metric (winner outlined in black). $\lambda_{\min}$/$\log\det$: higher is better; all others: lower is better.}
\label{fig:tab21_ksweep_aprilgrid_noiseless}
\end{figure*}


\begin{table*}[ht]
\centering
\setlength{\tabcolsep}{4pt}
\caption{ChArUco, near-noiseless ($\sigma{=}0.01$\,px). $N{=}720$, $K{=}20$. Focal/PP errors in pixels, distortion error unitless. Held-out: mean reprojection RMS on all $N{-}K$ unselected frames per method. Random: mean over 25 trials. Bold: best per column per $K$.}
\label{tab:ksweep_charuco_noiseless}
\scriptsize
\renewcommand{\arraystretch}{0.80}
\begin{tabular}{lrrrrrr}
\toprule
Method & $\epsilon_f$\,(px) & $\epsilon_{pp}$\,(px) & $\epsilon_d$ & $\|\Delta\theta\|$ & Train-RMS & Held-out \\
\midrule
\multicolumn{7}{l}{\textit{$K{=}5$}} \\[1pt]
FW-E & $1.91\!\times\!10^{1}$ & $1.85\!\times\!10^{1}$ & $7.95\!\times\!10^{-4}$ & $26.58$ & $1.37\!\times\!10^{-2}$ & $1.77\!\times\!10^{1}$ \\
\textbf{FW-A} & $2.44\!\times\!10^{1}$ & $5.04\!\times\!10^{1}$ & $\mathbf{2.42\!\times\!10^{-4}}$ & $56.02$ & $1.37\!\times\!10^{-2}$ & $5.11\!\times\!10^{1}$ \\
FW-D & $2.43\!\times\!10^{-2}$ & $6.85\!\times\!10^{-2}$ & $3.94\!\times\!10^{-4}$ & $7.27\!\times\!10^{-2}$ & $1.42\!\times\!10^{-2}$ & $6.89\!\times\!10^{-2}$ \\
\textbf{Greedy-E} & $3.27\!\times\!10^{-2}$ & $\mathbf{4.26\!\times\!10^{-2}}$ & $5.69\!\times\!10^{-4}$ & $\mathbf{5.37\!\times\!10^{-2}}$ & $1.44\!\times\!10^{-2}$ & $\mathbf{4.31\!\times\!10^{-2}}$ \\
\textbf{A-optimal} & $\mathbf{1.09\!\times\!10^{-2}}$ & $5.74\!\times\!10^{-2}$ & $4.71\!\times\!10^{-4}$ & $5.85\!\times\!10^{-2}$ & $1.41\!\times\!10^{-2}$ & $6.10\!\times\!10^{-2}$ \\
D-optimal & $3.11\!\times\!10^{-2}$ & $9.45\!\times\!10^{-2}$ & $9.17\!\times\!10^{-4}$ & $9.95\!\times\!10^{-2}$ & $1.42\!\times\!10^{-2}$ & $9.66\!\times\!10^{-2}$ \\
Random & $1.06\!\times\!10^{0}$ & $1.69\!\times\!10^{0}$ & $7.16\!\times\!10^{-4}$ & $1.995$ & $1.38\!\times\!10^{-2}$ & $1.70\!\times\!10^{0}$ \\
\textbf{Coverage} & $2.50\!\times\!10^{0}$ & $7.29\!\times\!10^{0}$ & $3.38\!\times\!10^{-4}$ & $7.706$ & $\mathbf{1.34\!\times\!10^{-2}}$ & $7.21\!\times\!10^{0}$ \\
\addlinespace[1pt]
\multicolumn{7}{l}{\textit{$K{=}10$}} \\[1pt]
\textbf{FW-E} & $8.98\!\times\!10^{-3}$ & $2.90\!\times\!10^{-2}$ & $\mathbf{1.88\!\times\!10^{-4}}$ & $3.03\!\times\!10^{-2}$ & $\mathbf{1.37\!\times\!10^{-2}}$ & $3.29\!\times\!10^{-2}$ \\
FW-A & $3.28\!\times\!10^{-2}$ & $6.96\!\times\!10^{-2}$ & $6.12\!\times\!10^{-4}$ & $7.69\!\times\!10^{-2}$ & $1.38\!\times\!10^{-2}$ & $7.23\!\times\!10^{-2}$ \\
\textbf{FW-D} & $\mathbf{3.60\!\times\!10^{-3}}$ & $3.12\!\times\!10^{-2}$ & $5.42\!\times\!10^{-4}$ & $3.15\!\times\!10^{-2}$ & $1.40\!\times\!10^{-2}$ & $3.55\!\times\!10^{-2}$ \\
Greedy-E & $1.19\!\times\!10^{-2}$ & $2.78\!\times\!10^{-2}$ & $4.29\!\times\!10^{-4}$ & $3.02\!\times\!10^{-2}$ & $1.40\!\times\!10^{-2}$ & $3.09\!\times\!10^{-2}$ \\
A-optimal & $1.18\!\times\!10^{-2}$ & $2.97\!\times\!10^{-2}$ & $5.02\!\times\!10^{-4}$ & $3.20\!\times\!10^{-2}$ & $1.39\!\times\!10^{-2}$ & $3.39\!\times\!10^{-2}$ \\
D-optimal & $8.05\!\times\!10^{-3}$ & $4.18\!\times\!10^{-2}$ & $7.01\!\times\!10^{-4}$ & $4.26\!\times\!10^{-2}$ & $1.39\!\times\!10^{-2}$ & $4.54\!\times\!10^{-2}$ \\
Random & $5.16\!\times\!10^{-2}$ & $9.03\!\times\!10^{-2}$ & $5.79\!\times\!10^{-4}$ & $0.1040$ & $1.39\!\times\!10^{-2}$ & $9.26\!\times\!10^{-2}$ \\
\textbf{Coverage} & $4.57\!\times\!10^{-3}$ & $\mathbf{2.63\!\times\!10^{-2}}$ & $9.42\!\times\!10^{-4}$ & $\mathbf{2.67\!\times\!10^{-2}}$ & $1.39\!\times\!10^{-2}$ & $\mathbf{2.84\!\times\!10^{-2}}$ \\
\addlinespace[1pt]
\multicolumn{7}{l}{\textit{$K{=}15$}} \\[1pt]
FW-E & $9.69\!\times\!10^{-3}$ & $3.33\!\times\!10^{-2}$ & $1.40\!\times\!10^{-4}$ & $3.47\!\times\!10^{-2}$ & $1.40\!\times\!10^{-2}$ & $3.59\!\times\!10^{-2}$ \\
\textbf{FW-A} & $\mathbf{8.55\!\times\!10^{-3}}$ & $\mathbf{2.35\!\times\!10^{-2}}$ & $2.22\!\times\!10^{-4}$ & $\mathbf{2.50\!\times\!10^{-2}}$ & $\mathbf{1.37\!\times\!10^{-2}}$ & $\mathbf{2.77\!\times\!10^{-2}}$ \\
FW-D & $1.60\!\times\!10^{-2}$ & $2.96\!\times\!10^{-2}$ & $9.93\!\times\!10^{-5}$ & $3.37\!\times\!10^{-2}$ & $1.40\!\times\!10^{-2}$ & $3.39\!\times\!10^{-2}$ \\
Greedy-E & $1.21\!\times\!10^{-2}$ & $2.70\!\times\!10^{-2}$ & $5.37\!\times\!10^{-5}$ & $2.97\!\times\!10^{-2}$ & $1.39\!\times\!10^{-2}$ & $3.12\!\times\!10^{-2}$ \\
\textbf{A-optimal} & $9.50\!\times\!10^{-3}$ & $3.51\!\times\!10^{-2}$ & $\mathbf{3.69\!\times\!10^{-5}}$ & $3.63\!\times\!10^{-2}$ & $1.39\!\times\!10^{-2}$ & $3.86\!\times\!10^{-2}$ \\
D-optimal & $1.56\!\times\!10^{-2}$ & $2.78\!\times\!10^{-2}$ & $1.28\!\times\!10^{-4}$ & $3.19\!\times\!10^{-2}$ & $1.40\!\times\!10^{-2}$ & $3.21\!\times\!10^{-2}$ \\
Random & $2.42\!\times\!10^{-2}$ & $4.50\!\times\!10^{-2}$ & $5.79\!\times\!10^{-4}$ & $5.11\!\times\!10^{-2}$ & $1.39\!\times\!10^{-2}$ & $4.73\!\times\!10^{-2}$ \\
Coverage & $3.09\!\times\!10^{-2}$ & $6.40\!\times\!10^{-2}$ & $8.58\!\times\!10^{-4}$ & $7.11\!\times\!10^{-2}$ & $1.38\!\times\!10^{-2}$ & $6.55\!\times\!10^{-2}$ \\
\addlinespace[1pt]
\multicolumn{7}{l}{\textit{$K{=}20$}} \\[1pt]
\textbf{FW-E} & $\mathbf{5.88\!\times\!10^{-3}}$ & $2.16\!\times\!10^{-2}$ & $1.63\!\times\!10^{-4}$ & $2.23\!\times\!10^{-2}$ & $1.39\!\times\!10^{-2}$ & $2.52\!\times\!10^{-2}$ \\
\textbf{FW-A} & $\mathbf{5.88\!\times\!10^{-3}}$ & $2.16\!\times\!10^{-2}$ & $1.63\!\times\!10^{-4}$ & $2.23\!\times\!10^{-2}$ & $1.39\!\times\!10^{-2}$ & $2.52\!\times\!10^{-2}$ \\
FW-D & $1.71\!\times\!10^{-2}$ & $2.25\!\times\!10^{-2}$ & $7.24\!\times\!10^{-5}$ & $2.83\!\times\!10^{-2}$ & $1.40\!\times\!10^{-2}$ & $2.66\!\times\!10^{-2}$ \\
\textbf{Greedy-E} & $8.43\!\times\!10^{-3}$ & $\mathbf{1.82\!\times\!10^{-2}}$ & $\mathbf{2.51\!\times\!10^{-5}}$ & $\mathbf{2.01\!\times\!10^{-2}}$ & $1.41\!\times\!10^{-2}$ & $\mathbf{2.25\!\times\!10^{-2}}$ \\
A-optimal & $8.02\!\times\!10^{-3}$ & $2.90\!\times\!10^{-2}$ & $1.12\!\times\!10^{-4}$ & $3.01\!\times\!10^{-2}$ & $1.39\!\times\!10^{-2}$ & $3.15\!\times\!10^{-2}$ \\
D-optimal & $1.61\!\times\!10^{-2}$ & $2.71\!\times\!10^{-2}$ & $7.70\!\times\!10^{-5}$ & $3.15\!\times\!10^{-2}$ & $1.40\!\times\!10^{-2}$ & $3.11\!\times\!10^{-2}$ \\
Random & $1.41\!\times\!10^{-2}$ & $2.60\!\times\!10^{-2}$ & $3.42\!\times\!10^{-4}$ & $2.96\!\times\!10^{-2}$ & $1.39\!\times\!10^{-2}$ & $2.95\!\times\!10^{-2}$ \\
\textbf{Coverage} & $2.17\!\times\!10^{-2}$ & $4.92\!\times\!10^{-2}$ & $7.13\!\times\!10^{-4}$ & $5.38\!\times\!10^{-2}$ & $\mathbf{1.38\!\times\!10^{-2}}$ & $5.25\!\times\!10^{-2}$ \\
\addlinespace[1pt]
\multicolumn{7}{l}{\textit{$K{=}30$}} \\[1pt]
FW-E & $4.46\!\times\!10^{-3}$ & $2.22\!\times\!10^{-2}$ & $8.80\!\times\!10^{-5}$ & $2.26\!\times\!10^{-2}$ & $1.40\!\times\!10^{-2}$ & $2.57\!\times\!10^{-2}$ \\
\textbf{FW-A} & $\mathbf{3.12\!\times\!10^{-3}}$ & $\mathbf{1.22\!\times\!10^{-2}}$ & $1.03\!\times\!10^{-4}$ & $\mathbf{1.26\!\times\!10^{-2}}$ & $\mathbf{1.39\!\times\!10^{-2}}$ & $\mathbf{1.76\!\times\!10^{-2}}$ \\
FW-D & $1.32\!\times\!10^{-2}$ & $1.71\!\times\!10^{-2}$ & $8.38\!\times\!10^{-5}$ & $2.16\!\times\!10^{-2}$ & $1.41\!\times\!10^{-2}$ & $2.18\!\times\!10^{-2}$ \\
Greedy-E & $7.99\!\times\!10^{-3}$ & $2.23\!\times\!10^{-2}$ & $8.73\!\times\!10^{-5}$ & $2.37\!\times\!10^{-2}$ & $1.41\!\times\!10^{-2}$ & $2.60\!\times\!10^{-2}$ \\
\textbf{A-optimal} & $4.00\!\times\!10^{-3}$ & $1.83\!\times\!10^{-2}$ & $\mathbf{1.84\!\times\!10^{-5}}$ & $1.87\!\times\!10^{-2}$ & $1.40\!\times\!10^{-2}$ & $2.22\!\times\!10^{-2}$ \\
D-optimal & $1.39\!\times\!10^{-2}$ & $1.68\!\times\!10^{-2}$ & $1.72\!\times\!10^{-4}$ & $2.18\!\times\!10^{-2}$ & $1.41\!\times\!10^{-2}$ & $2.11\!\times\!10^{-2}$ \\
Random & $1.22\!\times\!10^{-2}$ & $2.40\!\times\!10^{-2}$ & $3.27\!\times\!10^{-4}$ & $2.69\!\times\!10^{-2}$ & $1.39\!\times\!10^{-2}$ & $2.77\!\times\!10^{-2}$ \\
Coverage & $8.88\!\times\!10^{-3}$ & $2.95\!\times\!10^{-2}$ & $8.89\!\times\!10^{-4}$ & $3.09\!\times\!10^{-2}$ & $1.40\!\times\!10^{-2}$ & $3.43\!\times\!10^{-2}$ \\
\addlinespace[1pt]
\multicolumn{7}{l}{\textit{$K{=}40$}} \\[1pt]
FW-E & $8.57\!\times\!10^{-3}$ & $2.71\!\times\!10^{-2}$ & $1.10\!\times\!10^{-4}$ & $2.85\!\times\!10^{-2}$ & $1.40\!\times\!10^{-2}$ & $3.02\!\times\!10^{-2}$ \\
FW-A & $7.45\!\times\!10^{-3}$ & $1.91\!\times\!10^{-2}$ & $5.53\!\times\!10^{-5}$ & $\mathbf{2.05\!\times\!10^{-2}}$ & $1.40\!\times\!10^{-2}$ & $2.31\!\times\!10^{-2}$ \\
FW-D & $1.10\!\times\!10^{-2}$ & $1.86\!\times\!10^{-2}$ & $1.82\!\times\!10^{-4}$ & $2.16\!\times\!10^{-2}$ & $1.40\!\times\!10^{-2}$ & $2.31\!\times\!10^{-2}$ \\
\textbf{Greedy-E} & $\mathbf{4.47\!\times\!10^{-3}}$ & $2.30\!\times\!10^{-2}$ & $\mathbf{3.54\!\times\!10^{-5}}$ & $2.34\!\times\!10^{-2}$ & $1.41\!\times\!10^{-2}$ & $2.69\!\times\!10^{-2}$ \\
A-optimal & $8.00\!\times\!10^{-3}$ & $2.06\!\times\!10^{-2}$ & $1.59\!\times\!10^{-4}$ & $2.21\!\times\!10^{-2}$ & $1.40\!\times\!10^{-2}$ & $2.36\!\times\!10^{-2}$ \\
\textbf{D-optimal} & $1.10\!\times\!10^{-2}$ & $\mathbf{1.86\!\times\!10^{-2}}$ & $1.82\!\times\!10^{-4}$ & $2.16\!\times\!10^{-2}$ & $1.40\!\times\!10^{-2}$ & $\mathbf{2.31\!\times\!10^{-2}}$ \\
Random & $1.22\!\times\!10^{-2}$ & $2.35\!\times\!10^{-2}$ & $2.67\!\times\!10^{-4}$ & $2.65\!\times\!10^{-2}$ & $1.39\!\times\!10^{-2}$ & $2.72\!\times\!10^{-2}$ \\
\textbf{Coverage} & $1.33\!\times\!10^{-2}$ & $2.78\!\times\!10^{-2}$ & $1.12\!\times\!10^{-3}$ & $3.08\!\times\!10^{-2}$ & $\mathbf{1.39\!\times\!10^{-2}}$ & $3.14\!\times\!10^{-2}$ \\
\addlinespace[1pt]
\multicolumn{7}{l}{\textit{$K{=}50$}} \\[1pt]
FW-E & $7.07\!\times\!10^{-3}$ & $2.04\!\times\!10^{-2}$ & $7.97\!\times\!10^{-5}$ & $2.16\!\times\!10^{-2}$ & $1.40\!\times\!10^{-2}$ & $2.41\!\times\!10^{-2}$ \\
\textbf{FW-A} & $1.39\!\times\!10^{-2}$ & $\mathbf{1.54\!\times\!10^{-2}}$ & $1.77\!\times\!10^{-4}$ & $\mathbf{2.08\!\times\!10^{-2}}$ & $1.40\!\times\!10^{-2}$ & $\mathbf{2.04\!\times\!10^{-2}}$ \\
FW-D & $1.23\!\times\!10^{-2}$ & $1.69\!\times\!10^{-2}$ & $5.13\!\times\!10^{-5}$ & $2.09\!\times\!10^{-2}$ & $1.40\!\times\!10^{-2}$ & $2.14\!\times\!10^{-2}$ \\
\textbf{Greedy-E} & $\mathbf{6.52\!\times\!10^{-3}}$ & $2.18\!\times\!10^{-2}$ & $9.36\!\times\!10^{-5}$ & $2.27\!\times\!10^{-2}$ & $1.40\!\times\!10^{-2}$ & $2.55\!\times\!10^{-2}$ \\
A-optimal & $7.74\!\times\!10^{-3}$ & $2.24\!\times\!10^{-2}$ & $9.12\!\times\!10^{-5}$ & $2.37\!\times\!10^{-2}$ & $1.40\!\times\!10^{-2}$ & $2.51\!\times\!10^{-2}$ \\
\textbf{D-optimal} & $1.23\!\times\!10^{-2}$ & $1.69\!\times\!10^{-2}$ & $\mathbf{5.13\!\times\!10^{-5}}$ & $2.09\!\times\!10^{-2}$ & $1.40\!\times\!10^{-2}$ & $2.14\!\times\!10^{-2}$ \\
\textbf{Random} & $9.65\!\times\!10^{-3}$ & $1.98\!\times\!10^{-2}$ & $2.70\!\times\!10^{-4}$ & $2.20\!\times\!10^{-2}$ & $\mathbf{1.39\!\times\!10^{-2}}$ & $2.39\!\times\!10^{-2}$ \\
Coverage & $1.09\!\times\!10^{-2}$ & $2.16\!\times\!10^{-2}$ & $6.91\!\times\!10^{-4}$ & $2.42\!\times\!10^{-2}$ & $1.39\!\times\!10^{-2}$ & $2.49\!\times\!10^{-2}$ \\
\bottomrule
\end{tabular}
\end{table*}

\begin{figure*}[t]
\centering
\includegraphics[width=\textwidth]{tab22_ksweep_charuco_noiseless}
\caption{Method comparison: K-sweep: ChArUco, $\sigma{\approx}0$. Bars show each metric (winner outlined in black). $\lambda_{\min}$/$\log\det$: higher is better; all others: lower is better.}
\label{fig:tab22_ksweep_charuco_noiseless}
\end{figure*}


\begin{table*}[ht]
\centering
\setlength{\tabcolsep}{4pt}
\caption{AprilGrid, $\sigma{=}0.5$\,px. $N{=}720$, $K{=}20$. Focal/PP errors in pixels, distortion error unitless. Held-out: mean reprojection RMS on all $N{-}K$ unselected frames per method. Random: mean over 25 trials. Bold: best per column per $K$.}
\label{tab:ksweep_aprilgrid_sigma05}
\scriptsize
\renewcommand{\arraystretch}{0.80}
\begin{tabular}{lrrrrrr}
\toprule
Method & $\epsilon_f$\,(px) & $\epsilon_{pp}$\,(px) & $\epsilon_d$ & $\|\Delta\theta\|$ & Train-RMS & Held-out \\
\midrule
\multicolumn{7}{l}{\textit{$K{=}5$}} \\[1pt]
FW-E & $5.070$ & $4.310$ & $6.59\!\times\!10^{-3}$ & $6.654$ & $0.694$ & $4.119$ \\
FW-A & $0.640$ & $1.664$ & $0.0199$ & $1.783$ & $0.707$ & $1.744$ \\
\textbf{FW-D} & $0.504$ & $2.099$ & $0.0268$ & $2.159$ & $\mathbf{0.686}$ & $2.288$ \\
Greedy-E & $1.087$ & $2.059$ & $0.0775$ & $2.329$ & $0.694$ & $2.379$ \\
A-optimal & $1.028$ & $1.792$ & $4.81\!\times\!10^{-3}$ & $2.066$ & $0.695$ & $2.009$ \\
\textbf{D-optimal} & $\mathbf{0.352}$ & $\mathbf{0.628}$ & $\mathbf{4.61\!\times\!10^{-3}}$ & $\mathbf{0.7194}$ & $0.699$ & $\mathbf{0.943}$ \\
Random & $5.457$ & $8.338$ & $0.0343$ & $9.965$ & $0.696$ & $8.339$ \\
Coverage & $4.499$ & $10.267$ & $0.0587$ & $11.21$ & $0.706$ & $10.082$ \\
\addlinespace[1pt]
\multicolumn{7}{l}{\textit{$K{=}10$}} \\[1pt]
\textbf{FW-E} & $0.908$ & $0.795$ & $0.0229$ & $1.207$ & $0.699$ & $\mathbf{0.985}$ \\
\textbf{FW-A} & $\mathbf{0.190}$ & $1.276$ & $0.0436$ & $1.290$ & $0.708$ & $1.479$ \\
\textbf{FW-D} & $0.312$ & $\mathbf{0.785}$ & $9.11\!\times\!10^{-3}$ & $\mathbf{0.8445}$ & $0.698$ & $1.009$ \\
Greedy-E & $0.470$ & $1.683$ & $0.0167$ & $1.748$ & $0.694$ & $1.927$ \\
A-optimal & $0.468$ & $1.429$ & $0.0350$ & $1.504$ & $0.703$ & $1.620$ \\
\textbf{D-optimal} & $0.569$ & $1.213$ & $\mathbf{4.44\!\times\!10^{-3}}$ & $1.340$ & $\mathbf{0.690}$ & $1.327$ \\
Random & $2.153$ & $4.276$ & $0.0296$ & $4.788$ & $0.693$ & $4.419$ \\
Coverage & $1.034$ & $3.578$ & $0.0373$ & $3.725$ & $0.694$ & $3.622$ \\
\addlinespace[1pt]
\multicolumn{7}{l}{\textit{$K{=}15$}} \\[1pt]
FW-E & $0.138$ & $0.768$ & $0.0199$ & $0.7803$ & $0.695$ & $1.003$ \\
FW-A & $0.431$ & $0.470$ & $0.0372$ & $0.6389$ & $0.701$ & $0.743$ \\
FW-D & $0.543$ & $0.990$ & $7.83\!\times\!10^{-3}$ & $1.129$ & $0.689$ & $1.186$ \\
\textbf{Greedy-E} & $\mathbf{0.0767}$ & $0.731$ & $0.0140$ & $0.7347$ & $0.697$ & $1.009$ \\
\textbf{A-optimal} & $0.384$ & $\mathbf{0.226}$ & $0.0497$ & $\mathbf{0.4483}$ & $0.696$ & $\mathbf{0.681}$ \\
\textbf{D-optimal} & $0.153$ & $1.032$ & $\mathbf{7.72\!\times\!10^{-3}}$ & $1.043$ & $\mathbf{0.684}$ & $1.242$ \\
Random & $1.026$ & $2.232$ & $0.0233$ & $2.457$ & $0.699$ & $2.355$ \\
Coverage & $0.827$ & $2.978$ & $0.0239$ & $3.091$ & $0.694$ & $3.090$ \\
\addlinespace[1pt]
\multicolumn{7}{l}{\textit{$K{=}20$}} \\[1pt]
FW-E & $0.203$ & $0.648$ & $0.0117$ & $0.6794$ & $0.695$ & $0.904$ \\
\textbf{FW-A} & $0.360$ & $\mathbf{0.489}$ & $0.0316$ & $0.6083$ & $0.697$ & $\mathbf{0.819}$ \\
\textbf{FW-D} & $0.392$ & $1.239$ & $0.0133$ & $1.299$ & $\mathbf{0.693}$ & $1.365$ \\
\textbf{Greedy-E} & $\mathbf{0.150}$ & $0.530$ & $\mathbf{3.43\!\times\!10^{-3}}$ & $\mathbf{0.5505}$ & $0.696$ & $0.830$ \\
A-optimal & $0.328$ & $0.623$ & $0.0313$ & $0.7044$ & $0.696$ & $0.893$ \\
D-optimal & $0.202$ & $1.076$ & $4.42\!\times\!10^{-3}$ & $1.095$ & $0.693$ & $1.248$ \\
Random & $0.753$ & $1.706$ & $0.0248$ & $1.865$ & $0.696$ & $1.866$ \\
Coverage & $0.278$ & $1.586$ & $0.0299$ & $1.611$ & $0.702$ & $1.733$ \\
\addlinespace[1pt]
\multicolumn{7}{l}{\textit{$K{=}30$}} \\[1pt]
FW-E & $0.242$ & $0.528$ & $0.0127$ & $0.5813$ & $0.698$ & $0.836$ \\
\textbf{FW-A} & $0.581$ & $0.963$ & $0.0312$ & $1.126$ & $\mathbf{0.693}$ & $1.212$ \\
\textbf{FW-D} & $0.254$ & $\mathbf{0.236}$ & $0.0172$ & $0.3472$ & $0.694$ & $\mathbf{0.653}$ \\
\textbf{Greedy-E} & $\mathbf{0.0964}$ & $0.570$ & $\mathbf{4.22\!\times\!10^{-3}}$ & $0.5782$ & $0.698$ & $0.866$ \\
\textbf{A-optimal} & $0.581$ & $0.963$ & $0.0312$ & $1.126$ & $\mathbf{0.693}$ & $1.212$ \\
D-optimal & $0.135$ & $0.247$ & $0.0134$ & $\mathbf{0.2816}$ & $0.695$ & $0.658$ \\
Random & $0.693$ & $1.556$ & $0.0154$ & $1.703$ & $0.695$ & $1.686$ \\
Coverage & $0.490$ & $1.770$ & $0.0330$ & $1.837$ & $0.702$ & $1.850$ \\
\addlinespace[1pt]
\multicolumn{7}{l}{\textit{$K{=}40$}} \\[1pt]
FW-E & $0.525$ & $0.556$ & $0.0210$ & $0.7653$ & $0.693$ & $0.855$ \\
FW-A & $0.369$ & $0.991$ & $0.0313$ & $1.058$ & $0.693$ & $1.212$ \\
\textbf{FW-D} & $\mathbf{0.0990}$ & $0.250$ & $0.0214$ & $0.2693$ & $0.691$ & $0.675$ \\
\textbf{Greedy-E} & $0.312$ & $0.693$ & $\mathbf{5.98\!\times\!10^{-3}}$ & $0.7602$ & $0.693$ & $0.954$ \\
A-optimal & $0.421$ & $1.177$ & $0.0302$ & $1.251$ & $0.694$ & $1.376$ \\
\textbf{D-optimal} & $0.114$ & $\mathbf{0.212}$ & $0.0250$ & $\mathbf{0.2419}$ & $\mathbf{0.689}$ & $\mathbf{0.662}$ \\
Random & $0.538$ & $1.112$ & $0.0174$ & $1.235$ & $0.698$ & $1.324$ \\
Coverage & $0.661$ & $1.539$ & $0.0334$ & $1.675$ & $0.698$ & $1.620$ \\
\addlinespace[1pt]
\multicolumn{7}{l}{\textit{$K{=}50$}} \\[1pt]
FW-E & $0.456$ & $0.835$ & $4.38\!\times\!10^{-3}$ & $0.9516$ & $0.692$ & $1.049$ \\
FW-A & $0.234$ & $0.920$ & $0.0217$ & $0.9496$ & $0.694$ & $1.151$ \\
FW-D & $0.201$ & $0.448$ & $0.0206$ & $0.4918$ & $0.690$ & $0.764$ \\
\textbf{Greedy-E} & $0.423$ & $0.794$ & $\mathbf{3.33\!\times\!10^{-3}}$ & $0.8996$ & $0.693$ & $1.028$ \\
A-optimal & $0.376$ & $1.068$ & $0.0320$ & $1.133$ & $0.693$ & $1.285$ \\
\textbf{D-optimal} & $\mathbf{0.177}$ & $\mathbf{0.352}$ & $0.0225$ & $\mathbf{0.3946}$ & $\mathbf{0.689}$ & $\mathbf{0.718}$ \\
Random & $0.447$ & $1.034$ & $0.0119$ & $1.126$ & $0.696$ & $1.239$ \\
Coverage & $0.828$ & $1.779$ & $0.0357$ & $1.962$ & $0.696$ & $1.835$ \\
\bottomrule
\end{tabular}
\end{table*}

\begin{figure*}[t]
\centering
\includegraphics[width=\textwidth]{tab23_ksweep_aprilgrid_sigma05}
\caption{Method comparison: K-sweep: AprilGrid, $\sigma{=}0.5$\,px. Bars show each metric (winner outlined in black). $\lambda_{\min}$/$\log\det$: higher is better; all others: lower is better.}
\label{fig:tab23_ksweep_aprilgrid_sigma05}
\end{figure*}


\begin{table*}[ht]
\centering
\setlength{\tabcolsep}{4pt}
\caption{ChArUco, $\sigma{=}0.5$\,px. $N{=}720$, $K{=}20$. Focal/PP errors in pixels, distortion error unitless. Held-out: mean reprojection RMS on all $N{-}K$ unselected frames per method. Random: mean over 25 trials. Bold: best per column per $K$.}
\label{tab:ksweep_charuco_sigma05}
\scriptsize
\renewcommand{\arraystretch}{0.80}
\begin{tabular}{lrrrrrr}
\toprule
Method & $\epsilon_f$\,(px) & $\epsilon_{pp}$\,(px) & $\epsilon_d$ & $\|\Delta\theta\|$ & Train-RMS & Held-out \\
\midrule
\multicolumn{7}{l}{\textit{$K{=}5$}} \\[1pt]
FW-E & $39.414$ & $39.777$ & $0.0349$ & $56.00$ & $0.683$ & $38.214$ \\
FW-A & $18.881$ & $37.640$ & $0.0125$ & $42.11$ & $0.686$ & $38.302$ \\
FW-D & $1.520$ & $4.296$ & $0.0192$ & $4.557$ & $0.714$ & $4.382$ \\
\textbf{Greedy-E} & $\mathbf{0.656}$ & $\mathbf{1.262}$ & $0.0847$ & $\mathbf{1.425}$ & $0.709$ & $\mathbf{1.473}$ \\
\textbf{A-optimal} & $2.299$ & $4.583$ & $\mathbf{1.58\!\times\!10^{-3}}$ & $5.128$ & $0.691$ & $4.570$ \\
D-optimal & $1.676$ & $5.025$ & $5.48\!\times\!10^{-3}$ & $5.297$ & $0.723$ & $5.148$ \\
Random & $3.765$ & $5.510$ & $0.0367$ & $6.674$ & $0.691$ & $5.559$ \\
\textbf{Coverage} & $14.644$ & $41.072$ & $0.0170$ & $43.61$ & $\mathbf{0.670}$ & $40.625$ \\
\addlinespace[1pt]
\multicolumn{7}{l}{\textit{$K{=}10$}} \\[1pt]
\textbf{FW-E} & $0.452$ & $1.440$ & $\mathbf{9.23\!\times\!10^{-3}}$ & $1.510$ & $\mathbf{0.683}$ & $1.639$ \\
FW-A & $1.618$ & $3.452$ & $0.0313$ & $3.812$ & $0.692$ & $3.591$ \\
FW-D & $0.849$ & $2.631$ & $0.0122$ & $2.765$ & $0.693$ & $2.739$ \\
Greedy-E & $0.636$ & $2.100$ & $0.0214$ & $2.195$ & $0.703$ & $2.367$ \\
\textbf{A-optimal} & $\mathbf{0.0337}$ & $\mathbf{1.059}$ & $0.0138$ & $\mathbf{1.060}$ & $0.694$ & $\mathbf{1.281}$ \\
D-optimal & $0.718$ & $1.971$ & $0.0209$ & $2.098$ & $0.695$ & $2.132$ \\
Random & $2.626$ & $4.584$ & $0.0292$ & $5.282$ & $0.693$ & $4.697$ \\
Coverage & $0.226$ & $1.401$ & $0.0476$ & $1.420$ & $0.695$ & $1.501$ \\
\addlinespace[1pt]
\multicolumn{7}{l}{\textit{$K{=}15$}} \\[1pt]
FW-E & $0.485$ & $1.672$ & $7.34\!\times\!10^{-3}$ & $1.741$ & $0.699$ & $1.798$ \\
FW-A & $0.485$ & $1.672$ & $7.34\!\times\!10^{-3}$ & $1.741$ & $0.699$ & $1.798$ \\
\textbf{FW-D} & $0.738$ & $\mathbf{1.135}$ & $\mathbf{2.75\!\times\!10^{-3}}$ & $1.353$ & $0.702$ & $\mathbf{1.332}$ \\
Greedy-E & $0.546$ & $1.452$ & $6.05\!\times\!10^{-3}$ & $1.551$ & $0.698$ & $1.668$ \\
\textbf{A-optimal} & $\mathbf{0.209}$ & $1.180$ & $0.0277$ & $\mathbf{1.198}$ & $0.696$ & $1.390$ \\
\textbf{D-optimal} & $0.738$ & $\mathbf{1.135}$ & $\mathbf{2.75\!\times\!10^{-3}}$ & $1.353$ & $0.702$ & $\mathbf{1.332}$ \\
Random & $1.210$ & $2.254$ & $0.0290$ & $2.559$ & $0.693$ & $2.369$ \\
\textbf{Coverage} & $1.539$ & $3.186$ & $0.0422$ & $3.538$ & $\mathbf{0.690}$ & $3.260$ \\
\addlinespace[1pt]
\multicolumn{7}{l}{\textit{$K{=}20$}} \\[1pt]
\textbf{FW-E} & $0.293$ & $\mathbf{1.089}$ & $8.81\!\times\!10^{-3}$ & $\mathbf{1.128}$ & $0.695$ & $\mathbf{1.271}$ \\
\textbf{FW-A} & $0.293$ & $\mathbf{1.089}$ & $8.81\!\times\!10^{-3}$ & $\mathbf{1.128}$ & $0.695$ & $\mathbf{1.271}$ \\
\textbf{FW-D} & $0.859$ & $1.141$ & $\mathbf{3.39\!\times\!10^{-3}}$ & $1.428$ & $0.699$ & $1.342$ \\
Greedy-E & $0.424$ & $1.316$ & $5.01\!\times\!10^{-3}$ & $1.383$ & $0.700$ & $1.529$ \\
\textbf{A-optimal} & $\mathbf{0.191}$ & $1.182$ & $0.0128$ & $1.198$ & $\mathbf{0.690}$ & $1.395$ \\
D-optimal & $0.864$ & $1.150$ & $9.36\!\times\!10^{-3}$ & $1.438$ & $0.701$ & $1.348$ \\
Random & $0.707$ & $1.302$ & $0.0171$ & $1.481$ & $0.696$ & $1.479$ \\
Coverage & $1.077$ & $2.458$ & $0.0358$ & $2.684$ & $0.692$ & $2.627$ \\
\addlinespace[1pt]
\multicolumn{7}{l}{\textit{$K{=}30$}} \\[1pt]
FW-E & $0.411$ & $1.315$ & $2.69\!\times\!10^{-3}$ & $1.378$ & $0.699$ & $1.467$ \\
\textbf{FW-A} & $\mathbf{0.146}$ & $\mathbf{0.702}$ & $\mathbf{1.50\!\times\!10^{-3}}$ & $\mathbf{0.7174}$ & $0.697$ & $\mathbf{0.953}$ \\
FW-D & $0.660$ & $0.865$ & $4.51\!\times\!10^{-3}$ & $1.087$ & $0.705$ & $1.100$ \\
Greedy-E & $0.464$ & $1.049$ & $3.71\!\times\!10^{-3}$ & $1.147$ & $0.697$ & $1.215$ \\
\textbf{A-optimal} & $0.426$ & $1.072$ & $3.77\!\times\!10^{-3}$ & $1.154$ & $\mathbf{0.696}$ & $1.263$ \\
D-optimal & $0.700$ & $0.840$ & $8.74\!\times\!10^{-3}$ & $1.094$ & $0.704$ & $1.059$ \\
Random & $0.610$ & $1.201$ & $0.0163$ & $1.347$ & $0.697$ & $1.388$ \\
Coverage & $0.454$ & $1.475$ & $0.0448$ & $1.544$ & $0.698$ & $1.712$ \\
\addlinespace[1pt]
\multicolumn{7}{l}{\textit{$K{=}40$}} \\[1pt]
FW-E & $0.429$ & $1.360$ & $6.21\!\times\!10^{-3}$ & $1.426$ & $0.701$ & $1.516$ \\
FW-A & $0.409$ & $0.984$ & $0.0162$ & $1.065$ & $0.700$ & $1.145$ \\
FW-D & $0.553$ & $0.941$ & $9.30\!\times\!10^{-3}$ & $1.091$ & $0.700$ & $1.164$ \\
\textbf{Greedy-E} & $\mathbf{0.336}$ & $1.218$ & $\mathbf{1.25\!\times\!10^{-3}}$ & $1.264$ & $0.701$ & $1.401$ \\
\textbf{A-optimal} & $0.390$ & $\mathbf{0.779}$ & $0.0142$ & $\mathbf{0.8710}$ & $0.698$ & $\mathbf{0.980}$ \\
D-optimal & $0.584$ & $0.975$ & $6.38\!\times\!10^{-3}$ & $1.137$ & $0.700$ & $1.184$ \\
Random & $0.612$ & $1.177$ & $0.0134$ & $1.327$ & $0.695$ & $1.360$ \\
\textbf{Coverage} & $0.676$ & $1.408$ & $0.0565$ & $1.563$ & $\mathbf{0.695}$ & $1.592$ \\
\addlinespace[1pt]
\multicolumn{7}{l}{\textit{$K{=}50$}} \\[1pt]
FW-E & $0.353$ & $1.022$ & $4.43\!\times\!10^{-3}$ & $1.081$ & $0.698$ & $1.207$ \\
FW-A & $0.442$ & $1.121$ & $8.86\!\times\!10^{-3}$ & $1.205$ & $0.701$ & $1.250$ \\
FW-D & $0.665$ & $0.892$ & $3.58\!\times\!10^{-3}$ & $1.113$ & $0.697$ & $1.110$ \\
\textbf{Greedy-E} & $\mathbf{0.269}$ & $1.055$ & $4.45\!\times\!10^{-3}$ & $1.089$ & $0.701$ & $1.254$ \\
\textbf{A-optimal} & $0.272$ & $\mathbf{0.776}$ & $0.0132$ & $\mathbf{0.8227}$ & $0.697$ & $\mathbf{0.974}$ \\
\textbf{D-optimal} & $0.617$ & $0.856$ & $\mathbf{2.59\!\times\!10^{-3}}$ & $1.056$ & $0.698$ & $1.080$ \\
\textbf{Random} & $0.483$ & $0.990$ & $0.0135$ & $1.102$ & $\mathbf{0.694}$ & $1.199$ \\
Coverage & $0.555$ & $1.096$ & $0.0349$ & $1.229$ & $0.697$ & $1.262$ \\
\bottomrule
\end{tabular}
\end{table*}

\begin{figure*}[t]
\centering
\includegraphics[width=\textwidth]{tab24_ksweep_charuco_sigma05}
\caption{Method comparison: K-sweep: ChArUco, $\sigma{=}0.5$\,px. Bars show each metric (winner outlined in black). $\lambda_{\min}$/$\log\det$: higher is better; all others: lower is better.}
\label{fig:tab24_ksweep_charuco_sigma05}
\end{figure*}


\begin{table*}[ht]
\centering
\setlength{\tabcolsep}{4pt}
\caption{AprilGrid, $\sigma{=}1.0$\,px. $N{=}720$, $K{=}20$. Focal/PP errors in pixels, distortion error unitless. Held-out: mean reprojection RMS on all $N{-}K$ unselected frames per method. Random: mean over 25 trials. Bold: best per column per $K$.}
\label{tab:ksweep_aprilgrid_sigma10}
\scriptsize
\renewcommand{\arraystretch}{0.80}
\begin{tabular}{lrrrrrr}
\toprule
Method & $\epsilon_f$\,(px) & $\epsilon_{pp}$\,(px) & $\epsilon_d$ & $\|\Delta\theta\|$ & Train-RMS & Held-out \\
\midrule
\multicolumn{7}{l}{\textit{$K{=}5$}} \\[1pt]
\textbf{FW-E} & $68.909$ & $61.976$ & $0.0321$ & $92.68$ & $\mathbf{1.387}$ & $59.088$ \\
FW-A & $1.275$ & $3.355$ & $0.0404$ & $3.590$ & $1.414$ & $3.517$ \\
\textbf{FW-D} & $0.708$ & $\mathbf{1.239}$ & $\mathbf{8.07\!\times\!10^{-3}}$ & $\mathbf{1.427}$ & $1.398$ & $\mathbf{1.873}$ \\
\textbf{Greedy-E} & $\mathbf{0.439}$ & $3.874$ & $0.0508$ & $3.900$ & $1.410$ & $4.357$ \\
A-optimal & $2.888$ & $4.647$ & $0.0615$ & $5.471$ & $1.424$ & $4.710$ \\
\textbf{D-optimal} & $0.708$ & $\mathbf{1.239}$ & $8.07\!\times\!10^{-3}$ & $\mathbf{1.427}$ & $1.398$ & $\mathbf{1.873}$ \\
Random & $7.490$ & $13.412$ & $0.0686$ & $15.36$ & $1.391$ & $13.558$ \\
Coverage & $19.460$ & $49.894$ & $0.128$ & $53.55$ & $1.411$ & $49.145$ \\
\addlinespace[1pt]
\multicolumn{7}{l}{\textit{$K{=}10$}} \\[1pt]
FW-E & $1.820$ & $1.556$ & $0.0462$ & $2.396$ & $1.398$ & $1.934$ \\
\textbf{FW-A} & $\mathbf{0.404}$ & $2.612$ & $0.0895$ & $2.645$ & $1.417$ & $3.016$ \\
FW-D & $0.619$ & $1.564$ & $0.0179$ & $1.682$ & $1.395$ & $2.014$ \\
\textbf{Greedy-E} & $1.895$ & $\mathbf{1.046}$ & $9.38\!\times\!10^{-3}$ & $2.165$ & $1.404$ & $\mathbf{1.588}$ \\
A-optimal & $0.696$ & $1.241$ & $0.102$ & $\mathbf{1.426}$ & $1.411$ & $1.733$ \\
\textbf{D-optimal} & $1.138$ & $2.433$ & $\mathbf{8.57\!\times\!10^{-3}}$ & $2.686$ & $\mathbf{1.381}$ & $2.661$ \\
Random & $3.982$ & $8.046$ & $0.0587$ & $8.978$ & $1.386$ & $8.320$ \\
Coverage & $2.024$ & $7.029$ & $0.0737$ & $7.315$ & $1.389$ & $7.120$ \\
\addlinespace[1pt]
\multicolumn{7}{l}{\textit{$K{=}15$}} \\[1pt]
FW-E & $0.460$ & $2.142$ & $0.0311$ & $2.191$ & $1.404$ & $2.503$ \\
FW-A & $0.877$ & $0.996$ & $0.0749$ & $1.329$ & $1.403$ & $1.514$ \\
FW-D & $1.089$ & $1.994$ & $0.0155$ & $2.272$ & $1.378$ & $2.383$ \\
Greedy-E & $0.629$ & $0.536$ & $0.0438$ & $\mathbf{0.8274}$ & $1.421$ & $1.385$ \\
\textbf{A-optimal} & $0.773$ & $\mathbf{0.379}$ & $0.0999$ & $0.8666$ & $1.392$ & $\mathbf{1.338}$ \\
\textbf{D-optimal} & $\mathbf{0.153}$ & $2.536$ & $\mathbf{4.04\!\times\!10^{-3}}$ & $2.541$ & $\mathbf{1.372}$ & $2.938$ \\
Random & $2.041$ & $4.437$ & $0.0462$ & $4.884$ & $1.397$ & $4.683$ \\
Coverage & $1.659$ & $5.983$ & $0.0507$ & $6.209$ & $1.388$ & $6.210$ \\
\addlinespace[1pt]
\multicolumn{7}{l}{\textit{$K{=}20$}} \\[1pt]
FW-E & $0.559$ & $1.994$ & $0.0364$ & $2.071$ & $1.398$ & $2.404$ \\
FW-A & $0.698$ & $1.447$ & $0.0735$ & $1.608$ & $1.395$ & $1.955$ \\
\textbf{FW-D} & $0.783$ & $2.480$ & $0.0267$ & $2.601$ & $\mathbf{1.386}$ & $2.734$ \\
\textbf{Greedy-E} & $\mathbf{0.260}$ & $1.476$ & $0.0526$ & $1.499$ & $1.409$ & $2.002$ \\
\textbf{A-optimal} & $0.656$ & $\mathbf{1.242}$ & $0.0631$ & $\mathbf{1.406}$ & $1.392$ & $\mathbf{1.782}$ \\
\textbf{D-optimal} & $0.405$ & $2.157$ & $\mathbf{8.82\!\times\!10^{-3}}$ & $2.194$ & $1.386$ & $2.499$ \\
Random & $1.511$ & $3.417$ & $0.0498$ & $3.737$ & $1.392$ & $3.736$ \\
Coverage & $0.551$ & $3.232$ & $0.0621$ & $3.279$ & $1.405$ & $3.528$ \\
\addlinespace[1pt]
\multicolumn{7}{l}{\textit{$K{=}30$}} \\[1pt]
\textbf{FW-E} & $0.483$ & $1.035$ & $\mathbf{0.0256}$ & $1.142$ & $1.397$ & $1.656$ \\
\textbf{FW-A} & $1.160$ & $1.906$ & $0.0629$ & $2.232$ & $\mathbf{1.387}$ & $2.403$ \\
FW-D & $0.506$ & $0.464$ & $0.0345$ & $0.6873$ & $1.388$ & $1.305$ \\
\textbf{Greedy-E} & $\mathbf{0.445}$ & $0.853$ & $0.0400$ & $0.9631$ & $1.401$ & $1.530$ \\
\textbf{A-optimal} & $1.160$ & $1.906$ & $0.0629$ & $2.232$ & $\mathbf{1.387}$ & $2.403$ \\
\textbf{D-optimal} & $0.498$ & $\mathbf{0.159}$ & $0.0345$ & $\mathbf{0.5237}$ & $1.388$ & $\mathbf{1.260}$ \\
Random & $1.390$ & $3.117$ & $0.0308$ & $3.413$ & $1.389$ & $3.377$ \\
Coverage & $0.988$ & $3.569$ & $0.0668$ & $3.704$ & $1.403$ & $3.727$ \\
\addlinespace[1pt]
\multicolumn{7}{l}{\textit{$K{=}40$}} \\[1pt]
\textbf{FW-E} & $1.021$ & $1.028$ & $\mathbf{0.0263}$ & $1.449$ & $1.387$ & $1.657$ \\
FW-A & $0.738$ & $1.985$ & $0.0630$ & $2.119$ & $1.387$ & $2.426$ \\
\textbf{FW-D} & $\mathbf{0.196}$ & $0.494$ & $0.0429$ & $0.5334$ & $1.381$ & $1.348$ \\
Greedy-E & $0.963$ & $1.227$ & $0.0333$ & $1.560$ & $1.389$ & $1.764$ \\
A-optimal & $0.840$ & $2.355$ & $0.0608$ & $2.502$ & $1.387$ & $2.751$ \\
\textbf{D-optimal} & $0.231$ & $\mathbf{0.431}$ & $0.0502$ & $\mathbf{0.4915}$ & $\mathbf{1.379}$ & $\mathbf{1.326}$ \\
Random & $1.080$ & $2.234$ & $0.0349$ & $2.482$ & $1.395$ & $2.656$ \\
Coverage & $1.322$ & $3.080$ & $0.0671$ & $3.352$ & $1.396$ & $3.242$ \\
\addlinespace[1pt]
\multicolumn{7}{l}{\textit{$K{=}50$}} \\[1pt]
\textbf{FW-E} & $0.876$ & $1.566$ & $\mathbf{9.12\!\times\!10^{-3}}$ & $1.795$ & $1.385$ & $1.998$ \\
FW-A & $0.467$ & $1.835$ & $0.0434$ & $1.894$ & $1.387$ & $2.295$ \\
\textbf{FW-D} & $0.354$ & $0.705$ & $0.0450$ & $0.7905$ & $\mathbf{1.379}$ & $\mathbf{1.437}$ \\
Greedy-E & $0.902$ & $1.452$ & $0.0217$ & $1.709$ & $1.386$ & $1.940$ \\
A-optimal & $0.751$ & $2.136$ & $0.0641$ & $2.265$ & $1.385$ & $2.568$ \\
\textbf{D-optimal} & $\mathbf{0.354}$ & $\mathbf{0.705}$ & $0.0450$ & $\mathbf{0.7905}$ & $\mathbf{1.379}$ & $\mathbf{1.437}$ \\
Random & $0.894$ & $2.066$ & $0.0238$ & $2.251$ & $1.393$ & $2.476$ \\
Coverage & $1.653$ & $3.557$ & $0.0715$ & $3.922$ & $1.393$ & $3.669$ \\
\bottomrule
\end{tabular}
\end{table*}

\begin{figure*}[t]
\centering
\includegraphics[width=\textwidth]{tab25_ksweep_aprilgrid_sigma10}
\caption{Method comparison: K-sweep: AprilGrid, $\sigma{=}1.0$\,px. Bars show each metric (winner outlined in black). $\lambda_{\min}$/$\log\det$: higher is better; all others: lower is better.}
\label{fig:tab25_ksweep_aprilgrid_sigma10}
\end{figure*}


\begin{table*}[ht]
\centering
\setlength{\tabcolsep}{4pt}
\caption{ChArUco, $\sigma{=}1.0$\,px. $N{=}720$, $K{=}20$. Focal/PP errors in pixels, distortion error unitless. Held-out: mean reprojection RMS on all $N{-}K$ unselected frames per method. Random: mean over 25 trials. Bold: best per column per $K$.}
\label{tab:ksweep_charuco_sigma10}
\scriptsize
\renewcommand{\arraystretch}{0.80}
\begin{tabular}{lrrrrrr}
\toprule
Method & $\epsilon_f$\,(px) & $\epsilon_{pp}$\,(px) & $\epsilon_d$ & $\|\Delta\theta\|$ & Train-RMS & Held-out \\
\midrule
\multicolumn{7}{l}{\textit{$K{=}5$}} \\[1pt]
FW-E & $33.582$ & $38.011$ & $0.0853$ & $50.72$ & $1.366$ & $36.858$ \\
FW-A & $9.730$ & $16.826$ & $0.0246$ & $19.44$ & $1.371$ & $17.418$ \\
FW-D & $3.038$ & $8.527$ & $0.0342$ & $9.052$ & $1.428$ & $8.697$ \\
\textbf{Greedy-E} & $\mathbf{1.022}$ & $5.605$ & $0.100$ & $5.699$ & $1.381$ & $6.150$ \\
\textbf{A-optimal} & $1.709$ & $\mathbf{4.847}$ & $0.0415$ & $\mathbf{5.139}$ & $1.380$ & $\mathbf{5.091}$ \\
\textbf{D-optimal} & $3.357$ & $10.045$ & $\mathbf{0.0104}$ & $10.59$ & $1.446$ & $10.293$ \\
Random & $7.050$ & $10.277$ & $0.0745$ & $12.46$ & $1.382$ & $10.392$ \\
\textbf{Coverage} & $3.267$ & $8.521$ & $0.0256$ & $9.125$ & $\mathbf{1.341}$ & $8.698$ \\
\addlinespace[1pt]
\multicolumn{7}{l}{\textit{$K{=}10$}} \\[1pt]
\textbf{FW-E} & $0.900$ & $2.876$ & $0.0185$ & $3.014$ & $\mathbf{1.367}$ & $3.272$ \\
FW-A & $3.233$ & $6.902$ & $0.0637$ & $7.622$ & $1.384$ & $7.181$ \\
FW-D & $1.701$ & $5.253$ & $0.0239$ & $5.521$ & $1.385$ & $5.467$ \\
\textbf{Greedy-E} & $1.459$ & $4.147$ & $\mathbf{3.47\!\times\!10^{-3}}$ & $4.396$ & $1.371$ & $4.542$ \\
\textbf{A-optimal} & $0.649$ & $\mathbf{2.406}$ & $0.0169$ & $\mathbf{2.493}$ & $1.377$ & $\mathbf{2.828}$ \\
D-optimal & $1.442$ & $3.928$ & $0.0414$ & $4.184$ & $1.390$ & $4.250$ \\
Random & $5.237$ & $9.124$ & $0.0587$ & $10.52$ & $1.387$ & $9.352$ \\
\textbf{Coverage} & $\mathbf{0.451}$ & $2.871$ & $0.0957$ & $2.908$ & $1.390$ & $3.066$ \\
\addlinespace[1pt]
\multicolumn{7}{l}{\textit{$K{=}15$}} \\[1pt]
FW-E & $0.975$ & $3.366$ & $0.0152$ & $3.504$ & $1.397$ & $3.618$ \\
FW-A & $0.975$ & $3.366$ & $0.0152$ & $3.504$ & $1.397$ & $3.618$ \\
\textbf{FW-D} & $1.479$ & $2.269$ & $\mathbf{5.39\!\times\!10^{-3}}$ & $2.708$ & $1.403$ & $2.663$ \\
\textbf{Greedy-E} & $0.842$ & $3.789$ & $8.60\!\times\!10^{-3}$ & $3.881$ & $\mathbf{1.377}$ & $4.191$ \\
\textbf{A-optimal} & $\mathbf{0.632}$ & $2.650$ & $0.0445$ & $2.725$ & $1.387$ & $3.030$ \\
\textbf{D-optimal} & $1.479$ & $\mathbf{2.269}$ & $5.39\!\times\!10^{-3}$ & $\mathbf{2.708}$ & $1.403$ & $\mathbf{2.663}$ \\
Random & $2.423$ & $4.526$ & $0.0582$ & $5.134$ & $1.385$ & $4.756$ \\
Coverage & $3.061$ & $6.329$ & $0.0823$ & $7.031$ & $1.379$ & $6.479$ \\
\addlinespace[1pt]
\multicolumn{7}{l}{\textit{$K{=}20$}} \\[1pt]
\textbf{FW-E} & $\mathbf{0.590}$ & $2.206$ & $0.0186$ & $2.284$ & $1.389$ & $2.569$ \\
\textbf{FW-A} & $\mathbf{0.590}$ & $2.206$ & $0.0186$ & $2.284$ & $1.389$ & $2.569$ \\
\textbf{FW-D} & $1.719$ & $2.295$ & $\mathbf{6.47\!\times\!10^{-3}}$ & $2.868$ & $1.398$ & $2.699$ \\
Greedy-E & $0.944$ & $3.243$ & $0.0327$ & $3.377$ & $1.382$ & $3.668$ \\
\textbf{A-optimal} & $0.725$ & $\mathbf{2.088}$ & $0.0101$ & $\mathbf{2.210}$ & $\mathbf{1.381}$ & $\mathbf{2.483}$ \\
D-optimal & $1.731$ & $2.309$ & $0.0187$ & $2.886$ & $1.402$ & $2.705$ \\
Random & $1.415$ & $2.607$ & $0.0343$ & $2.967$ & $1.391$ & $2.962$ \\
Coverage & $2.141$ & $4.913$ & $0.0716$ & $5.360$ & $1.385$ & $5.256$ \\
\addlinespace[1pt]
\multicolumn{7}{l}{\textit{$K{=}30$}} \\[1pt]
FW-E & $0.826$ & $2.654$ & $6.35\!\times\!10^{-3}$ & $2.780$ & $1.398$ & $2.958$ \\
\textbf{FW-A} & $\mathbf{0.292}$ & $\mathbf{1.411}$ & $\mathbf{3.25\!\times\!10^{-3}}$ & $\mathbf{1.441}$ & $1.394$ & $\mathbf{1.912}$ \\
FW-D & $1.321$ & $1.745$ & $9.54\!\times\!10^{-3}$ & $2.189$ & $1.411$ & $2.215$ \\
Greedy-E & $0.771$ & $2.793$ & $6.27\!\times\!10^{-3}$ & $2.898$ & $1.401$ & $3.112$ \\
\textbf{A-optimal} & $0.702$ & $1.613$ & $7.17\!\times\!10^{-3}$ & $1.759$ & $\mathbf{1.387}$ & $2.023$ \\
D-optimal & $1.403$ & $1.695$ & $0.0177$ & $2.201$ & $1.407$ & $2.131$ \\
Random & $1.222$ & $2.404$ & $0.0327$ & $2.697$ & $1.394$ & $2.779$ \\
Coverage & $0.920$ & $2.943$ & $0.0898$ & $3.085$ & $1.397$ & $3.416$ \\
\addlinespace[1pt]
\multicolumn{7}{l}{\textit{$K{=}40$}} \\[1pt]
FW-E & $0.860$ & $2.727$ & $0.0135$ & $2.860$ & $1.401$ & $3.037$ \\
FW-A & $0.817$ & $1.973$ & $0.0329$ & $2.136$ & $1.400$ & $2.295$ \\
FW-D & $1.107$ & $1.894$ & $0.0188$ & $2.194$ & $1.400$ & $2.340$ \\
\textbf{Greedy-E} & $\mathbf{0.608}$ & $2.459$ & $0.0135$ & $2.533$ & $1.396$ & $2.801$ \\
\textbf{A-optimal} & $0.837$ & $\mathbf{1.719}$ & $0.0253$ & $\mathbf{1.912}$ & $1.394$ & $\mathbf{2.080}$ \\
\textbf{D-optimal} & $1.170$ & $1.960$ & $\mathbf{0.0129}$ & $2.283$ & $1.401$ & $2.378$ \\
Random & $1.225$ & $2.359$ & $0.0270$ & $2.658$ & $1.391$ & $2.724$ \\
\textbf{Coverage} & $1.367$ & $2.851$ & $0.114$ & $3.164$ & $\mathbf{1.390}$ & $3.218$ \\
\addlinespace[1pt]
\multicolumn{7}{l}{\textit{$K{=}50$}} \\[1pt]
FW-E & $0.706$ & $2.040$ & $9.41\!\times\!10^{-3}$ & $2.159$ & $1.396$ & $2.411$ \\
FW-A & $0.884$ & $2.243$ & $0.0179$ & $2.411$ & $1.402$ & $2.502$ \\
FW-D & $1.332$ & $1.791$ & $7.09\!\times\!10^{-3}$ & $2.232$ & $1.393$ & $2.227$ \\
Greedy-E & $0.964$ & $2.585$ & $0.0179$ & $2.759$ & $1.397$ & $2.894$ \\
\textbf{A-optimal} & $\mathbf{0.584}$ & $\mathbf{1.685}$ & $0.0243$ & $\mathbf{1.783}$ & $1.392$ & $\mathbf{2.048}$ \\
\textbf{D-optimal} & $1.237$ & $1.720$ & $\mathbf{5.16\!\times\!10^{-3}}$ & $2.119$ & $1.395$ & $2.166$ \\
\textbf{Random} & $0.966$ & $1.982$ & $0.0271$ & $2.205$ & $\mathbf{1.389}$ & $2.399$ \\
Coverage & $1.126$ & $2.224$ & $0.0705$ & $2.494$ & $1.393$ & $2.553$ \\
\bottomrule
\end{tabular}
\end{table*}

\begin{figure*}[t]
\centering
\includegraphics[width=\textwidth]{tab26_ksweep_charuco_sigma10}
\caption{Method comparison: K-sweep: ChArUco, $\sigma{=}1.0$\,px. Bars show each metric (winner outlined in black). $\lambda_{\min}$/$\log\det$: higher is better; all others: lower is better.}
\label{fig:tab26_ksweep_charuco_sigma10}
\end{figure*}


\end{document}
